% આ ભાગનો પૂર્ણ સંપાદનીય પાઠ છે. શૈલી, ફોન્ટ, ચિત્રો અને પ્રકરણસંદર્ભ માટે સંપૂર્ણ સ્રોત ZIP તથા reader/full-reader.tex જુઓ. આ એકલી સ્વતંત્ર કમ્પાઇલ થતી મુખ્ય ફાઇલ નથી.

% BEGIN OLP-0208
\olpart{cmp}{સંગણનીયતા}
\phantomsection\label{guunit:OLP-0208}


\begin{editorial}
  આ ભાગ Jeremy Avigadની સંગણનીયતા સિદ્ધાંત વિષયક નોંધો પર આધારિત છે.
  હજી સુધી માત્ર પુનરાવર્તી વિધેયોના પ્રકરણમાં કસરતો છે. પ્રેરણા,
  ઉદાહરણો, વિગતો અને કસરતો ઉમેરીને આ સમગ્ર ભાગને વિસ્તારી શકાય તેમ છે.
\end{editorial}

\begingroup
\makeatletter\def\ol@sectioncs{section}\makeatother

% BEGIN OLP-0209
\olchapter{cmp}{rec}{પુનરાવર્તી વિધેયો}
\olchapterid{cmp}{rec}
\phantomsection\label{guunit:OLP-0209}


\begin{editorial}
  આ Jeremy Avigadની પુનરાવર્તી વિધેયો વિષયક નોંધો છે, જેને Richard Zachએ
  સુધારી અને વિસ્તારી છે. આ પ્રકરણમાં કેટલીક કસરતો છે અને વાક્યરચનાના
  અંકગણિતીકરણની ચર્ચા માટેનો આધાર આપવા તેને સ્વતંત્ર રીતે પણ સમાવી શકાય છે.
\end{editorial}

\begingroup
\makeatletter\def\ol@sectioncs{section}\makeatother

% BEGIN OLP-0210
\olfileid{cmp}{rec}{int}
\olsection{પરિચય}
\phantomsection\label{guunit:OLP-0210}


સંગણનીયતાનો ગણિતીય સિદ્ધાંત વિકસાવવા માટે સૌપ્રથમ સંગણનીયતાનું એક
\emph{નિદર્શ} ઘડવું પડે. આજે આપણે સંગણનીયતાને કમ્પ્યુટર જે કામ કરે છે
તેના પ્રકાર તરીકે સમજીએ છીએ, અને કમ્પ્યુટર પ્રતીકો સાથે કામ કરે છે. પરંતુ
સંગણનીયતાના સિદ્ધાંતોના વિકાસની શરૂઆતમાં સંગણનાનું આદર્શ ઉદાહરણ
\emph{સંખ્યાત્મક} ગણતરી હતું. ગણિતશાસ્ત્રીઓને હંમેશાં સંખ્યા-સિદ્ધાંતનાં
સંગણન કરી શકાય એવાં વિધેયોમાં, એટલે કે $f\colon \Nat^n \to \Nat$
પ્રકારનાં વિધેયોમાં, રસ રહ્યો છે. તેથી સંગણનીયતાના સિદ્ધાંતની શરૂઆતમાં
આવાં વિધેયોનો અભ્યાસ થયો તે આશ્ચર્યજનક નથી. પ્રાકૃતિક સંખ્યાઓના સરવાળા,
ગુણાકાર અને ઘાતાંક જેવા સૌથી પરિચિત સંગણનીય સંખ્યાત્મક વિધેયોમાં એક
રસપ્રદ સામ્ય છે: તેમને \emph{પુનરાવર્તી રીતે} વ્યાખ્યાયિત કરી શકાય છે.
આથી પુનરાવર્તી વ્યાખ્યાઓને આધારે \emph{સંગણનીય વિધેય}ની સામાન્ય વ્યાખ્યા
આપવાનો પ્રયાસ સ્વાભાવિક છે. સંખ્યા-સિદ્ધાંતનાં વિધેયોને પુનરાવર્તી રીતે
વ્યાખ્યાયિત કરવાની ઘણી રીતોમાંથી અહીં એક ખાસ સરળ પદ્ધતિ કેન્દ્રસ્થાને
છે: કહેવાતું \emph{આદિમ પુનરાવર્તન}.

સંગણનીય વિધેયો ઉપરાંત સંગણનીય ગણો અને સંબંધોમાં પણ આપણને રસ હોઈ શકે.
આપેલી સંખ્યા ગણનો ઘટક છે કે નહીં, તેનો જવાબ સંગણનથી મેળવી
શકાય તો તે ગણ સંગણનીય છે. તેવી જ રીતે, $\tuple{n_1, \dots, n_k}$
બહુજોડ સંબંધનો ઘટક છે કે નહીં, તે સંગણનથી નક્કી કરી શકાય તો
તે સંબંધ સંગણનીય છે. ગણ અથવા સંબંધના \emph{લક્ષણ વિધેય}ને ધ્યાનમાં
લેવાથી, સંગણનીય ગણો અને સંબંધોની ચર્ચા સંગણનીય વિધેયોની ચર્ચામાં
સમાવી શકાય છે. તેથી આદિમ પુનરાવર્તી સંબંધો પણ વ્યાખ્યાયિત કરી શકાય;
દા.ત., ``$n$ વડે $m$ શેષ વિના ભાગાય છે'' એવો સંબંધ આદિમ પુનરાવર્તી છે.

તેમ છતાં, આદિમ પુનરાવર્તી વિધેયો---માત્ર આદિમ પુનરાવર્તનથી વ્યાખ્યાયિત
થતાં વિધેયો---જ બધા સંગણનીય સંખ્યા-સિદ્ધાંતના વિધેયો નથી. આદિમ
પુનરાવર્તનના અનેક સામાન્યકરણો વિચારવામાં આવ્યા છે. પરંતુ વિધેયોનું
સંગણન કરવા માટે સૌથી શક્તિશાળી અને વ્યાપકપણે સ્વીકારાયેલી વધારાની રીત
અનિયત મર્યાદા સુધીની શોધ છે. તેનાથી \emph{આંશિક પુનરાવર્તી વિધેયો}ની
વ્યાખ્યા મળે છે અને તેની સાથે સંકળાયેલી \emph{સામાન્ય પુનરાવર્તી
વિધેયો}ની વ્યાખ્યા પણ મળે છે. સામાન્ય પુનરાવર્તી વિધેયો સંગણનીય છે
અને દરેક ઇનપુટ માટે વ્યાખ્યાયિત છે; તેમની વ્યાખ્યા એવા આંશિક
પુનરાવર્તી વિધેયોને ચોક્કસ રીતે ઓળખાવે છે જે વાસ્તવમાં દરેક ઇનપુટ માટે
વ્યાખ્યાયિત હોય. પુનરાવર્તી વિધેયો સંગણનનાં બીજા દરેક નિદર્શનું
(ટ્યુરિંગ યંત્રો, લૅમ્બડા કલન વગેરેનું) અનુકરણ કરી શકે છે. તેથી તેઓ
સંગણનનાં સ્વીકૃત અનેક નિદર્શોમાંનું એક નિદર્શ છે.
% END OLP-0210
\endgroup


\begingroup
\makeatletter\def\ol@sectioncs{section}\makeatother

% BEGIN OLP-0211
\olfileid{cmp}{rec}{pre}
\olsection{આદિમ પુનરાવર્તન}
\phantomsection\label{guunit:OLP-0211}


પ્રાકૃતિક સંખ્યાઓની એક ખાસિયત એ છે કે $0$થી શરૂ કરીને ઉત્તરગામી
ક્રિયા~$+1$ સાન્ત વખત કરતાં દરેક પ્રાકૃતિક સંખ્યા સુધી પહોંચી શકાય
છે---દરેક પ્રાકૃતિક સંખ્યા કાં તો~$0$ છે અથવા~$0$ના ઉત્તરગામીના
\dots{} ઉત્તરગામી છે. આ હકીકતના આધારે $h\colon\Nat \to \Nat$
વિધેય આપવાની એક રીત આવી છે: (a)~આર્ગ્યુમેન્ટ~$0$ માટે $h$નું મૂલ્ય
જણાવો, અને (b)~$h(x)$નું મૂલ્ય જાણ્યા પછી $h(x+1)$નું મૂલ્ય કેવી રીતે
ગણવું તે પણ જણાવો. (a) સીધું $h(0)$ આપે છે, તેથી $h$નું મૂલ્ય~$0$ માટે
નક્કી થાય છે. હવે (b)ની સૂચના $x=0$ માટે વાપરીને~$h(0)$માંથી
$h(1) = h(0+1)$ ગણી શકાય. એ જ સૂચના $x=1$ માટે વાપરીને~$h(1)$માંથી
$h(2) = h(1+1)$ ગણાય, અને આમ આગળ વધીએ. દરેક પ્રાકૃતિક સંખ્યા~$x$
માટે અંતે એવો તબક્કો આવે છે જેમાં $h(x)$માંથી $h(x+1)$ નક્કી કરીએ
છીએ. તેથી દરેક $x \in \Nat$ માટે $h(x)$ વ્યાખ્યાયિત થાય છે.

\sourcecorrection{OLCMP-001}{સ્થિર અંગ્રેજી મૂળના આ વાક્યમાં
``$h(x+1)$માંથી $h(x)$ નક્કી કરીએ'' એવી દિશા ઊલટી લખાઈ છે.
અગાઉ આપેલી (b) સૂચના મુજબ $h(x)$માંથી $h(x+1)$ નક્કી થાય છે;
ગુજરાતી વાક્યમાં આ દિશા સુધારી છે.}

દાખલા તરીકે, ધારો કે નીચેના બે સમીકરણોથી $h\colon \Nat \to \Nat$
આપ્યું છે:
\begin{align*}
h(0) & =  1\\
h(x+1) & =  2 \cdot h(x)
\end{align*}
જો ગુણાકાર કેવી રીતે કરવો તે પહેલેથી જાણતા હોઈએ, તો આ સમીકરણો ઉપરની
(a) અને~(b) માટે જરૂરી માહિતી આપે છે. બીજું સમીકરણ વારંવાર વાપરતાં
મળે છે કે
\begin{align*}
  h(1) & = 2\cdot h(0) = 2,\\
  h(2) & = 2\cdot h(1) = 2\cdot 2,\\
  h(3) & = 2 \cdot h(2) = 2\cdot 2 \cdot 2,\\
  & \vdots
\end{align*}
આથી આપણે આપેલું વિધેય~$h$ એટલે $h(x) = 2^x$ છે.

પ્રાકૃતિક સંખ્યાઓની ઉપર જણાવેલી ખાસિયત ખાતરી આપે છે કે આ બે શરતો
પૂરી કરતું $h$ વિધેય એક જ છે. આવા બે સમીકરણોની જોડીને $h$ની
\emph{આદિમ પુનરાવર્તન દ્વારા વ્યાખ્યા} કહે છે. આ રીતે $h$ને
નક્કી કરીએ છીએ. તેને ``પુનરાવર્તી''
કહેવાય છે, કારણ કે વ્યાખ્યામાં---ખાસ કરીને બીજા સમીકરણની જમણી
બાજુએ---$h$ પોતે આવે છે. તેને ``આદિમ'' કહેવાય છે, કારણ કે
$h(x+1)$ નક્કી કરતી વખતે માત્ર તરત પહેલાંનું મૂલ્ય~$h(x)$ વાપરીએ
છીએ. $\Nat$ પર વિધેયને પુનરાવર્તી રીતે વ્યાખ્યાયિત કરવાની આ સૌથી
સરળ રીત છે.

સરવાળો અને ગુણાકાર જેવા વધુ મૂળભૂત વિધેયો પણ આદિમ પુનરાવર્તનથી
વ્યાખ્યાયિત કરી શકાય છે. જોકે અહીં સંબંધિત વિધેયોને $2$ આર્ગ્યુમેન્ટ
છે. આપણે એક આર્ગ્યુમેન્ટનું સ્થાન નિશ્ચિત રાખીએ અને બીજા પર
પુનરાવર્તન કરીએ. દાખલા તરીકે, $\Add(x, y)$ વ્યાખ્યાયિત કરવા~$x$
નિશ્ચિત રાખીને પહેલાં $y=0$ માટે મૂલ્ય આપીએ અને પછી~$y$ માટેના
મૂલ્ય પરથી $y+1$ માટેનું મૂલ્ય આપીએ. $x$ નિશ્ચિત છે, તેથી તે
વ્યાખ્યાકારક સમીકરણોની ડાબી અને જમણી બંને બાજુ આવે છે.
\begin{align*}
\Add(x,0) & =  x\\
\Add(x,y+1) & =  \Add(x,y)+1
\end{align*}
આ સમીકરણો દરેક $x$ \emph{અને}~$y$ માટે $\Add$નું મૂલ્ય આપે છે.
દાખલા તરીકે $\Add(2,3)$ શોધવા $x = 2$ રાખીને પહેલાં પ્રથમ સમીકરણથી
$\Add(2,0) = 2$ મળે છે. પછી બીજા સમીકરણને ક્રમે વાપરતાં
$\Add(2,1) = 2 + 1 = 3$, $\Add(2, 2) = 3 + 1 = 4$ અને $\Add(2,
3) = 4 + 1 = 5$ મળે છે.

$\Add$ની વ્યાખ્યાના બીજા સમીકરણની જમણી બાજુએ આપણે $+$ વાપર્યું,
પણ તે માત્ર~$1$ ઉમેરવા માટે છે. બીજા શબ્દોમાં, ઉત્તરગામી વિધેય
$\Succ(z) = z+1$ને અગાઉના મૂલ્ય $\Add(x,y)$ પર લાગુ કરીને
$\Add(x,y+1)$ આપ્યું છે. તેથી આ પુનરાવર્તી વ્યાખ્યાને અગાઉના
મૂલ્ય પર લાગુ થતા એક વિધેય દ્વારા આપવામાં આવેલી માની શકાય. છતાં
એ વિધેય માત્ર અગાઉના મૂલ્ય પર જ નહીં, પણ $x$ અને~$y$ પર પણ આધાર
રાખે એવી છૂટ આપવાથી કોઈ અડચણ નથી---અને ક્યારેક એવી છૂટ જરૂરી છે.
આ જુઓ:
\begin{align*}
  \Mult(x,0) & =  0 \\
  \Mult(x,y+1) & =  \Add(\Mult(x,y),x)
\end{align*}
અહીં અગાઉના મૂલ્ય $\Mult(x,y)$ અને પ્રથમ આર્ગ્યુમેન્ટ~$x$ પર
$\Add$ લાગુ કરીને $\Mult$ની આદિમ પુનરાવર્તી વ્યાખ્યા આપી છે.
આ વ્યાખ્યા દરેક $x$ અને~$y$ માટે $\Mult(x,y)$ નક્કી કરે છે.
દાખલા તરીકે $\Mult(2,3)$, ક્રમે $\Mult(2,0)$, $\Mult(2,1)$,
$\Mult(2,2)$ અને~$\Mult(2,3)$ ગણીને મળે છે:
\begin{align*}
  \Mult(2,0) & = 0\\
  \Mult(2,1) & = \Mult(2,0+1) =
  \Add(\Mult(2,0), 2) = \Add(0, 2) = 2\\
  \Mult(2,2) & = \Mult(2,1+1) =
  \Add(\Mult(2,1), 2) = \Add(2, 2) = 4\\
  \Mult(2,3) & = \Mult(2,2+1) =
  \Add(\Mult(2,2), 2) = \Add(4, 2) = 6
\end{align*}

સામાન્ય માળખું આ પ્રમાણે છે. $h(x_0, \dots, x_{k-1}, y)$ની આદિમ
પુનરાવર્તી વ્યાખ્યા આપવા બે સમીકરણો જોઈએ. પહેલું સમીકરણ~$h$નો
ઉલ્લેખ કર્યા વિના $h(x_0, \dots, x_{k-1}, 0)$નું મૂલ્ય આપે છે.
બીજું સમીકરણ $h(x_0, \dots, x_{k-1}, y)$, અન્ય આર્ગ્યુમેન્ટ
$x_0$, \dots,~$x_{k-1}$ અને~$y$ના આધારે
$h(x_0, \dots, x_{k-1}, y+1)$નું મૂલ્ય આપે છે. બીજા સમીકરણમાં
$h$નું માત્ર તરત પહેલાંનું મૂલ્ય જ વાપરી શકાય. જો બંને સમીકરણોની
જમણી બાજુએ આવતી ક્રિયાઓને અનુક્રમે $f$ અને~$g$ વિધેયો માનીએ,
તો આદિમ પુનરાવર્તનથી નવું વિધેય~$h$ વ્યાખ્યાયિત કરવાની સામાન્ય
પદ્ધતિ આ છે:
\begin{align*}
  h(x_0, \dots, x_{k-1}, 0) & = f(x_0, \dots, x_{k-1})\\
  h(x_0, \dots, x_{k-1}, y+1) & =
  g(x_0, \dots, x_{k-1}, y, h(x_0, \dots, x_{k-1}, y))
\end{align*}
$\Add$ માટે $k=1$, $f(x_0) = x_0$ (તદેવ વિધેય) અને
$g(x_0, y, z) = z + 1$ છે (પોતાના ત્રીજા આર્ગ્યુમેન્ટનો ઉત્તરગામી
આપતું $3$ આર્ગ્યુમેન્ટવાળું વિધેય):
\begin{align*}
  \Add(x_0, 0) & = f(x_0) = x_0\\
  \Add(x_0, y+1) & = g(x_0, y, \Add(x_0, y)) =
  \Succ(\Add(x_0, y))
\end{align*}
$\Mult$ માટે $f(x_0) = 0$ (હંમેશાં~$0$ આપતું અચળ વિધેય) અને
$g(x_0, y, z) = \Add(z,x_0)$ છે (પોતાના છેલ્લા અને પ્રથમ
આર્ગ્યુમેન્ટનો સરવાળો આપતું $3$ આર્ગ્યુમેન્ટવાળું વિધેય):
\begin{align*}
  \Mult(x_0,0) & =  f(x_0) = 0 \\
  \Mult(x_0,y+1) & =  g(x_0, y, \Mult(x_0,y)) =
  \Add(\Mult(x_0,y), x_0)
\end{align*}
% END OLP-0211
\endgroup


\begingroup
\makeatletter\def\ol@sectioncs{section}\makeatother

% BEGIN OLP-0212
\olfileid{cmp}{rec}{com}
\olsection{સંયોજન}
\phantomsection\label{guunit:OLP-0212}


પ્રાકૃતિક સંખ્યાઓનાં એક-આર્ગ્યુમેન્ટવાળાં બે વિધેયો $f$ અને $g$ હોય,
તો તેમનું સંયોજન કરી શકાય: $h(x) = f(g(x))$. આમ નવું વિધેય~$h(x)$,
$f$ અને~$g$ના \emph{સંયોજન}થી વ્યાખ્યાયિત થાય છે. હવે આ રીતને
એકથી વધુ આર્ગ્યુમેન્ટ ધરાવતાં વિધેયો માટે સામાન્ય બનાવવા માગીએ છીએ.

તેની એક રીત આ છે: ધારો કે $f$ને $k$ આર્ગ્યુમેન્ટ છે અને
$g_0$, \dots, $g_{k-1}$ એવાં $k$ વિધેયો છે, જેમાંના દરેકને
$n$ આર્ગ્યુમેન્ટ છે. તો નવું $n$-આર્ગ્યુમેન્ટવાળું વિધેય $h$ આ રીતે
વ્યાખ્યાયિત કરી શકાય:
\[
h(x_0, \dots, x_{n-1}) =
f(g_0(x_0, \dots, x_{n-1}), \dots, g_{k-1}(x_0, \dots, x_{n-1}))
\]
$f$ અને દરેક $g_i$ સંગણનીય હોય તો $h$ પણ સંગણનીય છે:
$h(x_0, \dots, x_{n-1})$ ગણવા માટે પહેલાં દરેક $i = 0$,
\dots,~$k-1$ માટે $y_i = g_i(x_0, \dots, x_{n-1})$ ગણો.
પછી આ મૂલ્યો $f$માં આપીને
$h(x_0, \dots, x_{n-1}) = f(y_0, \dots, y_{k-1})$ ગણો.

\sourcecorrection{OLCMP-002}{સ્થિર અંગ્રેજી મૂળના આ વાક્યમાં
$h(x_0,\dots,x_{k-1})$ લખાયું છે. અહીં $h$ને $n$ આર્ગ્યુમેન્ટ છે;
અતઃ ડાબી બાજુનો અંતિમ સૂચક $n-1$ હોવો જોઈએ. ગુજરાતી લખાણમાં
આ જ સૂત્ર $h(x_0,\dots,x_{n-1})=f(y_0,\dots,y_{k-1})$ રાખ્યું છે.}

પહેલેથી મળેલાં વિધેયો વાપરીને નવું વિધેય ગણીએ ત્યારે શું થાય છે તેનું
આ વર્ણન કદાચ વધારે પડતું મર્યાદિત લાગે. એક કારણ એ છે કે ક્યારેક
આપણે વિધેયના બધા આર્ગ્યુમેન્ટ વાપરતા નથી: $\Add$ની આદિમ
પુનરાવર્તી વ્યાખ્યા માટે $g(x, y, z) = \Succ(z)$ આપ્યું ત્યારે પણ
એવું જ થયું હતું. ધારો કે નીચેના વિધેયો વાપરવાની છૂટ છે:
\[
\Proj{n}{i}(x_0, \dots, x_{n-1}) = x_i
\]
$\Proj{k}{i}$ને \emph{પ્રક્ષેપ} વિધેયો કહે છે:
$\Proj{n}{i}$ને $n$ આર્ગ્યુમેન્ટ છે. હવે $g$ને આ રીતે
વ્યાખ્યાયિત કરી શકાય:
\[
g(x, y, z) = \Succ(\Proj{3}{2}(x, y, z)).
\]
અહીં $f$ની ભૂમિકા $1$-આર્ગ્યુમેન્ટવાળું વિધેય $\Succ$ ભજવે છે,
તેથી $k=1$. અને $g_0$ની ભૂમિકા એક $3$-આર્ગ્યુમેન્ટવાળું વિધેય
$\Proj{3}{2}$ ભજવે છે. પરિણામે મળતું $3$-આર્ગ્યુમેન્ટવાળું વિધેય
તેના ત્રીજા આર્ગ્યુમેન્ટનો ઉત્તરગામી આપે છે.

પ્રક્ષેપ વિધેયો આર્ગ્યુમેન્ટોનો ક્રમ બદલવા અથવા બે આર્ગ્યુમેન્ટોને
એક જ આર્ગ્યુમેન્ટથી ભરવા વડે પણ નવાં વિધેયો વ્યાખ્યાયિત કરવા દે છે.
દાખલા તરીકે, $h(x) = \Add(x, x)$ને આ રીતે વ્યાખ્યાયિત કરી શકાય:
\[
h(x_0) = \Add(\Proj{1}{0}(x_0),\Proj{1}{0}(x_0)).
\]
અહીં $k=2$, $n=1$ છે. $f(y_0,y_1)$ની ભૂમિકા $\Add$ ભજવે છે;
$g_0(x_0)$ અને $g_1(x_0)$ બંનેની ભૂમિકા~$\Proj{1}{0}(x_0)$,
એટલે કે એક-આર્ગ્યુમેન્ટવાળું પ્રક્ષેપ વિધેય (તદેવ વિધેય), ભજવે છે.

જો $f(y_0, y_1)$ વિધેય પહેલેથી આપણી પાસે હોય, તો
$h(x_0, x_1) = f(x_1, x_0)$ને આ રીતે વ્યાખ્યાયિત કરી શકાય:
\[
h(x_0, x_1) = f(\Proj{2}{1}(x_0, x_1),\Proj{2}{0}(x_0, x_1)).
\]
અહીં $k=2$, $n = 2$ છે, અને $g_0$ તથા $g_1$ની ભૂમિકા અનુક્રમે
$\Proj{2}{1}$ અને~$\Proj{2}{0}$ ભજવે છે.

$g_0$, \dots,~$g_{k-1}$ બધાંને એકસરખા $n$ આર્ગ્યુમેન્ટ હોવા
જ જોઈએ, એ શરત અંગે પણ શંકા થઈ શકે. (યાદ રાખો કે વિધેયની
\emph{આર્ગ્યુમેન્ટ-સંખ્યા} એટલે તેના આર્ગ્યુમેન્ટોની સંખ્યા;
$n$-આર્ગ્યુમેન્ટવાળાં વિધેયની આ સંખ્યા~$n$ છે.) પરંતુ પ્રક્ષેપ
વિધેયો ઉમેરવાથી જરૂરી લવચીકતા મળે છે. દાખલા તરીકે, ધારો કે $f$
અને~$g$ $3$-આર્ગ્યુમેન્ટવાળાં વિધેયો છે, અને $h$ નીચે પ્રમાણે
વ્યાખ્યાયિત $2$-આર્ગ્યુમેન્ટવાળું વિધેય છે:
\[
h(x,y) = f(x,g(x,x,y),y).
\]
પ્રક્ષેપ વિધેયો વડે $h$ની વ્યાખ્યા આ રીતે ફરી લખી શકાય:
\[
h(x,y) = f(\Proj{2}{0}(x,y), g(\Proj{2}{0}(x,y), \Proj{2}{0}(x,y),
\Proj{2}{1}(x,y)), \Proj{2}{1}(x,y)).
\]
આથી $h$, $f$નું $\Proj{2}{0}$, $l$ અને $\Proj{2}{1}$ સાથેનું
સંયોજન છે, જ્યાં
\[
l(x,y) = g(\Proj{2}{0}(x,y),\Proj{2}{0}(x,y),\Proj{2}{1}(x,y)),
\]
એટલે કે $l$, $g$નું $\Proj{2}{0}$, $\Proj{2}{0}$ અને
~$\Proj{2}{1}$ સાથેનું સંયોજન છે.
% END OLP-0212
\endgroup


\begingroup
\makeatletter\def\ol@sectioncs{section}\makeatother

% BEGIN OLP-0213
\olfileid{cmp}{rec}{prf}
\olsection{આદિમ પુનરાવર્તી વિધેયો}
\phantomsection\label{guunit:OLP-0213}


આદિમ પુનરાવર્તન અને સંયોજન વડે પહેલેથી મળેલા વિધેયોમાંથી નવાં વિધેયો
કેવી રીતે વ્યાખ્યાયિત થાય છે, તે ફરી નોંધીએ.

\begin{defn}
  \ollabel{defn:primitive-recursion} ધારો કે $f$ એ $k$-આર્ગ્યુમેન્ટવાળું
  વિધેય છે ($k\ge 1$) અને $g$ એ $(k+2)$-આર્ગ્યુમેન્ટવાળું વિધેય છે.
  $f$ અને $g$માંથી \emph{આદિમ પુનરાવર્તન દ્વારા વ્યાખ્યાયિત વિધેય}
  એટલે નીચેના સમીકરણોથી વ્યાખ્યાયિત $(k+1)$-આર્ગ્યુમેન્ટવાળું
  વિધેય~$h$:
  \begin{align*}
    h(x_0,\dots,x_{k-1},0) & =  f(x_0,\dots,x_{k-1}) \\
    h(x_0,\dots,x_{k-1},y+1) & = g(x_0,\dots,x_{k-1}, y, h(x_0,\dots,x_{k-1}, y))
  \end{align*}
\end{defn}

\begin{defn}
  \ollabel{defn:composition} ધારો કે $f$ એ $k$-આર્ગ્યુમેન્ટવાળું
  વિધેય છે, અને $g_0$, \dots, $g_{k-1}$ એ બધાં
  $n$-આર્ગ્યુમેન્ટવાળાં એવાં $k$ વિધેયો છે. $f$ તથા
  $g_0$, \dots,~$g_{k-1}$માંથી \emph{સંયોજન દ્વારા વ્યાખ્યાયિત વિધેય}
  એટલે નીચેના સમીકરણથી વ્યાખ્યાયિત $n$-આર્ગ્યુમેન્ટવાળું વિધેય~$h$:
  \[
  h(x_0, \dots, x_{n-1}) =
  f(g_0(x_0, \dots, x_{n-1}), \dots, g_{k-1}(x_0, \dots, x_{n-1})).
  \]
\end{defn}

$\Succ$ અને દરેક પ્રાકૃતિક સંખ્યા $n$ તથા $i < n$ માટેનાં પ્રક્ષેપ
વિધેયો
\[
\Proj{n}{i}(x_0,\dots,x_{n-1}) = x_i,
\]
ઉપરાંત, આદિમ પુનરાવર્તી વિધેયોમાં આપણે $\Zero(x) = 0$ વિધેયનો પણ
સમાવેશ કરીશું.

\begin{defn}
  આદિમ પુનરાવર્તી વિધેયોનો ગણ એટલે $\Nat^n$માંથી $\Nat$માં જતાં
  વિધેયોનો એવો ગણ, જેને નીચેની શરતોથી આગમનાત્મક રીતે વ્યાખ્યાયિત
  કરીએ છીએ:
  \begin{enumerate}
  \item $\Zero$ આદિમ પુનરાવર્તી છે.
  \item $\Succ$ આદિમ પુનરાવર્તી છે.
  \item દરેક પ્રક્ષેપ વિધેય $\Proj{n}{i}$ આદિમ પુનરાવર્તી છે.
  \item જો $f$ એ $k$-આર્ગ્યુમેન્ટવાળું આદિમ પુનરાવર્તી વિધેય હોય અને
    $g_0$, \dots,~$g_{k-1}$ એ $n$-આર્ગ્યુમેન્ટવાળાં આદિમ
    પુનરાવર્તી વિધેયો હોય, તો $g_0$, \dots,~$g_{k-1}$ સાથે $f$નું
    સંયોજન પણ આદિમ પુનરાવર્તી છે.
  \item જો $f$ એ $k$-આર્ગ્યુમેન્ટવાળું આદિમ પુનરાવર્તી વિધેય હોય અને
    $g$ એ $k+2$-આર્ગ્યુમેન્ટવાળું આદિમ પુનરાવર્તી વિધેય હોય, તો
    $f$ અને $g$માંથી આદિમ પુનરાવર્તન દ્વારા વ્યાખ્યાયિત વિધેય પણ
    આદિમ પુનરાવર્તી છે.
\end{enumerate}
\end{defn}

\begin{explain}
વધુ ટૂંકમાં, આદિમ પુનરાવર્તી વિધેયોનો ગણ એવો સૌથી નાનો ગણ છે જેમાં
$\Zero$, $\Succ$ અને પ્રક્ષેપ વિધેયો~$\Proj{n}{j}$ છે અને જે
સંયોજન તથા આદિમ પુનરાવર્તન હેઠળ બંધ છે.

આદિમ પુનરાવર્તી વિધેયોનો ગણ ``તબક્કાઓ'' વડે પણ વર્ણવી શકાય. $S_0$
એ શરૂઆતનાં વિધેયોનો ગણ દર્શાવે: $\Zero$, $\Succ$ અને પ્રક્ષેપ
વિધેયો. આ તબક્કા~$0$નાં આદિમ પુનરાવર્તી વિધેયો છે. તબક્કો $S_i$
વ્યાખ્યાયિત થયા પછી, $S_i$માંનાં વિધેયો પર સંયોજન અથવા આદિમ
પુનરાવર્તનનો એક વાર ઉપયોગ કરીને મળતાં બધાં વિધેયોનો ગણ $S_{i+1}$
માનીએ. ત્યારે
\[
S = \bigcup_{i \in \Nat} S_i
\]
એ બધાં આદિમ પુનરાવર્તી વિધેયોનો ગણ છે.
\end{explain}

હવે તપાસીએ કે $\Add$ આદિમ પુનરાવર્તી વિધેય છે.

\begin{prop}
  સરવાળાનું વિધેય $\Add(x,y) = x+y$ આદિમ પુનરાવર્તી છે.
\end{prop}

\begin{proof}
આપણે $f$ અને~$g$ એવાં બે વિધેયોના આધારે $\Add$ની આદિમ પુનરાવર્તી
વ્યાખ્યા પહેલેથી આપી છે, જે \olref{defn:primitive-recursion}ના
માળખાને અનુસરે છે:
\begin{align*}
  \Add(x_0, 0) & = f(x_0) = x_0\\
  \Add(x_0, y+1) & = g(x_0, y, \Add(x_0, y)) =
  \Succ(\Add(x_0, y))
\end{align*}
તેથી $f$ અને $g$ પણ આદિમ પુનરાવર્તી હોય તો $\Add$ આદિમ
પુનરાવર્તી છે. $f(x_0) = x_0 = \Proj{1}{0}(x_0)$ છે અને પ્રક્ષેપ
વિધેયોને આદિમ પુનરાવર્તી ગણીએ છીએ, તેથી $f$ આદિમ પુનરાવર્તી છે.
$g$ એ $g(x_0, y, z)$ એવાં ત્રણ આર્ગ્યુમેન્ટવાળું વિધેય છે, જે આ
રીતે વ્યાખ્યાયિત છે:
\[
g(x_0, y, z) = \Succ(z).
\]
આટલાથી હજી $g$ આદિમ પુનરાવર્તી છે એવું સાબિત થતું નથી, કારણ કે
$g$ અને $\Succ$ એકસરખાં વિધેયો નથી: $\Succ$ને એક આર્ગ્યુમેન્ટ છે,
જ્યારે $g$ને ત્રણ આર્ગ્યુમેન્ટ હોવા જોઈએ. પરંતુ સંયોજન દ્વારા
$g$ને ઔપચારિક રીતે આમ વ્યાખ્યાયિત કરી શકીએ:
\[
g(x_0, y, z) = \Succ(\Proj{3}{2}(x_0, y, z))
\]
$\Succ$ અને $\Proj{3}{2}$ને આદિમ પુનરાવર્તી વિધેયો ગણીએ છીએ;
તેમના સંયોજનથી આ વિધેય વ્યાખ્યાયિત થાય છે, તેથી $g$ પણ આદિમ
પુનરાવર્તી છે.
\end{proof}

\begin{prop}
  \ollabel{prop:mult-pr}
  ગુણાકારનું વિધેય $\Mult(x,y) = x \cdot y$ આદિમ પુનરાવર્તી છે.
\end{prop}

\begin{proof}
  અભ્યાસ માટે.
\end{proof}

\begin{prob}
  $\Mult$ની આદિમ પુનરાવર્તી વ્યાખ્યા
  \olref[cmp][rec][prf]{defn:primitive-recursion}માં જરૂરી માળખામાં
મૂકી શકાય છે અને સંબંધિત વિધેયો $f$ તથા~$g$ આદિમ પુનરાવર્તી છે,
એવું બતાવીને \olref[cmp][rec][prf]{prop:mult-pr} સાબિત કરો.
\end{prob}

\begin{ex}
આદિમ પુનરાવર્તી વ્યાખ્યાનું આપણું સૌપ્રથમ ઉદાહરણ આ હતું:
\begin{align*}
h(0) & =  1 \\
h(y+1) & =  2 \cdot h(y).
\end{align*}
આ વિધેય \olref{defn:primitive-recursion}માં જરૂરી માળખામાં બંધબેસતું
નથી, કારણ કે અહીં $k=0$ છે. વ્યાખ્યામાં અચળ મૂલ્યો $1$ અને $2$ પણ
આવે છે. પહેલી મુશ્કેલી દૂર કરવા એક બિનવપરાયેલું આર્ગ્યુમેન્ટ ઉમેરીને
વિધેય~$h'$ વ્યાખ્યાયિત કરીએ:
\begin{align*}
h'(x_0, 0) & =  f(x_0) = 1 \\
h'(x_0, y+1) & =  g(x_0, y, h'(x_0, y)) = 2 \cdot h'(x_0, y).
\end{align*}
સંયોજન વડે $\Succ$ અને $\Zero$માંથી $f(x_0) = 1$ વ્યાખ્યાયિત
કરી શકાય: $f(x_0) = \Succ(\Zero(x_0))$. $g'(z) = 2 \cdot z$
અને પ્રક્ષેપ વિધેયોમાંથી સંયોજન વડે $g$ વ્યાખ્યાયિત કરી શકાય:
\begin{align*}
g(x_0, y, z) & = g'(\Proj{3}{2}(x_0, y, z))
\intertext{અને $g'$ને પણ સંયોજન વડે આમ વ્યાખ્યાયિત કરી શકાય:}
g'(z) & = \Mult(g''(z), \Proj{1}{0}(z))
\intertext{તથા}
g''(z) & = \Succ(f(z)),
\end{align*}
અહીં $f$ ઉપર મુજબ છે: $f(z) = \Succ(\Zero(z))$. હવે~$h'$
મળી ગયું છે, એટલે ફરી સંયોજન વાપરીને $h(y) =
h'(\Proj{1}{0}(y),\Proj{1}{0}(y))$ રાખી શકીએ. આથી $h$ને
મૂળભૂત વિધેયોમાંથી સંયોજનો અને આદિમ પુનરાવર્તનોની શ્રેણી વડે
વ્યાખ્યાયિત કરી શકાય છે; એટલે $h$ આદિમ પુનરાવર્તી છે.
\end{ex}
% END OLP-0213
\endgroup


\begingroup
\makeatletter\def\ol@sectioncs{section}\makeatother

% BEGIN OLP-0214
\olfileid{cmp}{rec}{not}
\olsection{આદિમ પુનરાવર્તી વિધેયોનાં સંકેતન}
\phantomsection\label{guunit:OLP-0214}


આદિમ પુનરાવર્તી વિધેયોનું ચોક્કસ આગમનાત્મક વર્ણન આપ્યું છે તેનો એક
ફાયદો એ છે કે આપણે તેમને વ્યવસ્થિત રીતે દર્શાવી શકીએ છીએ.
દાખલા તરીકે, આવા દરેક વિધેયને નીચે મુજબ એક ``સંકેતન'' આપી શકીએ.
શૂન્ય વિધેય, ઉત્તરગામી વિધેય અને પ્રક્ષેપ વિધેયો માટે અનુક્રમે
$\Zero$, $\Succ$ અને $\Proj{n}{i}$ સંકેતો વાપરીએ. હવે ધારો કે
$k$-આર્ગ્યુમેન્ટવાળું વિધેય~$f$ અને $n$-આર્ગ્યુમેન્ટવાળાં વિધેયો
$g_0$, \dots,~$g_{k-1}$ના સંયોજનથી $h$ વ્યાખ્યાયિત છે, અને આ
વિધેયોને અનુક્રમે $F$, $G_0$, \dots,~$G_{k-1}$ સંકેતનો આપ્યાં છે.
તો નવો સંકેત $\fn{Comp}_{k,n}$ વાપરીને $h$ વિધેયને
$\fn{Comp}_{k,n}[F,G_0,\dots,G_{k-1}]$ વડે દર્શાવી શકીએ.

આદિમ પુનરાવર્તનથી વ્યાખ્યાયિત વિધેયો માટે પણ આવાં જ સંકેતનો
વાપરી શકીએ. ધારો કે $(k+1)$-આર્ગ્યુમેન્ટવાળું વિધેય~$h$,
$k$-આર્ગ્યુમેન્ટવાળા વિધેય~$f$ અને $(k+2)$-આર્ગ્યુમેન્ટવાળા
વિધેય~$g$માંથી આદિમ પુનરાવર્તન દ્વારા વ્યાખ્યાયિત છે; અને $f$ તથા
$g$ને અપાયેલાં સંકેતનો અનુક્રમે $F$ અને~$G$ છે. ત્યારે~$h$ને
અપાતું સંકેતન $\fn{Rec}_k[F,G]$ છે.

યાદ કરો કે સરવાળાનું વિધેય આદિમ પુનરાવર્તનથી આમ વ્યાખ્યાયિત છે:
\begin{align*}
  \Add(x_0, 0) & = \Proj{1}{0}(x_0) = x_0\\
  \Add(x_0, y+1) & = \Succ(\Proj{3}{2}(x_0, y, \Add(x_0, y))) = \Add(x_0, y) +1
\end{align*}
અહીં $f$નું સ્થાન $\Proj{1}{0}$ લે છે અને $g$નું સ્થાન
$\Succ(\Proj{3}{2}(x_0, y, z))$ લે છે. પછીના વિધેયને
$\fn{Comp}_{1,3}[\Succ,\Proj{3}{2}]$ સંકેતન મળે છે, કારણ કે તે
$1$-આર્ગ્યુમેન્ટવાળા વિધેય~$\Succ$ અને
$3$-આર્ગ્યુમેન્ટવાળા વિધેય~$\Proj{3}{2}$ના સંયોજનથી વ્યાખ્યાયિત
થાય છે. આ રીતે સરવાળાનું વિધેય નીચેના સંકેતનથી દર્શાવી શકીએ:
\[
\fn{Rec}_1[\Proj{1}{0},\fn{Comp}_{1,3}[\Succ,\Proj{3}{2}]].
\]
આવાં સંકેતનો ક્યારેક ઉપયોગી થાય છે; દાખલા તરીકે, આદિમ પુનરાવર્તી
વિધેયોની પરિગણના કરતી વખતે.

\begin{prob}
  $\Mult$ માટે સંપૂર્ણ આદિમ પુનરાવર્તી સંકેતન આપો.
\end{prob}
% END OLP-0214
\endgroup


\begingroup
\makeatletter\def\ol@sectioncs{section}\makeatother

% BEGIN OLP-0215
\olfileid{cmp}{rec}{cmp}
\olsection{આદિમ પુનરાવર્તી વિધેયો સંગણનીય છે}
\phantomsection\label{guunit:OLP-0215}


ધારો કે વિધેય $h$ આદિમ પુનરાવર્તનથી આ રીતે વ્યાખ્યાયિત છે:
\begin{eqnarray*}
h(\vec x, 0)     & = & f(\vec x) \\
h(\vec x, y + 1) & = & g(\vec x, y, h(\vec x, y))
\end{eqnarray*}
અને ધારો કે $f$ તથા $g$ સંગણનીય વિધેયો છે. (અહીં $x_0$,
\dots, $x_{k-1}$ના સંક્ષેપ તરીકે $\vec x$ વાપરીએ છીએ.) ત્યારે
$h(\vec x, 0)$ સ્પષ્ટ રીતે ગણી શકાય છે, કારણ કે તે માત્ર
$f(\vec x)$ છે, જે સંગણનીય છે એવું ધાર્યું છે. પછી $h(\vec x, 1)$
પણ ગણી શકાય: $1 = 0 + 1$ હોવાથી $h(\vec x, 1)$ માત્ર નીચેનું મૂલ્ય છે:
\begin{align*}
 h(\vec x, 1) & = g(\vec x, 0, h(\vec x, 0)) =  g(\vec x, 0, f(\vec x)).
\intertext{આ જ રીતે આગળ વધીને આપણે ગણી શકીએ છીએ:}
h(\vec x, 2) & = g(\vec x, 1, h(\vec x, 1)) = g(\vec x, 1, g(\vec x, 0, f(\vec x)))\\
h(\vec x, 3) & = g(\vec x, 2, h(\vec x, 2)) = g(\vec x, 2, g(\vec x, 1, g(\vec x, 0, f(\vec x))))\\
h(\vec x, 4) & = g(\vec x, 3, h(\vec x, 3)) = g(\vec x, 3, g(\vec x, 2, g(\vec x, 1, g(\vec x, 0, f(\vec x)))))\\
& \vdots
\end{align*}
તેથી સામાન્ય રીતે $h(\vec x, y)$ ગણવા માટે ક્રમે
$h(\vec x, 0)$, $h(\vec x, 1)$, \dots, ગણીએ અને
$h(\vec x, y)$ સુધી પહોંચીએ.

આથી $f$ અને $g$ સંગણનીય હોય, તો આદિમ પુનરાવર્તી વ્યાખ્યા નવું
સંગણનીય વિધેય આપે છે. એ જ રીતે, $f$ અને $g_i$ સંગણનીય હોય, તો
તેમના સંયોજનથી મળતું વિધેય પણ સંગણનીય છે.

મૂળભૂત વિધેયો $\Zero$, $\Succ$ અને $\Proj{n}{i}$ સંગણનીય છે, તથા
સંગણનીય વિધેયોના સંયોજન અને આદિમ પુનરાવર્તનથી ફરી સંગણનીય વિધેયો
મળે છે. તેથી દરેક આદિમ પુનરાવર્તી વિધેય સંગણનીય છે.
% END OLP-0215
\endgroup


\begingroup
\makeatletter\def\ol@sectioncs{section}\makeatother

% BEGIN OLP-0216
\olfileid{cmp}{rec}{exa}
\olsection{આદિમ પુનરાવર્તી વિધેયોનાં ઉદાહરણો}
\phantomsection\label{guunit:OLP-0216}



આદિમ પુનરાવર્તી વિધેયોનાં કેટલાંક ઉદાહરણો આપણી પાસે છે: સરવાળા
અને ગુણાકારનાં વિધેયો~$\Add$ અને $\Mult$. તદેવ વિધેય
$\fn{id}(x) = x$ આદિમ પુનરાવર્તી છે, કારણ કે તે માત્ર
$\Proj{1}{0}$ છે. અચળ વિધેયો $\fn{const}_n(x) = n$ પણ આદિમ
પુનરાવર્તી છે, કારણ કે $\Zero$ અને $\Succ$ના ક્રમિક સંયોજનથી
તેમને વ્યાખ્યાયિત કરી શકાય છે. આદિમ પુનરાવર્તી વ્યાખ્યાઓમાં અચળ
મૂલ્યો વાપરવા હોય ત્યારે આ ઉપયોગી છે. દાખલા તરીકે,
$f(x) = 2 \cdot x$ વ્યાખ્યાયિત કરવું હોય તો $\fn{const}_n(x)$ અને
ગુણાકારના સંયોજનથી $f(x) =
\Mult(\fn{const}_2(x), \Proj{1}{0}(x))$ મળે છે. હવે પછી આ રીત
વાપરતા રહીશું.

\begin{prop}
  ઘાતાંકનું વિધેય $\fn{exp}(x, y) = x^y$ આદિમ પુનરાવર્તી છે.
\end{prop}

\begin{proof}
  $\fn{exp}$ને આદિમ પુનરાવર્તનથી આમ વ્યાખ્યાયિત કરી શકીએ:
  \begin{align*}
    \fn{exp}(x, 0) & = 1\\
    \fn{exp}(x, y+1) & = \Mult(x, \fn{exp}(x,y)).
    \intertext{ચોક્કસ કહીએ તો, આ આદિમ પુનરાવર્તી વિધેયોમાંથી
      બનેલી પુનરાવર્તી વ્યાખ્યા હજી નથી. ઔપચારિક રીતે, જોકે,
      આપણે આ રીતે લખીએ:}
    \fn{exp}(x, 0) & = f(x)\\
    \fn{exp}(x, y+1) & = g(x, y, \fn{exp}(x,y)).
    \intertext{અહીં}
    f(x) & = \Succ(\Zero(x)) = 1\\
    g(x, y, z) & = \Mult(\Proj{3}{0}(x, y, z), \Proj{3}{2}(x, y, z)) = x \cdot z
  \end{align*}
  તેથી $f$ અને $g$ આદિમ પુનરાવર્તી વિધેયોના સંયોજનથી વ્યાખ્યાયિત
  થાય છે.
\end{proof}

\begin{prop}
  નીચે મુજબ વ્યાખ્યાયિત પૂર્વગામી વિધેય $\fn{pred}(y)$ આદિમ
  પુનરાવર્તી છે:
  \[
  \fn{pred}(y) = \begin{cases}
    0 & \text{જો $y=0$ હોય}\\
    y-1 & \text{અન્યથા}
  \end{cases}
  \]
\end{prop}

\begin{proof}
 નોંધો કે
 \begin{align*}
   \fn{pred}(0) & = 0 \text{ અને}\\
   \fn{pred}(y+1) & = y.
 \end{align*}
 આ લગભગ આદિમ પુનરાવર્તી વ્યાખ્યા છે. ચોક્કસ રીતે જોઈએ તો તે આદિમ
 પુનરાવર્તનની વ્યાખ્યાના માળખામાં બંધબેસતું નથી, કારણ કે એમાં ઓછામાં
 ઓછું એક વધારાનું આર્ગ્યુમેન્ટ~$x$ જરૂરી છે. વળી,
 $\fn{pred}(y+1)$ વ્યાખ્યાયિત કરતાં ખરેખર $\fn{pred}(y)$ વપરાતું જ
 નથી. પરંતુ પહેલાં $\fn{pred}'(x, y)$ને આમ વ્યાખ્યાયિત કરી શકીએ:
 \begin{align*}
   \fn{pred}'(x, 0) & = \Zero(x) = 0,\\
   \fn{pred}'(x, y+1) & = \Proj{3}{1}(x, y, \fn{pred'}(x, y)) = y.
 \end{align*}
અને પછી સંયોજનથી $\fn{pred}$ વ્યાખ્યાયિત કરી શકીએ; દાખલા તરીકે,
$\fn{pred}(x) = \fn{pred}'(\Zero(x), \Proj{1}{0}(x))$.
\end{proof}

\begin{prop}
  ક્રમગુણિતનું વિધેય $\fn{fac}(x) = \fact{x} = 1 \cdot 2 \cdot 3
  \cdot \dots \cdot x$ આદિમ પુનરાવર્તી છે.
\end{prop}

\begin{proof}
  સ્પષ્ટ દેખાતી આદિમ પુનરાવર્તી વ્યાખ્યા આ છે:
  \begin{align*}
    \fn{fac}(0) &= 1\\
    \fn{fac}(y+1) & = \fn{fac}(y) \cdot (y+1).
    \intertext{ઔપચારિક રીતે પહેલાં બે આર્ગ્યુમેન્ટવાળું વિધેય $h$
      વ્યાખ્યાયિત કરવું પડે:}
    h(x, 0) & = \fn{const}_1(x)\\
    h(x, y+1) & = g(x, y, h(x, y))
    \intertext{અહીં $g(x, y, z) = \Mult(\Proj{3}{2}(x, y, z),
      \Succ(\Proj{3}{1}(x, y, z)))$ છે; પછી રાખીએ:}
    \fn{fac}(y) & = h(\Proj{1}{0}(y), \Proj{1}{0}(y)) = h(y,y).
  \end{align*}
  હવે પછી આપણે આ બાબતમાં થોડા ઉદાર રહીશું અને સંયોજન તથા આદિમ
  પુનરાવર્તન વડે મળતી ઔપચારિક વ્યાખ્યાઓ દરેક વખતે આપીશું નહીં.
\end{proof}

\begin{prop}
  શૂન્યે અટકતી બાદબાકી $x \tsub y$, જે આ રીતે વ્યાખ્યાયિત છે,
  આદિમ પુનરાવર્તી છે:
  \[
  x \tsub y = \begin{cases}
    0 & \text{જો $x < y$ હોય}\\
    x-y & \text{અન્યથા}
  \end{cases}
  \]
\end{prop}

\begin{proof}
  આપણી પાસે આ સમીકરણો છે:
  \begin{align*}
    x \tsub 0 & = x\\
    x \tsub (y+1) & = \fn{pred}(x \tsub y)
  \end{align*}
\end{proof}

\begin{prop}
 $x$ અને $y$ વચ્ચેનું અંતર $\left|x-y\right|$ આદિમ પુનરાવર્તી છે.
\end{prop}

\begin{proof}
  $\left| x-y \right| = (x \tsub y) + (y \tsub x)$ છે. તેથી અંતરને
  સરવાળા~$+$ અને શૂન્યે અટકતી બાદબાકી~$\tsub$ના સંયોજનથી
  વ્યાખ્યાયિત કરી શકાય છે; બંને આદિમ પુનરાવર્તી છે.
\end{proof}

\begin{prop}
  $x$ અને $y$માંથી મોટામાં મોટું મૂલ્ય $\fn{max}(x,y)$ આદિમ
  પુનરાવર્તી છે.
\end{prop}

\begin{proof}
  $+$ અને $\tsub$ના સંયોજનથી $\fn{max}(x,y)$ને આમ વ્યાખ્યાયિત
  કરી શકાય:
  \[
  \fn{max}(x,y) \defis x + (y \tsub x).
  \]
  જો $x$ મોટામાં મોટું હોય, એટલે $x \ge y$, તો $y \tsub x = 0$;
  તેથી $x + (y \tsub x) = x + 0 = x$. જો $y$ મોટામાં મોટું હોય,
  તો $y \tsub x = y - x$; તેથી $x + (y \tsub x) = x + (y - x) = y$.
\end{proof}

\begin{prop}
  \ollabel{prop:min-pr}
  $x$ અને $y$માંથી નાનામાં નાનું મૂલ્ય $\fn{min}(x,y)$ આદિમ
  પુનરાવર્તી છે.
\end{prop}

\begin{proof}
  અભ્યાસ માટે.
\end{proof}

\begin{prob}
  \olref[cmp][rec][exa]{prop:min-pr} સાબિત કરો.
\end{prob}

\begin{prob}
નીચેનું વિધેય આદિમ પુનરાવર્તી છે એવું બતાવો: \[
f(x, y) =
2^{(2^{\iddots^{2^{x}}})}\raisebox{1ex}{\bigg\rbrace}
\raisebox{1ex}{\text {$y$ વખત $2$}}\]
\end{prob}

\begin{prob}
પૂર્ણાંક ભાગાકાર $d(x, y) = \lfloor x/y \rfloor$ (એટલે કે ભાગફળના
દશાંશ બિંદુ પછીનો ભાગ અવગણીને મળતો ભાગાકાર) આદિમ પુનરાવર્તી છે
એવું બતાવો. $y = 0$ હોય ત્યારે આપણે $d(x, y) = 0$ નક્કી કરીએ છીએ.
આદિમ પુનરાવર્તન અને સંયોજન વાપરીને~$d$ની સ્પષ્ટ વ્યાખ્યા આપો.
\end{prob}

\begin{prop}
આદિમ પુનરાવર્તી વિધેયોનો ગણ નીચેની બે ક્રિયાઓ હેઠળ બંધ છે:
\begin{enumerate}
\item સાન્ત સરવાળા: જો $f(\vec x, z)$ આદિમ પુનરાવર્તી હોય, તો
આ વિધેય પણ આદિમ પુનરાવર્તી છે:
\[
g(\vec x, y) \defis \sum_{z = 0}^y f(\vec x, z).
\]
\item સાન્ત ગુણાકાર: જો $f(\vec x, z)$ આદિમ પુનરાવર્તી હોય, તો
આ વિધેય પણ આદિમ પુનરાવર્તી છે:
\[
h(\vec x, y) \defis \prod_{z = 0}^y f(\vec x, z).
\]
\end{enumerate}
\end{prop}

\begin{proof}
દાખલા તરીકે, સાન્ત સરવાળાઓને નીચેના સમીકરણોથી પુનરાવર્તી રીતે
વ્યાખ્યાયિત કરીએ છીએ:
\begin{align*}
  g(\vec x, 0) & = f(\vec x, 0)\\
  g(\vec x, y+1) & = g(\vec x, y) + f(\vec x, y+1).
\end{align*}
\end{proof}
% END OLP-0216
\endgroup


\begingroup
\makeatletter\def\ol@sectioncs{section}\makeatother

% BEGIN OLP-0217
\olfileid{cmp}{rec}{prr}
\olsection{આદિમ પુનરાવર્તી સંબંધો}
\phantomsection\label{guunit:OLP-0217}



\begin{defn}
સંબંધ $R(\vec x)$નું લક્ષણ વિધેય
\[
\Char{R}(\vec x) = \left\{
  \begin{array}{ll}
  1 & \mbox{જો $R(\vec x)$ લાગુ પડે} \\
  0 & \mbox{અન્યથા}
  \end{array}
\right.
\]
આદિમ પુનરાવર્તી હોય, તો આ સંબંધને આદિમ પુનરાવર્તી
કહેવાય છે.
\end{defn}

બીજા શબ્દોમાં, આદિમ પુનરાવર્તી સંબંધ $R(\vec x)$ કહીએ ત્યારે
$\Char{R}(\vec x) = 1$ સ્વરૂપના સંબંધની વાત કરીએ છીએ. અહીં
$\Char{R}$ એવું આદિમ પુનરાવર્તી વિધેય છે જે કોઈ પણ ઇનપુટ પર
1 અથવા 0 આપે છે. દાખલા તરીકે, $x = 0$ હોય ત્યારે અને માત્ર
ત્યારે જ લાગુ પડતો સંબંધ $\fn{IsZero}(x)$, નીચેના આદિમ
પુનરાવર્તનથી વ્યાખ્યાયિત વિધેય $\Char{\fn{IsZero}}$ને અનુરૂપ છે:
\begin{align*}
\Char{\fn{IsZero}}(0) & = 1,\\
\Char{\fn{IsZero}}(x+1) & = 0.
\end{align*}

સંબંધોનું અન્ય આદિમ પુનરાવર્તી વિધેયો સાથે સંયોજન પણ કરી શકાય છે.
તેથી નીચેના સંબંધો પણ આદિમ પુનરાવર્તી છે:
\begin{enumerate}
\item સમાનતાનો સંબંધ $x = y$, જેને $\fn{IsZero}(\left|x -
  y\right|)$થી વ્યાખ્યાયિત કરીએ છીએ.
\item નાનું અથવા સમાન હોવાનો સંબંધ $x \leq y$, જેને
  $\fn{IsZero}(x \tsub y)$થી વ્યાખ્યાયિત કરીએ છીએ.
\end{enumerate}

\sourcecorrection{OLCMP-003}{સ્થિર અંગ્રેજી મૂળમાં બીજા મુદ્દાને
``less-than relation'' કહે છે, પરંતુ આપેલું સૂત્ર $x \leq y$ અને
$\fn{IsZero}(x \tsub y)$ બન્ને ``નાનું અથવા સમાન'' દર્શાવે છે.
ગુજરાતી નામમાં સૂત્ર પ્રમાણે સમાનતાનો સમાવેશ રાખ્યો છે.}

\begin{prop}
  આદિમ પુનરાવર્તી સંબંધોનો ગણ બુલિયન ક્રિયાઓ હેઠળ બંધ છે. એટલે
  જો $P(\vec x)$ અને $Q(\vec x)$ આદિમ પુનરાવર્તી હોય, તો નીચેના
  સંબંધો પણ એવા જ છે:
  \begin{enumerate}
  \item $\lnot P(\vec x)$
  \item $P(\vec x) \land Q(\vec x)$
  \item $P(\vec x) \lor Q(\vec x)$
  \item $P(\vec x) \lif Q(\vec x)$
  \end{enumerate}
\end{prop}

\begin{proof}
  ધારો કે $P(\vec x)$ અને $Q(\vec x)$ આદિમ પુનરાવર્તી છે, એટલે
  તેમનાં લક્ષણ વિધેયો $\Char{P}$ અને $\Char{Q}$ આદિમ પુનરાવર્તી છે.
  આપણે $\lnot P(\vec x)$ વગેરેના લક્ષણ વિધેયો પણ આદિમ પુનરાવર્તી છે
  તે બતાવવાનું છે.
  \[
  \Char{\lnot P}(\vec x) = \begin{cases}
    0 & \text{જો $\Char{P}(\vec x) = 1$ હોય}\\
    1 & \text{અન્યથા}
  \end{cases}
  \]
  $\Char{\lnot P}(\vec x)$ને $1 \tsub \Char{P}(\vec x)$ તરીકે
  વ્યાખ્યાયિત કરી શકીએ.
  \[
  \Char{P \land Q}(\vec x) = \begin{cases}
    1 & \text{જો $\Char{P}(\vec x) = \Char{Q}(\vec x) = 1$ હોય}\\
    0 & \text{અન્યથા}
  \end{cases}
  \]
  $\Char{P \land Q}(\vec x)$ને $\Char{P}(\vec x) \cdot
  \Char{Q}(\vec x)$ અથવા $\fn{min}(\Char{P}(\vec x), \Char{Q}(\vec
  x))$ તરીકે વ્યાખ્યાયિત કરી શકીએ. એ જ રીતે,
  \begin{align*}
    \Char{P \lor Q}(\vec x) & = \fn{max}(\Char{P}(\vec x), \Char{Q}(\vec x)) \text{ અને}\\
    \Char{P \lif Q}(\vec x) & = \fn{max}(1 \tsub
  \Char{P}(\vec x), \Char{Q}(\vec x)).
  \end{align*}
\end{proof}

\begin{prop}
  આદિમ પુનરાવર્તી સંબંધોનો ગણ સીમિત પરિમાણીકરણ હેઠળ બંધ છે. એટલે
  જો $R(\vec x, z)$ આદિમ પુનરાવર્તી સંબંધ હોય, તો નીચેના સંબંધો
  પણ એવા જ છે:
  \begin{align*}
    & \bforall{z < y}{R(\vec x, z)} \text{ અને}\\
    & \bexists{z < y}{R(\vec x, z)}.
  \end{align*}
  $\bforall{z < y}{R(\vec x, z)}$ એ $\vec x$ અને $y$ માટે ત્યારે
  અને માત્ર ત્યારે જ લાગુ પડે જ્યારે $y$થી નાની દરેક~$z$ માટે
  $R(\vec x, z)$ લાગુ પડે; $\bexists{z < y}{R(\vec x, z)}$ માટે પણ
  અનુરૂપ અર્થ છે.
\end{prop}

\begin{proof}
  રૂઢિ પ્રમાણે $\bforall{z < 0}{R(\vec x, z)}$ને સત્ય માનીએ છીએ
  (કારણ કે $0$થી નાની કોઈ $z$ છે જ નહીં) અને
  $\bexists{z < 0}{R(\vec x, z)}$ને અસત્ય માનીએ છીએ. સીમિત
  સર્વપરિમાણક સાન્ત ગુણાકાર અથવા ક્રમિક લઘુત્તમની જેમ કામ કરે છે.
  એટલે $P(\vec x, y) \defiff \bforall{z <
  y}{R(\vec x, z)}$ હોય તો $\Char{P}(\vec x, y)$ને આમ વ્યાખ્યાયિત
  કરી શકાય:
  \begin{align*}
    \Char{P}(\vec x, 0) & = 1\\
    \Char{P}(\vec x, y+1) & =
    \fn{min}(\Char{P}(\vec x, y), \Char{R}(\vec x, y)).
  \end{align*}
  સીમિત અસ્તિત્વપરિમાણીકરણને એ જ રીતે $\fn{max}$ વડે વ્યાખ્યાયિત
  કરી શકાય. અથવા તેને
  $\bexists{z < y}{R(\vec x, z)}
  \liff \lnot \bforall{z < y}{\lnot R(\vec x, z)}$ સમકક્ષતા
  વાપરીને સીમિત સર્વપરિમાણીકરણમાંથી મેળવી શકાય. નોંધો કે,
  દાખલા તરીકે, $\bexists{x \leq y}{\dots
  x\dots}$ સ્વરૂપનો સીમિત પરિમાણક
  $\bexists{x < y+1}{\dots x \dots}$ને સમકક્ષ છે.
\end{proof}

\begin{prob}
  ત્રણ આર્ગ્યુમેન્ટવાળો સંબંધ $x \equiv y \mod n$ (મોડ્યુલો~$n$
  મુજબ સર્વાંગસમતા) આદિમ પુનરાવર્તી છે એવું બતાવો.
\end{prob}

બીજું ઉપયોગી આદિમ પુનરાવર્તી વિધેય શરતી વિધેય
$\fn{cond}(x,y,z)$ છે, જે આમ વ્યાખ્યાયિત છે:
\begin{align*}
  \fn{cond}(x,y,z) & = \begin{cases}
  y & \text{જો $x = 0$ હોય} \\
  z & \text{અન્યથા}.
\end{cases}
\intertext{આ વિધેયને પુનરાવર્તી રીતે આમ વ્યાખ્યાયિત કરીએ:}
\fn{cond}(0,y,z) & = y,\\
\fn{cond}(x+1,y,z) & = z.
\end{align*}
આનો ઉપયોગ કરીને આદિમ પુનરાવર્તી સંબંધોના આધારે કિસ્સાવાર
વ્યાખ્યાયિત વિધેયોનું સમર્થન કરી શકીએ:

\begin{prop}
જો $g_0(\vec x)$, \dots,~$g_m(\vec x)$ આદિમ પુનરાવર્તી વિધેયો
હોય અને $R_0(\vec
x)$, \dots, $R_{m-1}(\vec x)$ આદિમ પુનરાવર્તી સંબંધો હોય, તો
નીચે મુજબ વ્યાખ્યાયિત વિધેય $f$ પણ આદિમ પુનરાવર્તી છે:
\[
f(\vec x) = \begin{cases}
    g_0(\vec x) & \text{જો $R_0(\vec{x})$ લાગુ પડે} \\
    g_1(\vec x) & \text{જો $R_1(\vec{x})$ લાગુ પડે અને $R_0(\vec{x})$ ન પડે} \\
    \vdots & \\
    g_{m-1}(\vec x) & \text{જો $R_{m-1}(\vec{x})$ લાગુ પડે અને અગાઉના
      કોઈ ન પડે}
    \\
    g_m(\vec x) & \mbox{અન્યથા}
\end{cases}
\]
\end{prop}

\begin{proof}
  $m = 1$ હોય ત્યારે આ માત્ર નીચે મુજબ વ્યાખ્યાયિત વિધેય છે:
  \[
  f(\vec x) = \fn{cond}(\Char{\lnot R_0}(\vec x),g_0(\vec x),g_1(\vec
  x)).
  \]
  $m$ એ $1$થી મોટું હોય ત્યારે આ સ્વરૂપની વ્યાખ્યાઓનું સંયોજન
  કરી શકાય છે.
\end{proof}
% END OLP-0217
\endgroup


\begingroup
\makeatletter\def\ol@sectioncs{section}\makeatother

% BEGIN OLP-0218
\olfileid{cmp}{rec}{bmi}
\olsection{સીમાબદ્ધ લઘુત્તમીકરણ}
\phantomsection\label{guunit:OLP-0218}


\begin{explain}
કોઈ ગુણધર્મ અથવા સંબંધ~$P$ સંતોષતી સૌથી નાની સંખ્યા તરીકે વિધેય
વ્યાખ્યાયિત કરવું ઘણી વાર ઉપયોગી છે. $P$નો નિર્ણય કરી શકાતો હોય,
તો $0$, $1$, $2$, \dots{} બધી શક્ય સંખ્યાઓ ક્રમે અજમાવતાં~$P$
સંતોષતી સૌથી નાની સંખ્યા મળે અને આ વિધેય ગણી શકાય. આવી અનિયત
મર્યાદા સુધીની શોધ આપણને આદિમ પુનરાવર્તી વિધેયોના ક્ષેત્રની બહાર લઈ
જાય છે. પરંતુ જો આપણે \emph{સ્વતંત્ર રીતે આપેલી કોઈ સીમાથી નાની}
સૌથી નાની સંખ્યા જ શોધવી હોય, તો આદિમ પુનરાવર્તી ક્ષેત્રમાં રહીએ
છીએ. બીજા શબ્દોમાં, અને થોડું વધુ સામાન્ય રીતે, ધારો કે
સંબંધ~$R(x,z)$ આદિમ પુનરાવર્તી છે. $x$ અને~$y$ને એવા સૌથી નાના
$z < y$ પર મોકલતું વિધેય વિચારો કે જેના માટે $R(x, z)$ લાગુ પડે.
$R(x, 0)$, $R(x, 1)$, \dots, $R(x, y-1)$ની તપાસ કરીને તેને પણ
ગણી શકાય. પણ તે આદિમ પુનરાવર્તી શા માટે છે?
\end{explain}

\begin{prop}
જો $R(\vec x, z)$ આદિમ પુનરાવર્તી હોય, તો વિધેય
$m_R(\vec{x}, y)$ પણ આદિમ પુનરાવર્તી છે. $R(\vec x, z)$ લાગુ પડે
એવી કોઈ સંખ્યા $y$થી નાની હોય, તો આ વિધેય સૌથી નાનો એવો~$z$ આપે
છે; અન્યથા $y$ આપે છે. વિધેય~$m_R$ને આ રીતે લખીશું:
\[
\bmin{z < y}{R(\vec{x}, z)},
\]
\end{prop}

\begin{proof}
$z < 0$ એવી સંખ્યા જ નથી; તેથી $R(\vec{x}, z)$ લાગુ પડે
એવી~$z < 0$ પણ હોઈ શકે નહીં. એટલે $m_R(\vec x, 0) = 0$.

સીમા $y + 1$ સ્વરૂપની હોય ત્યારે ત્રણ કિસ્સા છે:
\begin{enumerate}
\item $R(\vec{x}, z)$ લાગુ પડે એવી $z < y$ છે. ત્યારે
$m_R(\vec{x}, y+1) = m_R(\vec{x}, y)$.
\item એવી કોઈ~$z<y$ નથી, પરંતુ $R(\vec{x}, y)$ લાગુ પડે છે.
ત્યારે $m_R(\vec{x}, y+1) = y$.
\item $R(\vec{x}, z)$ લાગુ પડે એવી $z < y+1$ નથી. ત્યારે
$m_R(\vec{x}, y+1) = y+1$.
\end{enumerate}
\sourcecorrection{OLCMP-004}{સ્થિર અંગ્રેજી મૂળના ત્રીજા કિસ્સામાં
ડાબી બાજુ $m_R(\vec z, y+1)$ છપાયું છે. બાકીના બે કિસ્સા, પ્રસ્તાવનું
વિધેય અને તરત નીચેનું સમીકરણ બધે $\vec x$ છે; ગુજરાતી ત્રીજા
કિસ્સામાં પણ $m_R(\vec{x}, y+1)$ રાખ્યું છે.}
આથી $m_R(\vec x, 0)$ને આદિમ પુનરાવર્તન વડે નીચે મુજબ વ્યાખ્યાયિત
કરી શકીએ:
\begin{align*}
m_R(\vec{x}, 0) & = 0\\
m_R(\vec{x}, y+1) & =
\begin{cases}
m_R(\vec{x}, y) & \text{જો $m_R(\vec x, y) \neq y$ હોય}\\
y & \text{જો $m_R(\vec x, y) = y$ હોય અને $R(\vec{x}, y)$ લાગુ પડે}\\
y+1 & \text{અન્યથા.}
\end{cases}
\end{align*}
નોંધો કે $R(\vec x, z)$ લાગુ પડે એવી $z<y$ છે ત્યારે અને માત્ર
ત્યારે જ $m_R(\vec x,
y) \neq y$ છે.
\end{proof}

\begin{prob}
ધારો કે $R(\vec x, z)$ આદિમ પુનરાવર્તી છે. વિધેય
$m'_R(\vec{x}, y)$ને $\Char{R}$માંથી આદિમ પુનરાવર્તન વડે
વ્યાખ્યાયિત કરો. $R(\vec{x}, z)$ લાગુ પડે એવી કોઈ સંખ્યા $y$થી
નાનો હોય, તો તે સૌથી નાનો એવો~$z$ આપે અને અન્યથા $0$ આપે.
\end{prob}
% END OLP-0218
\endgroup


\begingroup
\makeatletter\def\ol@sectioncs{section}\makeatother

% BEGIN OLP-0219
\olfileid{cmp}{rec}{pri}
\olsection{અવિભાજ્ય સંખ્યાઓ}
\phantomsection\label{guunit:OLP-0219}


સીમિત પરિમાણીકરણ અને સીમાબદ્ધ લઘુત્તમીકરણથી ઘણા સ્વાભાવિક વિધેયો
અને સંબંધો આદિમ પુનરાવર્તી છે તે બતાવવા માટે પૂરતાં સાધનો મળે છે.
દાખલા તરીકે, ``$x$ એ $y$નો અવયવ છે'' એવો સંબંધ વિચારો, જેને
$x \mid y$ લખીએ છીએ. $y$ને~$x$થી ભાગતાં શેષ ન રહે, એટલે $y$ એ
~$x$નો પૂર્ણાંક ગુણિત હોય, ત્યારે $x \mid y$ લાગુ પડે છે. (ભાગાકારક
શૂન્ય ન હોય અને $y$ને $x$થી ભાગતાં શેષ $> 0$ હોય, તો $x \nmid y$
લખીએ છીએ. શૂન્ય ભાગાકારક માટે નીચેનું સમીકરણ સંબંધ નક્કી કરે છે.)
બીજા શબ્દોમાં, કોઈ~$z$ માટે $x \cdot z = y$ હોય
ત્યારે અને માત્ર ત્યારે જ $x \mid y$ છે. એવો $z$ અસ્તિત્વમાં
હોય, તો કોઈ એવો $z$ $\leq y$ પસંદ કરી શકાય છે. તેથી કોઈ
$z \le y$ માટે $x \cdot z = y$ હોય ત્યારે અને માત્ર ત્યારે જ
$x \mid y$ છે. સમાનતા~$=$ અને ગુણાકારમાંથી સીમિત
અસ્તિત્વપરિમાણીકરણ વડે $x \mid y$ સંબંધને આમ વ્યાખ્યાયિત કરી
શકીએ:
\[
x \mid y \defiff \bexists{z \leq y}{(x \cdot z) = y}.
\]
આથી $x \mid y$ આદિમ પુનરાવર્તી છે એવું આપણે બતાવ્યું છે.

\sourcecorrection{OLCMP-005}{સ્થિર અંગ્રેજી મૂળના કૌંસમાં
``$x$ને $y$થી ભાગતાં'' શેષ એવું લખાયું છે. $x \mid y$ના પહેલાંના
વાક્ય અને વ્યાખ્યાકારક સમીકરણ પ્રમાણે ભાગવાનું $y$ને $x$થી છે.
શૂન્ય ભાગાકારક માટે ભાગાકારનો શેષ વ્યાખ્યાયિત નથી; ગુજરાતી કૌંસમાં
આ અપવાદ પણ સ્પષ્ટ કર્યો છે.}
\sourcecorrection{OLCMP-006}{સ્થિર અંગ્રેજી મૂળ કહે છે કે
$x \cdot z = y$ સંતોષતો ``કોઈ પણ'' $z$ જરૂર $\leq y$ છે. $x=y=0$
હોય ત્યારે મોટા $z$ પણ સમીકરણ સંતોષે છે. સીમિત પરિમાણીકરણ માટે
માત્ર એવો ઓછામાં ઓછો એક $z \leq y$ મળે એટલું જરૂરી અને સાચું છે;
ગુજરાતી વાક્યમાં આ દાવો રાખ્યો છે.}

પ્રાકૃતિક સંખ્યા~$x$ \emph{અવિભાજ્ય} કહેવાય જો તે $0$ કે $1$
ન હોય અને માત્ર $1$ તથા પોતાનાથી જ વિભાજ્ય હોય. બીજા શબ્દોમાં,
અવિભાજ્ય સંખ્યાઓ માટે $y \mid x$ હોય ત્યારે $y = 1$ અથવા~$y=x$
હોય છે. $x$ અવિભાજ્ય છે કે નહીં તે તપાસવા $y \le x$ માટે જ
$y \mid x$ તપાસવું પૂરતું છે, કારણ કે $y > x$ હોય તો આપોઆપ
$y \nmid x$. તેથી $x$ અવિભાજ્ય હોય ત્યારે અને માત્ર ત્યારે જ
લાગુ પડતો સંબંધ $\fn{Prime}(x)$ આમ વ્યાખ્યાયિત કરી શકાય:
\[
\fn{Prime}(x) \defiff x \geq 2 \land \bforall{y \leq x}{(y \mid x \lif y
  = 1 \lor y = x)}
\]
અને આથી તે આદિમ પુનરાવર્તી છે.

અવિભાજ્ય સંખ્યાઓ $2$, $3$, $5$, $7$, $11$ વગેરે છે. આ શ્રેણીમાં
$x$મી અવિભાજ્ય સંખ્યા આપતું વિધેય $p(x)$ વિચારો: $p(0) =
2$, $p(1) = 3$, $p(2) = 5$ વગેરે. (સગવડ માટે આપણે ઘણી વાર
$p(x)$ને $p_x$ પણ લખીશું: $p_0=2$, $p_1=3$ વગેરે.)

$x$થી મોટી પહેલી અવિભાજ્ય સંખ્યા આપતું વિધેય
$\fn{nextPrime(x)}$ હોય, તો $p$ને આદિમ પુનરાવર્તન વડે સરળતાથી
વ્યાખ્યાયિત કરી શકાય:
\begin{align*}
  p(0) & = 2\\
  p(x+1) & = \fn{nextPrime}(p(x))
\end{align*}
$\fn{nextPrime}(x)$ એ $y > x$ અને $y$ અવિભાજ્ય હોય એવી સૌથી
નાની $y$ છે; અનિયત મર્યાદા સુધીની શોધ વડે તેને સરળતાથી ગણી
શકાય છે. પરંતુ યુક્લિડના પરિણામને કારણે તેને સીમાબદ્ધ
લઘુત્તમીકરણ વડે પણ વ્યાખ્યાયિત કરી શકાય: $x$ અને $\fact{x}+1$
વચ્ચે હંમેશાં એક અવિભાજ્ય સંખ્યા મળે છે.
\[
  \fn{nextPrime(x)} =
  \bmin{y \leq \fact{x}+1}{(y > x \land \fn{Prime}(y))}.
\]
આથી $\fn{nextPrime}(x)$ અને પરિણામે $p(x)$ માત્ર સંગણનીય જ નહીં,
પણ આદિમ પુનરાવર્તી છે.

(જિજ્ઞાસા હોય તો યુક્લિડના પરિણામની ટૂંકી સાબિતી આ છે. ધારો કે
$p_n$ એ $\le x$ એવી સૌથી મોટી અવિભાજ્ય સંખ્યા છે; અને
$\le x$ એવી બધી અવિભાજ્ય સંખ્યાઓનો ગુણાકાર $p =
p_0\cdot p_1 \cdot \dots \cdot p_n$ વિચારો. $p+1$ પોતે અવિભાજ્ય
હોય અથવા $x$ અને~$p+1$ વચ્ચે કોઈ અવિભાજ્ય સંખ્યા હોય. શા માટે?
ધારો કે $p+1$ અવિભાજ્ય નથી. ત્યારે કોઈ અવિભાજ્ય સંખ્યા
$q \mid p+1$ છે, જ્યાં $q < p+1$. $\le x$ એવી એકેય અવિભાજ્ય
સંખ્યા $p+1$ને ભાગતી નથી. ($p$ની વ્યાખ્યા પ્રમાણે દરેક
$p_i \le x$ એ~$p$ને ભાગે છે, એટલે શેષ~$0$ છે. તેથી દરેક
$p_i \le x$ વડે $p+1$ને ભાગતાં શેષ~$1$ રહે છે; એટલે
$p_i \nmid p+1$.) આથી $q$ એ $> x$ અને $< p+1$ એવી અવિભાજ્ય
સંખ્યા છે. તથા $p \le \fact{x}$ હોવાથી $> x$ અને
$\le \fact{x}+1$ એવી અવિભાજ્ય સંખ્યા મળે છે.)

\begin{prob}
સીમાબદ્ધ લઘુત્તમીકરણ વડે પૂર્ણાંક ભાગાકાર $d(x, y)$ વ્યાખ્યાયિત
કરો.
\end{prob}
% END OLP-0219
\endgroup


\begingroup
\makeatletter\def\ol@sectioncs{section}\makeatother

% BEGIN OLP-0220
\olfileid{cmp}{rec}{seq}
\olsection{અનુક્રમો}
\phantomsection\label{guunit:OLP-0220}


આદિમ પુનરાવર્તી વિધેયોનો ગણ ખાસ્સો સમર્થ છે. પરંતુ
\emph{અનુક્રમો}ને સંભાળવાની યોગ્ય રીત વિકસાવીએ પછી આપણે વધુ કરી
શકીશું. પ્રાકૃતિક સંખ્યાઓના સાન્ત અનુક્રમોને પ્રાકૃતિક સંખ્યાઓ સાથે
આ રીતે ઓળખીશું: અનુક્રમ $\langle a_0, a_1, a_2, \dots, a_k \rangle$
નીચેની સંખ્યાને અનુરૂપ છે:
\[
p_0^{a_0+1} \cdot p_1^{a_1+1} \cdot p_2^{a_2+1} \cdot \dots \cdot
p_k^{a_k+1}.
\]
ઘાતમાં એક ઉમેરીએ છીએ, જેથી દાખલા તરીકે
$\langle 2, 7, 3\rangle$ અને $\langle 2, 7, 3, 0, 0 \rangle$ના
સંખ્યાત્મક કોડ જુદા રહે. ખાલી અનુક્રમના કોડ તરીકે $0$ અને~$1$ બંને
લઈ શકાય; ચોક્કસતા માટે $\emptyseq$ એ~$0$ દર્શાવે એમ માનીએ.

અનુક્રમોનું આ સંકેતીકરણ સફળ થાય છે તેનું કારણ અંકગણિતનું કહેવાતું
મૂળભૂત પ્રમેય છે: દરેક પ્રાકૃતિક સંખ્યા $n \ge 2$ને નીચેના
સ્વરૂપમાં એક જ રીતે લખી શકાય:
\[
n = p_0^{a_0} \cdot p_1^{a_1} \cdot \dots \cdot p_k^{a_k}
\]
અહીં $a_k \ge 1$ છે. આથી નિશ્ચિત થાય છે કે વિધેય
$\tuple{}(a_0,
\dots, a_k) = \tuple{a_0, \dots, a_k}$ એક-એક છે: જુદા
અનુક્રમો જુદી સંખ્યાઓ પર જાય છે; દરેક સંખ્યાને વધુમાં વધુ એક
અનુક્રમ અનુરૂપ છે.

હવે બતાવીશું કે અનુક્રમની લંબાઈ શોધવી, તેનો $i$મો ઘટક શોધવો,
તેમાં ઘટક અંતે ઉમેરવો અને બે અનુક્રમો જોડવા---આ બધી ક્રિયાઓ આદિમ
પુનરાવર્તી છે.

\begin{prop}
  અનુક્રમ $s$ની લંબાઈ આપતું વિધેય $\len{s}$ આદિમ પુનરાવર્તી છે.
\end{prop}

\begin{proof}
  સંબંધ $R(i, s)$ને આ રીતે વ્યાખ્યાયિત કરીએ:
  \[
  R(i, s) \text{ ત્યારે અને માત્ર ત્યારે જ જ્યારે }
  p_i \mid s \land p_{i+1} \nmid s.
  \]
  સ્પષ્ટ છે કે $R$ આદિમ પુનરાવર્તી છે. જ્યારે પણ $s$ એ બિનખાલી
  અનુક્રમનો કોડ હોય, એટલે
  \[
  s = p_0^{a_0+1} \cdot \dots \cdot p_{k}^{a_{k}+1},
  \]
  ત્યારે $p_i$ એ એવી સૌથી મોટી અવિભાજ્ય સંખ્યા હોય કે
  $p_i \mid s$, તો $R(i,s)$ લાગુ પડે છે, એટલે $i = k$.
  તેથી $p_i$ એ~$s$ને ભાગતી
  સૌથી મોટી અવિભાજ્ય સંખ્યા હોય ત્યારે અને માત્ર ત્યારે જ $s$ની
  લંબાઈ $i+1$ છે. આથી આમ રાખી શકીએ:
  \[
  \len{s} =
  \begin{cases}
    0 & \text{જો $s = 0$ અથવા $s = 1$ હોય} \\
    1 + \bmin{i < s}{R(i, s)} & \text{અન્યથા}
  \end{cases}
  \]
  અહીં સીમાબદ્ધ લઘુત્તમીકરણ વાપરી શકીએ છીએ, કારણ કે $s$ એ અનુક્રમનો
  કોડ હોય ત્યારે $R(i,s)$ સંતોષતો એક જ $i$ હોય છે; અને એવો $i$ હોય તો
  પોતે~$s$થી નાનો છે.
\end{proof}

\begin{prop}
  અનુક્રમ $s$ના અંતે $a$ ઉમેરવાથી મળતું અનુક્રમ આપતું વિધેય
  $\fn{append}(s,a)$ આદિમ પુનરાવર્તી છે.
\end{prop}

\begin{proof}
  $\fn{append}$ને આમ વ્યાખ્યાયિત કરી શકીએ:
  \[
  \fn{append}(s,a) =
  \begin{cases}
    2^{a+1} & \text{જો $s = 0$ અથવા $s = 1$ હોય} \\
    s \cdot p_{\len{s}}^{a+1} & \text{અન્યથા.}
  \end{cases}
  \]
\end{proof}

\begin{prop}
  $s$નો $i$મો ઘટક આપતું વિધેય $\fn{element}(s,i)$ આદિમ
  પુનરાવર્તી છે; અહીં શરૂઆતના ઘટકને $0$મો કહીએ છીએ. જો $i$ એ
  $s$ની લંબાઈ જેટલો અથવા તેનાથી મોટો હોય, તો વિધેય $0$ આપે છે.
\end{prop}

\begin{proof}
$a$ એ~$s$નો $i$મો ઘટક હોય ત્યારે અને માત્ર ત્યારે જ
$p_i^{a+1}$ એ~$s$ને ભાગતો $p_i$નો સૌથી મોટો ઘાત છે, એટલે
$p_i^{a+1} \mid s$ પરંતુ $p_i^{a+2} \nmid s$. તેથી:
  \[
  \fn{element}(s,i) =
  \begin{cases}
    0 & \mbox{જો $i \geq \len{s}$ હોય} \\
    \bmin{a < s}{(p_i^{a+2} \nmid s)} & \text{અન્યથા.}
  \end{cases}
  \]
\end{proof}

ઉપર વ્યાખ્યાયિત વિધેયોનાં ઔપચારિક નામોને બદલે હવે વધુ ટૂંકાં
સંકેતનો લાવીએ. $\fn{element}(s,i)$ને બદલે $(s)_i$ લખીશું, અને
નીચેના સંક્ષેપ તરીકે $\tuple{s_0, \dots, s_k}$ લખીશું:
\[
\fn{append}(\fn{append}(\dots \fn{append}(\emptyseq,s_0)
\dots),s_k).
\]
નોંધો કે $s$ની લંબાઈ~$k$ હોય, તો $s$ના ઘટકો
$(s)_0$, \dots,~$(s)_{k-1}$ છે.

\begin{prop}
બે અનુક્રમો જોડતું વિધેય $\fn{concat}(s,t)$ આદિમ પુનરાવર્તી છે.
\end{prop}

\begin{proof}
  આપણે એવું વિધેય $\fn{concat}$ જોઈએ છે કે
  \[
    \fn{concat}(\tuple{a_0, \dots, a_k}, \tuple{b_0, \dots, b_l}) =
    \tuple{a_0, \dots, a_k, b_0, \dots, b_l}.
  \]
  સહાયક વિધેય $\fn{hconcat}(s,t,n)$ વાપરીશું, જે $t$ના પહેલાં
  $n$ ઘટકોને~$s$ના અંતે જોડે છે. તેને આદિમ પુનરાવર્તન વડે આમ
  વ્યાખ્યાયિત કરી શકાય:
  \begin{align*}
    \fn{hconcat}(s,t,0) & = s\\
    \fn{hconcat}(s,t,n+1) & = \fn{append}(\fn{hconcat}(s,t,n),(t)_n)
    \intertext{પછી $\fn{concat}$ને આમ વ્યાખ્યાયિત કરીએ:}
    \fn{concat}(s,t) & = \fn{hconcat}(s,t,\len{t}).
  \end{align*}
\end{proof}

$\fn{concat}(s,t)$ને બદલે આપણે $s \concat t$ લખીશું.

અનુક્રમની લંબાઈ અને તેના સૌથી મોટા ઘટકના આધારે તેનો સંખ્યાત્મક
કોડ ઉપરથી સીમિત કરી શકીએ, એ આપણને ઉપયોગી થશે. ધારો કે $s$ એ
લંબાઈ~$k$ ધરાવતો અનુક્રમ છે અને તેનો દરેક ઘટક કોઈ સંખ્યા~$x$થી
નાનો અથવા સમાન છે. તો~$s$ના વધુમાં વધુ $k$ અવિભાજ્ય અવયવો છે,
જે દરેક વધુમાં વધુ~$p_{k-1}$ છે; અને $s$ના અવિભાજ્ય અવયવીકરણમાં
દરેકનો ઘાત વધુમાં વધુ $x+1$ છે. બીજા શબ્દોમાં, જો આપણે
\[
\fn{sequenceBound}(x,k) = p_{k-1}^{k \cdot (x+1)},
\]
વ્યાખ્યાયિત કરીએ, તો ઉપર વર્ણવેલા અનુક્રમ~$s$નો સંખ્યાત્મક કોડ
વધુમાં વધુ~$\fn{sequenceBound}(x,k)$ છે.
\sourcecorrection{OLCMP-007}{સ્થિર અંગ્રેજી મૂળના આ સૂત્રમાં ખાલી
અનુક્રમ માટે $k=0$ લેતાં $p_{k-1}$ અનિર્ધારિત છે. ખાલી અનુક્રમનો કોડ
$0$ છે; આ સીમા-વિધેયને $k=0$ માટે મૂલ્ય $1$ આપીએ, જેથી ખાલી
અનુક્રમનો કોડ તેના કરતાં નાનો રહે. પ્રદર્શિત સૂત્ર $k\geq 1$ માટે છે.}

અનુક્રમો પર આવી સીમા હોવાથી સીમિત શોધ વડે નવાં વિધેયો
વ્યાખ્યાયિત કરી શકીએ. દાખલા તરીકે, $\fn{concat}$ને સીમિત શોધ
વડે વ્યાખ્યાયિત કરી શકીએ. આપણે શોધતા હોઈએ તે પદાર્થની (જોડાયેલા
અનુક્રમની સંખ્યાની) આદિમ પુનરાવર્તી \emph{વર્ણનાત્મક શરત} અને
ક્યાં સુધી શોધવું તેની સીમા લખવી પૂરતી છે. નીચેનું ચાલે છે:
\begin{align*}
  \fn{concat}(s,t) = {} & \bmin{v < \fn{sequenceBound}(s+t,\len{s} +
    \len{t})}{} \\
  & \quad(\len{v} = \len{s} + \len{t} \land {}\\
    & \qquad \bforall{i < \len{s}}{((v)_i = (s)_i) \land {} \\
      & \qquad \bforall{j < \len{t}}{((v)_{\len{s}+j} = (t)_j)})}
\end{align*}

\begin{prob}
$\fn{sconcat}(s)$ એવું આદિમ પુનરાવર્તી વિધેય છે કે જેનું
નીચેનું ગુણધર્મ છે, તે બતાવો:
\[
\fn{sconcat}(\tuple{s_0, \dots, s_k}) = s_0 \concat \dots \concat s_k.
\]
\end{prob}

\begin{prob}
$\fn{tail}(s)$ એવું આદિમ પુનરાવર્તી વિધેય છે કે જેનું
નીચેનું ગુણધર્મ છે, તે બતાવો:
\begin{align*}
  \fn{tail}(\emptyseq) & = 0 \text{ અને}\\
  \fn{tail}(\tuple{s_0, \dots, s_{k}}) & = \tuple{s_1, \dots, s_{k}}.
\end{align*}
\end{prob}

\begin{prop}
  \ollabel{prop:subseq}
  $s$ના $i$મા ઘટકથી શરૂ થતો લંબાઈ~$n$નો ઉપઅનુક્રમ આપતું વિધેય
  $\fn{subseq}(s, i, n)$ આદિમ પુનરાવર્તી છે.
\end{prop}

\begin{proof}
  અભ્યાસ માટે.
\end{proof}

\begin{prob}
  \olref[cmp][rec][seq]{prop:subseq} સાબિત કરો.
\end{prob}
% END OLP-0220
\endgroup


\begingroup
\makeatletter\def\ol@sectioncs{section}\makeatother

% BEGIN OLP-0221
\olfileid{cmp}{rec}{tre}
\olsection{વૃક્ષો}
\phantomsection\label{guunit:OLP-0221}


ક્યારેક વૃક્ષોને પ્રાકૃતિક સંખ્યાઓ તરીકે રજૂ કરવું ઉપયોગી છે. જેમ
અનુક્રમોને સંખ્યાઓથી રજૂ કરીએ છીએ અને તેમનાં ગુણધર્મો તથા ક્રિયાઓને
તેમના કોડ પરના આદિમ પુનરાવર્તી સંબંધો અને વિધેયોથી રજૂ કરીએ છીએ,
તેમ વૃક્ષો માટે પણ કરી શકીએ. આ માટે અનુક્રમો અને તેમના કોડ
વાપરીશું. વૃક્ષમાં માત્ર એક જ ગાંઠ હોઈ શકે છે (જેને નામચિહ્ન હોય
કે ન હોય), અથવા કેટલીક ઉપવૃક્ષો સાથે જોડાયેલી ગાંઠ હોઈ શકે છે
(જેને નામચિહ્ન હોય કે ન હોય). એ ગાંઠને વૃક્ષનું \emph{મૂળ} અને
તેની સાથે જોડાયેલાં ઉપવૃક્ષોને તેનાં \emph{તત્કાલ ઉપવૃક્ષો} કહીએ.

વૃક્ષને પુનરાવર્તી રીતે અનુક્રમ $\tuple{k, d_1, \dots, d_k}$ વડે
સંકેતીકૃત કરીએ છીએ. અહીં $k$ એ તત્કાલ ઉપવૃક્ષોની સંખ્યા છે અને
$d_1$, \dots,~$d_k$ તેમના કોડ છે. ગાંઠોને નામચિહ્નો હોય, તો
તેમને તત્કાલ ઉપવૃક્ષોની પાછળ ઉમેરી શકાય. તેથી માત્ર એક ગાંઠનું,
નામચિહ્ન~$l$ ધરાવતું વૃક્ષ $\tuple{0,l}$થી સંકેતીકૃત થશે.
નામચિહ્ન~$l_1$ ધરાવતા મૂળને નામચિહ્નો $l_2$, $l_3$ ધરાવતી બે
એકગાંઠવાળી ઉપવૃક્ષો જોડાયેલી હોય, તો તેનો કોડ
$\tuple{2,
  \tuple{0, l_2}, \tuple{0, l_3}, l_1}$ થશે.

\begin{prop}
  \ollabel{prop:subtreeseq}
  વિધેય $\fn{SubtreeSeq}(t)$ આદિમ પુનરાવર્તી છે. તે એવા
  અનુક્રમનો કોડ આપે છે જેના ઘટકો કોડ~$t$ ધરાવતા વૃક્ષનાં બધાં
  ઉપવૃક્ષોના કોડ છે.
\end{prop}

\begin{proof}
  પહેલાં નોંધો કે $\fn{ISubtrees}(t) = \fn{subseq}(t, 1, (t)_0)$
  આદિમ પુનરાવર્તી છે અને વૃક્ષ~$t$નાં તત્કાલ ઉપવૃક્ષોના કોડ આપે
  છે. હવે સહાયક વિધેય $\fn{hSubtreeSeq}(t,n)$ વ્યાખ્યાયિત કરી
  શકીએ, જે મૂળથી વધુમાં વધુ $n$ કડીઓ દૂરનાં બધાં ઉપવૃક્ષોનો
  અનુક્રમ ગણે છે. મૂળથી $0$ કડી દૂરનાં~$t$નાં ઉપવૃક્ષોનો
  અનુક્રમ---અર્થાત્~$t$ના મૂળથી શરૂ થતો અનુક્રમ---માત્ર~$t$ ધરાવે છે.
  વૃક્ષ~$t$ના મૂળથી વધુમાં વધુ $n+1$ કડીઓ દૂરનાં બધાં ઉપવૃક્ષોનો
  અનુક્રમ મેળવવા, વધુમાં વધુ $n$ કડીઓ દૂરનાં ઉપવૃક્ષોના અનુક્રમ
  સાથે આ $n$ સ્તર સુધીનાં ઉપવૃક્ષોનાં બધાં તત્કાલ ઉપવૃક્ષોનો
  અનુક્રમ જોડીએ. આ બધાની યાદી બનાવવા નોંધો કે
  $f(x)$ એ અનુક્રમોના કોડ આપતું આદિમ પુનરાવર્તી વિધેય હોય, તો
  પહેલાં $k$ ઘટકો પરનું જોડાણ $g_f(s, k) = f((s)_0) \concat
  \dots \concat f((s)_{k-1})$ પણ આદિમ પુનરાવર્તી છે; શૂન્ય
  ઘટકો માટે તે ખાલી અનુક્રમ છે:
    \begin{align*}
      g_f(s, 0) & = \emptyseq\\
      g_f(s, k+1) & = g_f(s, k) \concat f((s)_k)
    \end{align*}
    \sourcecorrection{OLCMP-008}{સ્થિર અંગ્રેજી મૂળ
    $g_f(s,k)$માં $0$થી $k$ સુધી, એટલે $k+1$ ઘટકો જોડે છે અને
    સમીકરણોમાં $g$ લખે છે; પછી $h(s)$ માટે $\len{s}$ વાપરે છે.
    ગુજરાતી વ્યાખ્યામાં પહેલાં $k$ ઘટકો જોડીએ છીએ, $k=0$ માટે
    ખાલી અનુક્રમ રાખીએ છીએ, અને સમીકરણોમાં $g_f$ રાખીએ છીએ.
    તેથી પછીનો $g_{\fn{ISubtrees}}(s,\len{s})$ ખરેખર $s$ના
    બધા ઘટકો પર જ ચાલે છે.}
    દાખલા તરીકે, $s$ એ વૃક્ષોનો અનુક્રમ હોય, તો
    $h(s) = g_{\fn{ISubtrees}}(s, \len{s})$ એ~$s$ના ઘટકોનાં
    તત્કાલ ઉપવૃક્ષોનો અનુક્રમ આપે છે. આ વડે $\fn{hSubtreeSeq}$ને
    આમ વ્યાખ્યાયિત કરી શકીએ:
    \begin{align*}
      \fn{hSubtreeSeq}(t, 0) & = \tuple{t} \\
      \fn{hSubtreeSeq}(t, n+1) & = \fn{hSubtreeSeq}(t, n) \concat
      h(\fn{hSubtreeSeq}(t, n)).
    \end{align*}
    \sourcecorrection{OLCMP-009}{સ્થિર અંગ્રેજી મૂળ આ સહાયક
    વિધેયને મૂળથી બરાબર $n$ કડીઓ દૂરનાં ઉપવૃક્ષોનો અનુક્રમ કહે છે.
    પ્રદર્શિત પુનરાવર્તી સમીકરણ અગાઉની સમગ્ર યાદી સાથે નવા
    ઉપવૃક્ષો જોડે છે; તેથી તેમાં વધુમાં વધુ $n$ કડીઓ દૂરનાં
    ઉપવૃક્ષો રહે છે. ગુજરાતી વર્ણનમાં આ સંચિત અર્થ રાખ્યો છે.}
    કોડ~$t$ ધરાવતા વૃક્ષમાં ઉપવૃક્ષોનું મહત્તમ સ્તર, એટલે મૂળથી
    છેડાની ગાંઠ સુધીનું મહત્તમ અંતર, કોડ~$t$થી ઉપરથી સીમિત છે.
    તેથી કોડ~$t$ ધરાવતા વૃક્ષનાં બધાં ઉપવૃક્ષોના કોડનો અનુક્રમ
    $\fn{hSubtreeSeq}(t, t)$ આપે છે.
\end{proof}

\begin{prob}
  \olref[cmp][rec][tre]{prop:subtreeseq}ની સાબિતીમાં
  $\fn{hSubtreeSeq}$ની વ્યાખ્યા સામાન્ય રીતે પુનરાવૃત્તિઓ
  ધરાવે છે. મળતી યાદીમાં કોઈ ઉપવૃક્ષનો કોડ માત્ર એક વાર જ આવે,
  તેની ખાતરી આપતી બીજી વ્યાખ્યા આપો.
\end{prob}
% END OLP-0221
\endgroup


\begingroup
\makeatletter\def\ol@sectioncs{section}\makeatother

% BEGIN OLP-0222
\olfileid{cmp}{rec}{ore}
\olsection{પુનરાવર્તનનાં અન્ય સ્વરૂપો}
\phantomsection\label{guunit:OLP-0222}


જોડ-નિરૂપણ અને અનુક્રમ-નિરૂપણ વડે આદિમ પુનરાવર્તનનાં વધુ અસામાન્ય
(અને ઉપયોગી) સ્વરૂપોનું સમર્થન કરી શકીએ છીએ. દાખલા તરીકે, નીચેની
વ્યાખ્યામાં જેમ થાય છે તેમ, બે વિધેયો એકસાથે વ્યાખ્યાયિત કરવાં ઘણી
વાર ઉપયોગી છે:
\begin{align*}
h_0(\vec x, 0) & = f_0(\vec x) \\
h_1(\vec x, 0) & = f_1(\vec x) \\
h_0(\vec x, y+1) & = g_0(\vec x, y, h_0(\vec x, y), h_1(\vec x, y)) \\
h_1(\vec x, y+1) & = g_1(\vec x, y, h_0(\vec x, y), h_1(\vec x, y))
\end{align*}
આ \emph{એકસાથે પુનરાવર્તન}નું ઉદાહરણ છે. વિધેયો વ્યાખ્યાયિત
કરવાની બીજી ઉપયોગી રીતમાં, નીચે મુજબ $h(\vec x, y+1)$નું મૂલ્ય
અગાઉનાં \emph{બધાં} મૂલ્યો $h(\vec x, 0)$, \dots,~$h(\vec x, y)$ના
આધારે આપીએ છીએ:
\begin{align*}
h(\vec x, 0) & = f(\vec x) \\
h(\vec x, y+1) & = g(\vec x, y, \tuple{h(\vec x, 0), \dots, h(\vec x, y)}).
\end{align*}
આ વિચારને નીચેની યોજના વધુ ટૂંકમાં વ્યક્ત કરે છે:
\[
h(\vec x, y) = g(\vec x, y, \tuple{h(\vec x, 0), \dots, h(\vec x, y-1)})
\]
અહીં $y$ એ $0$ હોય ત્યારે $g$નું છેલ્લું આર્ગ્યુમેન્ટ માત્ર ખાલી
અનુક્રમ છે એમ સમજવું. બંને રજૂઆતોમાં વિચાર એ છે કે ``ઉત્તરગામી
તબક્કો'' ગણતી વખતે $h$ અત્યાર સુધી ગણાયેલા મૂલ્યોનો આખો અનુક્રમ
વાપરી શકે છે. તેને \emph{અગાઉનાં બધાં મૂલ્યો પર આધારિત પુનરાવર્તન}
કહે છે. ખાસ ઉદાહરણ તરીકે, તે નીચેના પ્રકારની વ્યાખ્યાનું સમર્થન કરે
છે:
\begin{align*}
h(\vec x, y) & = \begin{cases}
  g(\vec x, y, h(\vec x, k(\vec x, y))) & \text{જો $k(\vec x, y) < y$ હોય} \\
  f(\vec x) & \text{અન્યથા}
\end{cases}
\end{align*}
બીજા શબ્દોમાં, $y$ પર $h$નું મૂલ્ય~$k$ વડે અપાયેલા પહેલાંના
\emph{કોઈ પણ} સ્થાને $h$ના મૂલ્ય પર આધારિત ગણી શકાય છે.


\begin{prob}
  અગાઉનાં બધાં મૂલ્યો પર આધારિત પુનરાવર્તન વડે શેષ વિધેય
  $r(x,y)$ વ્યાખ્યાયિત કરો. ($x$, $y$ પ્રાકૃતિક સંખ્યાઓ અને
  $y > 0$ હોય, તો $r(x,y)$ એ~$y$થી નાની એવી સંખ્યા છે કે કોઈ~$z$
  માટે $x = z\times y + r(x,y)$ થાય છે. ચોક્કસતા માટે $y=0$ હોય
  ત્યારે $r(x,0) = 0$ માનીએ.)
\end{prob}

આ વિધેયો સામાન્ય આદિમ પુનરાવર્તનથી કેવી રીતે મેળવી શકાય તે
વિચારો. આદિમ પુનરાવર્તનનું એક છેલ્લું સ્વરૂપ વધુ લવચીક છે: તેમાં
ચાલતાં ચાલતાં \emph{પ્રાચલો} (બાજુનાં મૂલ્યો) બદલવાની છૂટ છે:
\begin{align*}
h(\vec x, 0) & = f(\vec x) \\
h(\vec x, y+1) & = g(\vec x, y, h(k(\vec x), y))
\end{align*}
તેને પણ સામાન્ય આદિમ પુનરાવર્તનથી અનુકરિત કરી શકાય છે. (આ કરવું
સરળ નથી. સંકેત તરીકે ગણતરીને હાથથી ઊલટી દિશામાં ઉકેલી જુઓ.)
% END OLP-0222
\endgroup


\begingroup
\makeatletter\def\ol@sectioncs{section}\makeatother

% BEGIN OLP-0223
\olfileid{cmp}{rec}{npr}
\olsection{આદિમ પુનરાવર્તી ન હોય તેવાં વિધેયો}
\phantomsection\label{guunit:OLP-0223}


આદિમ પુનરાવર્તી વિધેયોમાં સહજ રીતે સંગણનીય લાગતાં બધાં વિધેયો
સમાઈ જતાં નથી. સહજ રીતે સ્પષ્ટ હોવું જોઈએ કે એક આર્ગ્યુમેન્ટવાળાં
બધાં આદિમ પુનરાવર્તી વિધેયોની યાદી
$f_0$, $f_1$, $f_2$,~\dots{} બનાવી શકાય છે, એવી રીતે કે ઇનપુટ
$y$ પર $f_x$નું મૂલ્ય અસરકારક રીતે ગણી શકાય; બીજા શબ્દોમાં, નીચે
મુજબ વ્યાખ્યાયિત વિધેય $g(x,y)$ સંગણનીય છે:
\[
g(x,y) = f_x(y)
\]
પરંતુ ત્યાર પછી નીચેનું વિધેય પણ સંગણનીય છે:
\begin{eqnarray*}
h(x) & = & g(x,x) + 1 \\
& = & f_x(x) +1.
\end{eqnarray*}
દરેક આદિમ પુનરાવર્તી વિધેય $f_i$ માટે $i$ પર $h$ અને $f_i$નાં
મૂલ્યો જુદાં પડે છે. તેથી $h$ સંગણનીય છે, પરંતુ આદિમ પુનરાવર્તી
નથી; $g$ વિશે પણ એ જ કહી શકાય. આ કૅન્ટૉરની વિકર્ણ-રચનાની
``અસરકારક'' આવૃત્તિ છે.

સંગણનીય છતાં આદિમ પુનરાવર્તી ન હોય તેવાં વિધેયોનાં વધુ સ્પષ્ટ
ઉદાહરણો આપી શકાય. દાખલા તરીકે, $g^n(x)$ સંકેત $n$ વખત $g$
વપરાય એવા $g(g(\dots g(x)))$ને દર્શાવે; અને વિધેયોનો અનુક્રમ
$g_0,g_1,\dots$ આમ વ્યાખ્યાયિત કરીએ:
\begin{eqnarray*}
g_0(x) & = & x+1 \\
g_{n + 1}(x) & = & g_n^x(x)
\end{eqnarray*}
દરેક વિધેય $g_n$ આદિમ પુનરાવર્તી છે, તે ચકાસી શકાય. દરેક
આગળનું વિધેય તેના પહેલાંના કરતાં ઘણું ઝડપથી વધે છે: $g_1(x)$ એ
$2x$ છે, $g_2(x)$ એ $2^x \cdot x$ છે અને $g_3(x)$ આશરે $x$
સંખ્યામાં $2$ની ઘાત-મિનારની જેમ વધે છે. અકરમાન--પેટર વિધેય
મૂળભૂત રીતે $G(x) = g_x(x)$ વિધેય છે, અને તે કોઈ પણ આદિમ
પુનરાવર્તી વિધેય કરતાં વધુ ઝડપથી વધે છે એવું બતાવી શકાય છે.

હવે આદિમ પુનરાવર્તી વિધેયોની પરિગણનાના પ્રશ્ન તરફ પાછા ફરીએ.
યાદ કરો કે દરેક આદિમ પુનરાવર્તી વિધેયને આપણે પ્રતીકાત્મક
સંકેતન આપ્યું છે; તેથી સંકેતનોની પરિગણના કરવી પૂરતી છે. દરેક
સંકેતન $F$ને પ્રાકૃતિક સંખ્યા $\#(F)$ પુનરાવર્તી રીતે આ મુજબ
આપી શકીએ:
\begin{eqnarray*}
\#(0) & = & \langle 0 \rangle \\
\#(S) & = & \langle 1 \rangle \\
\#(\Proj{n}{i}) & = & \langle 2, n, i \rangle \\
\#(\fn{Comp}_{k,l}[H,G_0,\dots,G_{k-1}]) & = & \langle
3,k,l,\#(H),\#(G_0),\dots,\#(G_{k-1}) \rangle \\
\#(\fn{Rec}_l[G,H]) & = & \langle 4, l, \#(G), \#(H) \rangle
\end{eqnarray*}
અહીં પાછલા વિભાગના કોડ વાપરીને સંખ્યાઓના દરેક અનુક્રમને પ્રાકૃતિક
સંખ્યા તરીકે જોવાની હકીકત કામમાં લઈએ છીએ. પરિણામે દરેક કોડને
પ્રાકૃતિક સંખ્યા મળે છે. અલબત્ત, કેટલાક અનુક્રમો (અને તેથી કેટલીક
સંખ્યાઓ) કોઈ સંકેતનને અનુરૂપ નથી. પરંતુ $i$ કોઈ સંકેતનનો કોડ
હોય, તો $f_i$ એ $i$ વડે કોડ થયેલું સંકેતન ધરાવતું એક-આર્ગ્યુમેન્ટવાળું આદિમ
પુનરાવર્તી વિધેય માનીએ; અન્યથા તે હંમેશાં $0$ આપતું વિધેય માનીએ.
આ રીતે એક-આર્ગ્યુમેન્ટવાળાં આદિમ પુનરાવર્તી વિધેયોની સ્પષ્ટ
પરિગણના મળે છે.

(હકીકતમાં, અચળ શૂન્ય વિધેય જેવાં કેટલાંક વિધેયો યાદીમાં એકથી વધુ
વખત આવશે. આ માત્ર આપણા સંકેતીકરણની આડઅસર નથી; અચળ શૂન્ય
વિધેયને એક કરતાં વધુ સંકેતનો મળે છે તેનું પણ આ પરિણામ છે.
પછી આપણે જોઈશું કે આવી પુનરાવૃત્તિઓ સંગણનીય રીતે દૂર કરી શકાતી
નથી. દાખલા તરીકે, આપેલું સંકેતન અચળ શૂન્ય વિધેય દર્શાવે છે કે
નહીં તે નક્કી કરતું કોઈ સંગણનીય વિધેય નથી.)

હવે આપણે $g(x,y)$ને $f_x(y)$ વડે આપેલું માનીએ, જ્યાં $f_x$ હમણાં
વર્ણવેલી પરિગણનાનું વિધેય છે. $g(x,y)$ સંગણનીય છે તે કેવી રીતે
જાણીએ? સહજ રીતે આ સ્પષ્ટ છે: $g(x,y)$ ગણવા પહેલાં $x$નો કોડ
``ઉકેલો'' અને જુઓ કે તે એક-આર્ગ્યુમેન્ટવાળા વિધેયનું સંકેતન છે
કે નહીં. હોય તો તે વિધેયનું ઇનપુટ~$y$ પર મૂલ્ય ગણો.

\begin{digress}
તમને કદાચ ખાતરી થઈ ગઈ હશે કે (થોડી મહેનત કરીને!) આ કામ કરતો
પ્રોગ્રામ (કહો કે Java કે C++માં) લખી શકાય. પછી ચર્ચ--ટ્યુરિંગ
માન્યતાનો આધાર લઈ શકીએ: સહજ અર્થમાં જે કંઈ સંગણનીય હોય, તેને
ટ્યુરિંગ મશીનથી ગણી શકાય.

અલબત્ત $g(x,y)$ સંગણનીય છે તે બતાવવાની વધુ સીધી રીત, તેને ગણતું
ટ્યુરિંગ મશીન સ્પષ્ટ રીતે વર્ણવવાની છે. તેમાં ખાસ કરીને
ચર્ચ--ટ્યુરિંગ માન્યતા અને સહજ સમજનો આધાર લેવાની જરૂર નહીં રહે.
ટૂંક સમયમાં આપણે એટલાં સાધનો વિકસાવીશું કે ટ્યુરિંગ મશીન પર
\emph{અનુકરિત} થઈ શકે તેવા સંગણનના એક નમૂનાનો આધાર લઈને
$g(x,y)$ સંગણનીય છે એવું બતાવી શકીશું: એ નમૂનો પુનરાવર્તી
વિધેયોનો છે.
\end{digress}
% END OLP-0223
\endgroup


\begingroup
\makeatletter\def\ol@sectioncs{section}\makeatother

% BEGIN OLP-0224
\olfileid{cmp}{rec}{par}
\olsection{આંશિક પુનરાવર્તી વિધેયો}
\phantomsection\label{guunit:OLP-0224}



પુનરાવર્તી વિધેયોની વ્યાખ્યા તરફ વધતાં પહેલાં નોંધો કે આદિમ
પુનરાવર્તી ન હોય એવાં સંગણનીય વિધેયો અસ્તિત્વમાં છે એવી આપણી
સાબિતી ખરેખર ઘણું વધુ બતાવે છે. દલીલ સરળ હતી: આપણે માત્ર એ
હકીકત વાપરી કે વિધેયો $f_0,f_1,\dots$ને એવી રીતે પરિગણિત કરી
શકાય કે $x$ અને $y$ના વિધેય તરીકે $f_x(y)$ સંગણનીય હોય. તેથી
આ દલીલ \emph{આ રીતે પરિગણિત કરી શકાય એવા વિધેયોના કોઈ પણ વર્ગ}
માટે લાગુ પડે છે. હવે મુશ્કેલી એ છે કે આપણે સંગણનીય વિધેયોનું
સ્પષ્ટ વર્ણન કરવા માગીએ છીએ, પણ સંગણનીય વિધેયોના સંગ્રહનું કોઈ
પણ આવું સ્પષ્ટ વર્ણન સંપૂર્ણ હોઈ શકે નહીં!

ઉકેલ એ છે કે \emph{આંશિક} વિધેયોને પણ સામેલ કરીએ. આપણે જોઈશું કે
આંશિક સંગણનીય વિધેયોની પરિગણના \emph{શક્ય} છે.
 ખરેખર, ટ્યુરિંગ મશીનોને વ્યવસ્થિત રીતે પરિગણિત
  કરી શકાય છે, તેથી આ વાત લગભગ આપણે પહેલેથી જાણીએ છીએ.
પછી આપણે વિકર્ણ-દલીલ તરફ પાછા ફરીશું અને આંશિક વિધેયોનો
સમાવેશ કરતાં એ દલીલ કેમ ચાલતી નથી તે તપાસીશું.

હવે પ્રશ્ન એ છે કે બધા આંશિક પુનરાવર્તી વિધેયો મેળવવા માટે
આદિમ પુનરાવર્તી વિધેયોમાં શું ઉમેરવું જોઈએ? બે બાબતો કરવી પડશે:
\begin{enumerate}
\item આદિમ પુનરાવર્તી વિધેયોની વ્યાખ્યા બદલીને તેમાં આંશિક
  વિધેયોને પણ આવવા દેવા.
\item વ્યાખ્યામાં કંઈક \emph{ઉમેરવું}, જેથી નવા આંશિક વિધેયોનો
  સમાવેશ થાય.
\end{enumerate}

પહેલું સહેલું છે. પહેલાંની જેમ શૂન્ય વિધેય, ઉત્તરગામી વિધેય અને
પ્રક્ષેપોથી શરૂ કરીને સંયોજન અને આદિમ પુનરાવર્તન હેઠળ બંધ કરીશું.
ફરક માત્ર એટલો કે વ્યાખ્યામાંના કેટલાક પદોનું મૂલ્ય ન હોય એવી
શક્યતા સમાવવા સંયોજન અને આદિમ પુનરાવર્તનની વ્યાખ્યાઓ બદલવી પડશે.
$f$ અને $g$ આંશિક વિધેયો હોય, તો $f(x) \fdefined$થી જણાવીશું કે
$f$ એ $x$ પર વ્યાખ્યાયિત છે, એટલે $x$ એ $f$ના પ્રદેશમાં છે;
અને $f(x)
\fundefined$થી તેનો વિરોધી અર્થ, એટલે $f$ એ~$x$ પર
વ્યાખ્યાયિત નથી, જણાવીશું. $f(x) \simeq g(x)$થી જણાવીશું કે
$f(x)$ અને $g(x)$ બંને અવ્યાખ્યાયિત છે, અથવા બંને વ્યાખ્યાયિત
અને સમાન છે. વધુ જટિલ પદો માટે પણ આ સંકેતનો વાપરીશું. એક
રૂઢિ અપનાવીએ: $h$ અને $g_0$, \dots,~$g_k$ બધાં આંશિક
વિધેયો હોય, તો
\[
h(g_0(\vec x),\dots,g_k(\vec x))
\]
ત્યારે અને માત્ર ત્યારે જ વ્યાખ્યાયિત છે જ્યારે દરેક $g_i$ એ
$\vec x$ પર વ્યાખ્યાયિત હોય અને $h$ એ
$g_0(\vec x)$, \dots,~$g_k(\vec x)$ પર વ્યાખ્યાયિત હોય. આ
સમજ સાથે આંશિક વિધેયો માટે સંયોજન અને આદિમ પુનરાવર્તનની
વ્યાખ્યાઓ અગાઉ જેવી જ છે; માત્ર ``$=$''ની જગ્યાએ
``$\simeq$'' મૂકવું પડે છે.

આદિમ પુનરાવર્તી વિધેયોની વ્યાખ્યામાં આંશિક વિધેયો મેળવવા જે
ઉમેરીશું તે \emph{અનિયત મર્યાદા સુધીની શોધક્રિયા} છે. પ્રાકૃતિક
સંખ્યાઓ પર $f(x,\vec z)$ કોઈ પણ આંશિક વિધેય હોય, તો
$\mu x \; f(x,\vec z)$ને આ રીતે વ્યાખ્યાયિત કરીએ:
\begin{quote}
  જો એવો $x$ અસ્તિત્વમાં હોય કે $f(0,\vec z), f(1,\vec z),
  \dots, f(x,\vec
  z)$ બધાં વ્યાખ્યાયિત હોય અને $f(x,\vec z) = 0$ હોય, તો એવા
  $x$માંથી સૌથી નાનો.
\end{quote}
અન્યથા $\mu x \; f(x,\vec z)$ અવ્યાખ્યાયિત છે એમ સમજવું.
આથી $\mu x \; f(x,\vec z)$ એકમાત્ર રીતે વ્યાખ્યાયિત થાય છે.

\begin{explain}
નોંધો કે આપણી વ્યાખ્યા ટ્યુરિંગ મશીનો, અથવા
અલ્ગોરિધમો કે સંગણનના કોઈ ખાસ નમૂનાનો ઉલ્લેખ કરતી નથી. પરંતુ
સંયોજન અને આદિમ પુનરાવર્તનની જેમ અનિયત મર્યાદા સુધીની શોધ પાછળ
ક્રિયાત્મક, સંગણનાત્મક સમજ છે. આંશિક વિધેયની સંગણનીયતામાં જે
આર્ગ્યુમેન્ટો પર વિધેય અવ્યાખ્યાયિત છે તે એવા ઇનપુટને અનુરૂપ છે
જેના પર ગણતરી અટકતી નથી. $\mu x \;
f(x,\vec z)$ ગણવાની રીત આ છે: $f(0,\vec z), f(1,\vec z),
f(2,\vec z)$ વગેરેની ગણતરી 0 મૂલ્ય મળે ત્યાં સુધી કરો.
વચલી કોઈ ગણતરી ન અટકે, તો $\mu x \; f(x,\vec z)$ની ગણતરી પણ
નહીં અટકે.
\end{explain}

$R(x,\vec z)$ કોઈ પણ સંબંધ હોય, તો $\mu x \; R(x,\vec z)$ને
$\mu x \; (1 \tsub \Char{R}(x,\vec z))$ તરીકે વ્યાખ્યાયિત
કરીએ છીએ. બીજા શબ્દોમાં, $\mu x \; R(x,\vec z)$ એ $R(x,\vec z)$
લાગુ પડે એવું $x$નું સૌથી નાનું મૂલ્ય આપે છે. તેથી $f(x,\vec z)$
સંપૂર્ણ વિધેય હોય, તો $\mu x \; f(x,\vec
z)$ એ $\mu x \; (f(x,\vec z) = 0)$ જેવું જ છે. જોકે આપણી
મૂળ વ્યાખ્યા વધુ સામાન્ય છે: તેમાં $f(x,\vec z)$ દરેક જગ્યાએ
વ્યાખ્યાયિત ન હોય એવી છૂટ છે; તેની સામે સંબંધનું લક્ષણ વિધેય
હંમેશાં સંપૂર્ણ હોય છે.

\begin{defn}
\emph{આંશિક પુનરાવર્તી વિધેયો}નો ગણ એટલે પ્રાકૃતિક સંખ્યાઓમાંથી
પ્રાકૃતિક સંખ્યાઓમાં જતાં (જુદી જુદી આર્ગ્યુમેન્ટ-સંખ્યા ધરાવતાં)
આંશિક વિધેયોનો સૌથી નાનો એવો ગણ, જેમાં શૂન્ય વિધેય,
ઉત્તરગામી વિધેય અને પ્રક્ષેપો છે તથા જે સંયોજન, આદિમ પુનરાવર્તન
અને અનિયત મર્યાદા સુધીની શોધ હેઠળ બંધ છે.
\end{defn}

અલબત્ત, કેટલાંક આંશિક પુનરાવર્તી વિધેયો સંપૂર્ણ પણ નીવડે છે,
એટલે દરેક આર્ગ્યુમેન્ટ માટે વ્યાખ્યાયિત હોય છે.

\begin{defn}
\ollabel{defn:recursive-fn}
\emph{પુનરાવર્તી વિધેયો}નો ગણ એટલે સંપૂર્ણ હોય એવા આંશિક
પુનરાવર્તી વિધેયોનો ગણ.
\end{defn}

પુનરાવર્તી વિધેયને ક્યારેક ``સંપૂર્ણ પુનરાવર્તી'' પણ કહે છે,
જેથી તે દરેક જગ્યાએ વ્યાખ્યાયિત છે તે પર ભાર પડે.
% END OLP-0224
\endgroup


\begingroup
\makeatletter\def\ol@sectioncs{section}\makeatother

% BEGIN OLP-0225
\olfileid{cmp}{rec}{nft}
\olsection{સામાન્ય સ્વરૂપ પ્રમેય}
\phantomsection\label{guunit:OLP-0225}


\begin{thm}[ક્લીનીનું સામાન્ય સ્વરૂપ પ્રમેય]
\ollabel{thm:kleene-nf}
એવો આદિમ પુનરાવર્તી સંબંધ $T(e, x, s)$ અને એવું આદિમ
પુનરાવર્તી વિધેય $U(s)$ છે કે નીચેનો ગુણધર્મ મળે છે: $f$ કોઈ પણ
એક-આર્ગ્યુમેન્ટવાળું આંશિક પુનરાવર્તી વિધેય હોય, તો કોઈ~$e$ માટે
\[
f(x) \simeq U(\umin{s}{T(e, x, s)})
\]
દરેક $x$ માટે સાચું છે.
\sourcecorrection{OLCMP-010}{સ્થિર અંગ્રેજી મૂળ અહીં ``કોઈ પણ
આંશિક પુનરાવર્તી વિધેય'' કહે છે, પરંતુ પ્રદર્શિત સૂત્ર અને
$T(e,x,s)$ એક ઇનપુટ~$x$ માટે લખાયાં છે. ગુજરાતી વાક્ય આ લખેલા
સ્વરૂપને એક-આર્ગ્યુમેન્ટવાળા વિધેય માટે સ્પષ્ટ કરે છે; અનેક
આર્ગ્યુમેન્ટનું સ્વરૂપ જુદાં ટ્યુપલ-સંકેતીકરણથી મળે છે.}
\end{thm}

\begin{explain}
સામાન્ય સ્વરૂપ પ્રમેયની સાબિતી વિગતભરી છે, પરંતુ મૂળ વિચાર સરળ
છે. દરેક આંશિક પુનરાવર્તી વિધેયને એક \emph{સૂચક}~$e$ હોય છે:
સહજ રીતે, તે તેના કાર્યક્રમ અથવા વ્યાખ્યાનો કોડ કરતી સંખ્યા છે.
$f(x)
\fdefined$ હોય, તો ગણતરી વ્યવસ્થિત રીતે નોંધાઈને કોઈ સંખ્યા~$s$
વડે સંકેતીકૃત થઈ શકે. $s$ એ ઇનપુટ~$x$ પર~$f$ની ગણતરીનો કોડ
છે કે નહીં તે માત્ર $x$ અને વ્યાખ્યાના સૂચક~$e$ વડે આદિમ
પુનરાવર્તી રીતે ચકાસી શકાય છે. પરિણામે સંબંધ~$T$---``સૂચક~$e$
ધરાવતા વિધેયની ઇનપુટ~$x$ માટે ગણતરી છે અને $s$ એ ગણતરીનો
કોડ છે''---આદિમ પુનરાવર્તી છે. ગણતરીની સંપૂર્ણ નોંધ~$s$ મળી જાય
તો~$s$નું ``પરિણામ'' એટલે~$f(x)$નું મૂલ્ય, અને તેને પણ~$s$માંથી
આદિમ પુનરાવર્તી રીતે મેળવી શકાય છે.

સામાન્ય સ્વરૂપ પ્રમેય બતાવે છે કે કોઈ પણ આંશિક પુનરાવર્તી વિધેય
વ્યાખ્યાયિત કરવા અનિયત મર્યાદા સુધીની માત્ર એક શોધ જરૂરી છે.
મૂળભૂત રીતે, સૂચક~$e$ ધરાવતા વિધેયની ઇનપુટ~$x$ માટેની ગણતરીનો
કોડ મળે ત્યાં સુધી બધી સંખ્યાઓમાં શોધીએ છીએ. આંશિક પુનરાવર્તી
વિધેયોનાં ``નામો'' તરીકે સંખ્યાઓ~$e$ વાપરી શકીએ અને પ્રમેયના
સમીકરણથી વ્યાખ્યાયિત વિધેય~$f$ માટે $\cfind{e}$ લખી શકીએ.
નોંધો કે એક જ આંશિક પુનરાવર્તી વિધેયને એકથી વધુ સૂચકો હોઈ શકે
છે---હકીકતમાં, દરેકને અનંત સૂચકો હોય છે.
\end{explain}
% END OLP-0225
\endgroup


\begingroup
\makeatletter\def\ol@sectioncs{section}\makeatother

% BEGIN OLP-0226
\olfileid{cmp}{rec}{hlt}
\olsection{અટકવાની સમસ્યા}
\phantomsection\label{guunit:OLP-0226}


સામાન્ય રીતે \emph{અટકવાની સમસ્યા} એટલે સંગણનીય વિધેયનું
વર્ણન~$e$ (દાખલા તરીકે કાર્યક્રમ) અને સંખ્યા~$n$ આપેલાં હોય,
ત્યારે ઇનપુટ~$n$ પર તે વિધેયની ગણતરી અટકે છે કે નહીં, એટલે
પરિણામ આપે છે કે નહીં, તે નક્કી કરવાની સમસ્યા. એલન ટ્યુરિંગે
પ્રસિદ્ધ રીતે સાબિત કર્યું કે આ સમસ્યા પોતે સંગણનીય વિધેયથી
ઉકેલી શકાતી નથી; એટલે નીચેનું વિધેય સંગણનીય નથી:
\[
h(e, n) =
\begin{cases}
1 & \text{જો ગણતરી $e$ ઇનપુટ $n$ પર અટકે}\\
0 & \text{અન્યથા,}
\end{cases}
\]

આંશિક પુનરાવર્તી વિધેયોના સંદર્ભમાં કાર્યક્રમના વર્ણનનું કામ
ક્લીનીના સામાન્ય સ્વરૂપ પ્રમેયમાં મળતો સૂચક~$e$ કરી શકે છે.
$f$ આંશિક પુનરાવર્તી વિધેય હોય, તો સામાન્ય સ્વરૂપ પ્રમેયનું
સમીકરણ સંતોષતો કોઈ પણ $e$ એ~$f$નો સૂચક છે. સંખ્યા~$e$ આપેલી
હોય, તો સામાન્ય સ્વરૂપ પ્રમેય મુજબ
\[
\cfind{e}(x) \simeq U(\mu s \; T(e, x, s))
\]
આંશિક પુનરાવર્તી છે; અને દરેક આંશિક પુનરાવર્તી $f\colon \Nat
\to \Nat$ માટે એવો $e \in \Nat$ છે કે બધા~$x \in \Nat$ માટે
$\cfind{e}(x) \simeq
f(x)$. હકીકતમાં, આવા દરેક $f$ને એક નહીં પણ અનંત એવા~$e$
મળે છે. \emph{અટકવાનું વિધેય}~$h$ આમ વ્યાખ્યાયિત છે:
\[
h(e, x) =
\begin{cases}
1 & \text{જો $\cfind{e}(x) \fdefined$ હોય}\\
0 & \text{અન્યથા.}
\end{cases}
\]
નોંધો કે $\cfind{e}(x) \fundefined$ હોય ત્યારે $h(e, x) = 0$
છે; અને~$e$ બિલકુલ કોઈ આંશિક પુનરાવર્તી વિધેયનો સૂચક ન હોય
ત્યારે પણ એવું જ છે.

\begin{thm}
\ollabel{thm:halting-problem}
અટકવાનું વિધેય~$h$ આંશિક પુનરાવર્તી નથી.
\end{thm}

\begin{proof}
જો $h$ આંશિક પુનરાવર્તી હોત, તો નીચે મુજબ વિધેય વ્યાખ્યાયિત
કરી શકાત:
\[
d(y) =
\begin{cases}
1 & \text{જો $h(y, y) = 0$ હોય}\\
\umin{x}{x \neq x} & \text{અન્યથા.}
\end{cases}
\]
કોઈ સંખ્યા $x$ એ $x \neq x$ સંતોષતી નથી. તેથી
$\umin{x}{x \neq
x}$ અસ્તિત્વમાં નથી અને $h(y,y) \neq 0$ હોય ત્યારે અને માત્ર
ત્યારે જ $d(y) \fundefined$ છે. આ વ્યાખ્યામાંથી મળે છે:
\begin{enumerate}
\item $\cfind{y}(y) \fundefined$ હોય અથવા $y$ એ આંશિક
  પુનરાવર્તી વિધેયનો સૂચક ન હોય ત્યારે અને માત્ર ત્યારે જ
  $d(y) \fdefined$ છે.
\item $\cfind{y}(y) \fdefined$ હોય ત્યારે અને માત્ર ત્યારે જ
  $d(y) \fundefined$ છે.
\end{enumerate}
જો $h$ આંશિક પુનરાવર્તી હોત, તો $d$ પણ આંશિક પુનરાવર્તી હોત.
તેથી ક્લીનીના સામાન્ય સ્વરૂપ પ્રમેય મુજબ તેને સૂચક~$e_d$ મળે.
$h(e_d, e_d)$નું મૂલ્ય વિચારો. તેના બે શક્ય કિસ્સા છે: $0$ અને~$1$.
\begin{enumerate}
\item જો $h(e_d, e_d) = 1$ હોય, તો $\cfind{e_d}(e_d)
  \fdefined$. પરંતુ $\cfind{e_d} \simeq d$ છે; અને
  $h(e_d, e_d) = 0$ હોય ત્યારે અને માત્ર ત્યારે જ $d(e_d)$
  વ્યાખ્યાયિત છે. તેથી $h(e_d, e_d) \neq 1$.
\item જો $h(e_d, e_d) = 0$ હોય, તો કાં તો $e_d$ આંશિક પુનરાવર્તી
  વિધેયનો સૂચક નથી, અથવા છે અને $\cfind{e_d}(e_d)
  \fundefined$. પરંતુ ફરી $\cfind{e_d} \simeq d$ છે; અને
  $\cfind{e_d}(e_d) \fdefined$ હોય ત્યારે અને માત્ર ત્યારે જ
  $d(e_d)$ અવ્યાખ્યાયિત છે.
\end{enumerate}
અંતે એવું નીકળે છે કે $e_d$ આંશિક પુનરાવર્તી વિધેયનો સૂચક હોઈ
જ ન શકે. પરંતુ જો $h$ આંશિક પુનરાવર્તી હોત, તો $d$ પણ હોત;
એટલે $e_d$ને તેનો સૂચક માનતી આપણી વ્યાખ્યા માન્ય હોત. તેથી
$h$ આંશિક પુનરાવર્તી હોઈ શકે નહીં.
\end{proof}
% END OLP-0226
\endgroup


\begingroup
\makeatletter\def\ol@sectioncs{section}\makeatother

% BEGIN OLP-0227
\olfileid{cmp}{rec}{gen}
\olsection{સામાન્ય પુનરાવર્તી વિધેયો}
\phantomsection\label{guunit:OLP-0227}


સંપૂર્ણ વિધેયોનો ગણ મેળવવાની બીજી રીત પણ છે. સંપૂર્ણ વિધેય
$f(x,\vec z)$ને \emph{નિયમિત} કહીએ, જો પ્રાકૃતિક સંખ્યાઓના દરેક
અનુક્રમ $\vec z$ માટે એવો $x$ હોય કે $f(x,\vec z) = 0$ થાય.
બીજા શબ્દોમાં, નિયમિત વિધેયો ચોક્કસ એવાં વિધેયો છે જેના પર
અનિયત મર્યાદા સુધીની શોધ લાગુ કરતાં સંપૂર્ણ વિધેય મળે છે.
સાવધાનીપૂર્વક, અનિયત મર્યાદા સુધીની શોધને નિયમિત વિધેયો પૂરતી
મર્યાદિત કરી શકીએ:

\begin{defn}
\ollabel{defn:general-recursive}
\emph{સામાન્ય પુનરાવર્તી વિધેયો}નો ગણ એટલે પ્રાકૃતિક
સંખ્યાઓમાંથી પ્રાકૃતિક સંખ્યાઓમાં જતાં (જુદી જુદી
આર્ગ્યુમેન્ટ-સંખ્યા ધરાવતાં) વિધેયોનો સૌથી નાનો એવો ગણ, જેમાં
શૂન્ય વિધેય, ઉત્તરગામી વિધેય અને પ્રક્ષેપો છે તથા જે સંયોજન,
આદિમ પુનરાવર્તન અને \emph{નિયમિત} વિધેયો પર લાગુ પડતી અનિયત
મર્યાદા સુધીની શોધ હેઠળ બંધ છે.
\end{defn}

સ્પષ્ટ છે કે દરેક સામાન્ય પુનરાવર્તી વિધેય સંપૂર્ણ છે.
\olref{defn:general-recursive} અને
\olref[par]{defn:recursive-fn} વચ્ચેનો તફાવત એ છે કે પછીની
વ્યાખ્યામાં વચ્ચેના પગલાઓમાં આંશિક પુનરાવર્તી વિધેયો વાપરવાની
છૂટ છે; અંતે મળતું વિધેય સંપૂર્ણ હોય એટલી જ શરત છે. તેથી
``સામાન્ય'' શબ્દ ઐતિહાસિક અવશેષ છે અને નામ તરીકે ભ્રામક છે:
ઉપરથી જોતાં \olref{defn:general-recursive} એ
\olref[par]{defn:recursive-fn} કરતાં \emph{ઓછી} સામાન્ય લાગે છે.
પરંતુ સદનસીબે આ તફાવત માત્ર દેખીતો છે: વ્યાખ્યાઓ જુદી હોવા
છતાં સામાન્ય પુનરાવર્તી વિધેયોનો ગણ અને પુનરાવર્તી વિધેયોનો
ગણ એક જ છે.
% END OLP-0227
\endgroup


\OLEndChapterHook
% END OLP-0209
\endgroup


\begingroup
\makeatletter\def\ol@sectioncs{section}\makeatother

% BEGIN OLP-0228
\olchapter{cmp}{thy}{સંગણનીયતા સિદ્ધાંત}
\olchapterid{cmp}{thy}
\phantomsection\label{guunit:OLP-0228}


\begin{editorial}
  આ પ્રકરણની સામગ્રીની સમીક્ષા કરીને તેને વિસ્તૃત કરવાની જરૂર છે.
  ખાસ કરીને, તેમાં હજી કસરતો નથી.
\end{editorial}

\begingroup
\makeatletter\def\ol@sectioncs{section}\makeatother

% BEGIN OLP-0229
\olfileid{cmp}{thy}{int}
\olsection{પરિચય}
\phantomsection\label{guunit:OLP-0229}


તર્કશાસ્ત્રની \emph{સંગણનીયતા સિદ્ધાંત} નામની શાખા વિધેયો અને ગણોની
સંગણનીયતા અથવા સાપેક્ષ સંગણનીયતા અંગેના પ્રશ્નોનો અભ્યાસ કરે છે.
આ વિષય પર ક્લીનીનો પ્રભાવ એવો રહ્યો છે કે પહેલાં તેને \emph{પુનરાવર્તન
  સિદ્ધાંત} કહેવામાં આવતો હતો; આજે બંને નામ સામાન્ય રીતે વપરાય છે.

જો સંગણનાના કોઈ નિદર્શમાં $f\colon \Nat \pto \Nat$ પ્રકારનું વિધેય
ગણી શકાય, તો તેને \emph{આંશિક સંગણનીય} કહીએ. $f$ સંપૂર્ણ હોય તો
$f$ને માત્ર \emph{સંગણનીય} કહીએ. જે સંબંધ~$R$નું લક્ષણ વિધેય
$\Char{R}$ સંગણનીય હોય તેને પણ સંગણનીય કહીએ. $f$ અને~$g$ આંશિક
વિધેયો હોય, ત્યારે $f(x) \fdefined$નો અર્થ $x$ પર $f$ વ્યાખ્યાયિત
છે, એટલે $x$ $f$ના પ્રદેશમાં છે, એવો થશે; અને $f(x) \fundefined$નો
અર્થ $f$ $x$ પર વ્યાખ્યાયિત નથી એવો થશે. $f(x) \simeq g(x)$થી એવો અર્થ લઈશું કે
$f(x)$ અને $g(x)$ બંને અનિર્ધારિત છે, અથવા બંને વ્યાખ્યાયિત અને
સમાન છે.

સંગણનાના કોઈ એક ચોક્કસ નિદર્શનો ઉલ્લેખ કર્યા વિના પણ સંગણનીયતાનો
સિદ્ધાંત તપાસી શકાય છે. તે માટે સર્વવ્યાપી આંશિક સંગણનીય વિધેય
$\fn{Un}(k,x)$ છે એવું બતાવાય છે. તેના વડે આંશિક સંગણનીય વિધેયોની
પરિગણના કરી શકાય છે. $k$મા એક-આર્ગ્યુમેન્ટવાળા આંશિક સંગણનીય
વિધેય માટે $\cfind{k}$ સંકેત વાપરીશું, જ્યાં
$\cfind{k}(x) \simeq \fn{Un}(k,x)$ છે. (ક્લીનીએ આ માટે
$\{ k \}$ વાપર્યું હતું, પરંતુ તાજેતરમાં આ સંકેત ઓછો વપરાય છે.)
થોડું વધુ સામાન્ય રીતે, કોઈ પણ આર્ગ્યુમેન્ટ-સંખ્યાવાળાં આંશિક
સંગણનીય વિધેયોની એકસરખી રીતથી પરિગણના કરી શકાય છે. $k$મા
$n$-આર્ગ્યુમેન્ટવાળા આંશિક પુનરાવર્તી વિધેય માટે
$\cfind{k}[n]$ લખીશું.

$f(\vec x,y)$ સંપૂર્ણ અથવા આંશિક વિધેય હોય, તો
$\umin{y}{f(\vec x,y)}$ એ $\vec x$નું એવું વિધેય છે જે સૌથી નાનું
$y$ આપે છે જેના માટે $f(\vec x,y)=0$ છે, શરતે કે
$f(\vec x,0)$, \dots, $f(\vec x,y-1)$ બધાં વ્યાખ્યાયિત હોય.
આવું $y$ ન હોય તો $\umin{y}{f(\vec x,y)}$ અનિર્ધારિત છે.
$R(\vec x,y)$ સંબંધ હોય, તો $\umin{y}{R(\vec x,y)}$ એ
$R(\vec x,y)$ સાચું હોય એવું સૌથી નાનું~$y$ છે; બીજા શબ્દોમાં,
$1 \tsub \Char{R}(\vec x,y)=0$ હોય એવું સૌથી નાનું~$y$.

વિધેય સંગણનીય છે એવું બતાવવા બે રીતે આગળ વધી શકાય:
\begin{enumerate}
\item ચુસ્ત રીતે: ટ્યુરિંગ મશીન અથવા આંશિક પુનરાવર્તી વિધેય
  સ્પષ્ટ રીતે વર્ણવીને તે ઇચ્છિત વિધેય ગણે છે એવું બતાવવું;
\item અનૌપચારિક રીતે: વિધેય ગણતી કલનવિધિ વર્ણવવી અને ચર્ચની
  માન્યતાનો આધાર લેવો.
\end{enumerate}
આ બે રીતો વચ્ચે તીક્ષ્ણ સીમા નથી. કલનવિધિનું વિગતવાર વર્ણન એવું
માહિતી આપે કે સિદ્ધાંતમાં તેના માટે યોગ્ય ટ્યુરિંગ મશીન કે આંશિક
પુનરાવર્તી વ્યાખ્યાઓની શ્રેણી કેવી રીતે ઘડી શકાય તે પ્રમાણમાં સ્પષ્ટ
થવું જોઈએ. સંપૂર્ણ ઔપચારિક વ્યાખ્યાઓ ઘણી વાર સમજમાં ખાસ વધારો
કરતી નથી. તેથી બંને રીતો વચ્ચે યોગ્ય સમતુલા રાખીશું: ચોક્કસ
વ્યાખ્યાઓ મેળવી શકાય એવી ખાતરી આપતી, છતાં સહજ સમજમાં ઊતરે એવી
ચુસ્ત સાબિતીઓ શોધીશું.
% END OLP-0229
\endgroup


\begingroup
\makeatletter\def\ol@sectioncs{section}\makeatother

% BEGIN OLP-0230
\olfileid{cmp}{thy}{cod}
\olsection{સંગણનોનું સંકેતીકરણ}
\phantomsection\label{guunit:OLP-0230}


સંગણનાના દરેક નિદર્શમાં નીચેનું કરી શકાય છે:
\begin{enumerate}
\item સંગણનીય વિધેયોની \emph{વ્યાખ્યાઓ} વ્યવસ્થિત રીતે વર્ણવી શકાય.
  દાખલા તરીકે, ટ્યુરિંગ મશીનનાં વિનિર્દેશો, પુનરાવર્તી વ્યાખ્યાઓ
  અથવા કોઈ પ્રોગ્રામિંગ ભાષાના પ્રોગ્રામો આવી વ્યાખ્યાઓ આપે છે.
\item કોઈ આપેલી વ્યાખ્યા અને ઇનપુટ માટે વિધેયની સંગણનાનો પૂરો
  ઇતિહાસ વર્ણવી શકાય. દાખલા તરીકે, ટ્યુરિંગ મશીનની સંગણનાને
  તેના દરેક તબક્કાની ગોઠવણીના અનુક્રમથી વર્ણવી શકાય છે; તેમાં
  મશીનની સ્થિતિ અને ટેપ પરનું લખાણ આવે છે.
\item સંગણનાનો દેખીતી ઇતિહાસ ખરેખર આપેલી વ્યાખ્યાવાળું સંગણનીય
  વિધેય આપેલા ઇનપુટ પર કેવી રીતે ગણાય છે તેનો જ ઇતિહાસ છે કે
  નહીં તે તપાસી શકાય; અહીં એ ઇનપુટ પર વિધેય વ્યાખ્યાયિત છે,
  એટલે સંગણના અટકે છે.
\item સંગણનાના પૂરા ઇતિહાસના આવા વર્ણનમાંથી આપેલા ઇનપુટ માટે
  વિધેયનું મૂલ્ય કાઢી શકાય. દાખલા તરીકે, અટકતી ટ્યુરિંગ મશીનની
  સંગણનાના છેલ્લા તબક્કે ટેપ પરનું લખાણ એ મૂલ્ય છે.
\end{enumerate}

સંકેતીકરણથી સંગણનીય વિધેયના દરેક વર્ણનને એક સંખ્યાત્મક
\emph{સૂચક} આપી શકાય છે, એવી રીતે કે એ સૂચકમાંથી સૂચનાઓ
સંગણનીય રીતે ફરી મેળવી શકાય. એ જ રીતે સંગણનાનો પૂરો ઇતિહાસ
એક સંખ્યાથી સંકેતિત કરી શકાય. પરિણામે, ``$s$ એ ઇનપુટ~$x$
માટે સૂચક~$e$ ધરાવતા વિધેયની સંગણનાનો ઇતિહાસ સંકેતિત કરે છે''
એ અંકગણિતીય સંબંધ અને ``સંકેત~$s$ ધરાવતા સંગણના-અનુક્રમનું
પરિણામ'' એ વિધેય બંને સંગણનીય છે; ખરેખર, બંને આદિમ પુનરાવર્તી છે.

આ પાયાનો તથ્ય બહુ શક્તિશાળી છે. તે પસંદ કરેલા સંગણન-નિદર્શનથી
સ્વતંત્ર રીતે સંગણનીયતા વિશેના અનેક આશ્ચર્યકારક અને મહત્વના
પરિણામો સાબિત કરવા દે છે.
% END OLP-0230
\endgroup


\begingroup
\makeatletter\def\ol@sectioncs{section}\makeatother

% BEGIN OLP-0231
\olfileid{cmp}{thy}{nfm}
\olsection{સામાન્ય સ્વરૂપ પ્રમેય}
\phantomsection\label{guunit:OLP-0231}


ધારો કે સંગણનીય વિધેયોની વ્યાખ્યાઓ વર્ણવી શકીએ છીએ અને આપેલા
ઇનપુટ પર તે વિધેયની સંગણનાનો પૂરો ઇતિહાસ દર્શાવતું દેખીતું
વર્ણન સાચું છે કે નહીં તે તપાસી શકીએ છીએ. તો સંગણનાના નિદર્શનથી
સ્વતંત્ર રીતે, કોઈ પણ સંગણનીય વિધેય~$f$નું ઇનપુટ~$x$ પરનું
મૂલ્ય નીચેની રીતે નક્કી કરી શકીએ એવું સ્વાભાવિક લાગે છે:
\begin{enumerate}
  \item સંગણનાના ઇતિહાસનાં શક્ય બધાં વર્ણનોમાં શોધ કરવી.
  \item આપેલો ઇતિહાસ $f(x)$ની સંગણનાનો ઇતિહાસ છે કે નહીં તે તપાસવું.
  \item યોગ્ય ઇતિહાસ મળે તો તેમાંથી $f(x)$નું મૂલ્ય કાઢવું.
\end{enumerate}
ખરેખર આવું થઈ શકે છે તે ક્લીનીના સામાન્ય સ્વરૂપ પ્રમેયનો વિષય છે.

\begin{thm}[ક્લીનીનું સામાન્ય સ્વરૂપ પ્રમેય]
\ollabel{thm:normal-form}
આવો આદિમ પુનરાવર્તી સંબંધ~$T(e,x,s)$ અને આદિમ પુનરાવર્તી
વિધેય~$U(s)$ છે કે દરેક એક-આર્ગ્યુમેન્ટવાળા આંશિક સંગણનીય
વિધેય~$f$ માટે કોઈ~$e$ મળે છે અને દરેક~$x$ માટે
\[
f(x) \simeq U(\umin{s}{T(e,x,s)})
\]
થાય છે.
\sourcecorrection{OLCT-001}{સ્થિર અંગ્રેજી મૂળ ``કોઈ પણ આંશિક
સંગણનીય વિધેય'' કહે છે, પરંતુ પ્રદર્શિત $f(x)$ અને $T(e,x,s)$
એક-આર્ગ્યુમેન્ટવાળું સ્વરૂપ આપે છે. અહીં પ્રમેય એ ક્ષેત્રમાં
સ્પષ્ટ કર્યો છે; વધુ આર્ગ્યુમેન્ટો માટે ટ્યુપલ-કોડ વાપરી શકાય.}
\end{thm}

\begin{proof}[સાબિતીની રૂપરેખા]
સંગણનાના કોઈ પણ નિદર્શન માટે સંગણનીય વિધેય~$f$નું વર્ણન ચુસ્ત રીતે
વ્યાખ્યાયિત કરી શકાય અને તેને પ્રાકૃતિક સંખ્યા~$e$થી સંકેતિત કરી
શકાય. સૂચક~$e$વાળા વિધેયની ઇનપુટ~$x$ પરની સંગણનાની પ્રક્રિયા
નોંધતો ``સંગણના-અનુક્રમ'' પણ ચુસ્ત રીતે વ્યાખ્યાયિત કરી શકાય. એ
અનુક્રમને પણ એક સંખ્યા~$s$થી સંકેતિત કરી શકાય. આ બધું એવી રીતે
કરી શકાય કે
\begin{enumerate}
  \item $T(e,x,s)$ સંબંધ ત્યારે અને માત્ર ત્યારે સાચો હોય જ્યારે
  સંખ્યા~$s$ સૂચક~$e$વાળા વિધેયની ઇનપુટ~$x$ પરની સંગણનાનો
  અનુક્રમ સંકેતિત કરે; અને
  \item $U(s)$ વિધેય $s$થી સંકેતિત સંગણના-અનુક્રમને તે
  સંગણનાના અંતિમ પરિણામ સાથે જોડે.
\end{enumerate}
આ બંને સંગણનીય છે. ખરેખર, $T$ સંબંધ અને $U$ વિધેય આદિમ
પુનરાવર્તી છે.
\end{proof}

\begin{explain}
સામાન્ય સ્વરૂપ પ્રમેયની ચુસ્ત સાબિતી માટે સંગણનાનું એક નિદર્શન
નક્કી કરીને સંગણનીય વિધેયોનાં વર્ણનો તથા સંગણના-અનુક્રમોનું
સંકેતીકરણ વિગતવાર કરવું પડે અને $T$ સંબંધ તથા $U$ વિધેય
આદિમ પુનરાવર્તી છે તે તપાસવું પડે. મોટા ભાગના ઉપયોગોમાં $T$
અને~$U$ સંગણનીય છે તથા $U$ સંપૂર્ણ છે એટલું પૂરતું થાય છે.

પ્રમેયની સાબિતી સૂત્રરૂપે આમ યાદ રાખવી સારી: $\umin{s}{T(e,x,s)}$
સૂચક~$e$વાળા વિધેયની ઇનપુટ~$x$ પરની સંગણનાનો અનુક્રમ શોધે છે,
અને એવો અનુક્રમ મળે તો $U$ તેનું પરિણામ આપે છે.
\end{explain}

જો સંગણનાનું નિદર્શન આંશિક પુનરાવર્તી વિધેયોનું હોય (મૂળે
ક્લીનીએ એ જ વાપર્યું હતું), તો આ પ્રમેય બતાવે છે કે કોઈ પણ
વિધેયની વ્યાખ્યા માટે અનિયત મર્યાદા સુધીની શોધનો એક જ ઉપયોગ,
એટલે એક જ $\umin{y}{f(\vec x,y)}$ ક્રિયા, જરૂરી છે. આ અર્થમાં
દરેક આંશિક પુનરાવર્તી વિધેયને સામાન્ય સ્વરૂપ મળે છે.

$T$ અને $U$ વડે $\cfind{0}$, $\cfind{1}$, $\cfind{2}$, \dots{}ની
પરિગણના વ્યાખ્યાયિત કરી શકાય. હવેથી $T$ અને~$U$ની યોગ્ય પસંદગી
નક્કી થયેલી માનીશું અને
\[
\cfind{e}(x) \simeq U(\umin{s}{T(e,x,s)})
\]
એ સમીકરણને $\cfind{e}$ની \emph{વ્યાખ્યા} ગણીશું.

બીજું ઉપયોગી તથ્ય આ છે:

\begin{thm}
દરેક આંશિક સંગણનીય વિધેયને અનંત ઘણાં સૂચકો હોય છે.
\end{thm}

આ વાત સહજ રીતે પણ સ્પષ્ટ છે. સંગણનીય વિધેયના કોઈ વર્ણનમાંથી
એ જ વિધેય ગણતું બીજું વર્ણન બનાવી શકાય: દાખલા તરીકે, પહેલાં
સંગણનાને અસર ન કરે એવું કંઈક કરવું (જેમ કે $0=0$ છે કે નહીં
તપાસવું અથવા $5$ સુધી ગણવું). બદલાયેલા વર્ણનનો સૂચક મૂળના
સૂચકથી હંમેશાં જુદો હશે. છતાં બંને એક જ વિધેયના સૂચકો છે;
માત્ર તેની સંગણના થોડી જુદી રીતે થાય છે.
% END OLP-0231
\endgroup


\begingroup
\makeatletter\def\ol@sectioncs{section}\makeatother

% BEGIN OLP-0232
\olfileid{cmp}{thy}{smn}
\olsection{$s$-$m$-$n$ પ્રમેય}
\phantomsection\label{guunit:OLP-0232}


\begin{explain}
આગળનું પ્રમેય ``$s$-$m$-$n$ પ્રમેય'' તરીકે ઓળખાય છે; તેનું કારણ
થોડી વારમાં સ્પષ્ટ થશે. અઘરું કામ પ્રમેય ખરેખર શું કહે છે તે
સમજવાનું છે; વિધાન સમજાઈ જાય પછી તે ઘણું સ્વાભાવિક લાગશે.
\end{explain}

\begin{thm}
\ollabel{thm:s-m-n}
  પ્રાકૃતિક સંખ્યાઓની દરેક જોડી $n$ અને~$m$ માટે એવું આદિમ
  પુનરાવર્તી વિધેય~$s^m_n$ છે કે દરેક અનુક્રમ
  $e$, $a_0$, \dots, $a_{m-1}$, $y_0$, \dots, $y_{n-1}$ માટે
  \[
  \cfind{s^m_n(e, a_0, \dots, a_{m-1})}[n](y_0, \dots, y_{n-1}) \simeq
  \cfind{e}[m+n](a_0, \dots, a_{m-1}, y_0, \dots, y_{n-1}).
\]
\end{thm}

\begin{explain}
$s^m_n$ને \emph{પ્રોગ્રામો} પર કાર્ય કરતું માનવાથી મદદ મળે છે.
એટલે કે, $s^m_n$ને $(m+n)$-આર્ગ્યુમેન્ટવાળા વિધેયનો
પ્રોગ્રામ~$e$ અને નિશ્ચિત ઇનપુટ $a_0$, \dots, $a_{m-1}$ મળે;
તે બાકીના આર્ગ્યુમેન્ટોના $n$-આર્ગ્યુમેન્ટવાળા વિધેય માટે
$s^m_n(e,a_0,\dots,a_{m-1})$ પ્રોગ્રામ આપે છે.
\sourcecorrection{OLCT-002}{સ્થિર અંગ્રેજી મૂળ આ છેલ્લા
પ્રોગ્રામના સૂચકમાં અને ટ્યુરિંગ-મશીનના વર્ણનમાં $e$ને બદલે $x$
લખે છે. ઉપરના પ્રમેય અને તે જ વાક્યના પ્રોગ્રામ~$e$ મુજબ બંને
જગ્યાએ $e$ રાખ્યું છે.}
 જો $e$ને ટ્યુરિંગ મશીનનું વર્ણન માનીએ, તો
$s^m_n(e,a_0,\dots,a_{m-1})$ એ એવું ટ્યુરિંગ મશીન છે જે
$y_0$, \dots, $y_{n-1}$ ઇનપુટ પર ઇનપુટ-શબ્દમાળાની શરૂઆતમાં
$a_0$, \dots, $a_{m-1}$ જોડે છે અને પછી~$e$ ચલાવે છે.
દરેક $s^m_n$ એ યોગ્ય ટ્યુરિંગ મશીનનો કોડ શોધતું આદિમ
પુનરાવર્તી વિધેય છે.
\end{explain}
% END OLP-0232
\endgroup


\begingroup
\makeatletter\def\ol@sectioncs{section}\makeatother

% BEGIN OLP-0233
\olfileid{cmp}{thy}{uni}
\olsection{સર્વવ્યાપી આંશિક સંગણનીય વિધેય}
\phantomsection\label{guunit:OLP-0233}


\begin{thm}
\ollabel{thm:univ-comp}
સર્વવ્યાપી આંશિક સંગણનીય વિધેય~$\fn{Un}(e,x)$ છે. બીજા શબ્દોમાં,
એવું વિધેય $\fn{Un}(e,x)$ છે કે:
\begin{enumerate}
\item $\fn{Un}(e,x)$ આંશિક સંગણનીય છે.
\item $f(x)$ કોઈ પણ આંશિક સંગણનીય વિધેય હોય, તો એવી
પ્રાકૃતિક સંખ્યા $e$ છે કે દરેક~$x$ માટે $f(x) \simeq \fn{Un}(e,x)$.
\end{enumerate}
\end{thm}

\begin{proof}
$\fn{Un}(e,x) \simeq U(\umin{s}{T(e,x,s)})$ રાખો, જ્યાં $U$ અને $T$
ક્લીનીના સામાન્ય સ્વરૂપના પ્રમેય (\olref[nfm]{thm:normal-form})માં
આપેલાં છે.
\end{proof}

\begin{explain}
આ તો આંશિક સંગણનીય વિધેયોની અસરકારક પરિગણના છે એવું ચોક્કસ રીતે
કહે છે. વિચાર એવો છે કે $f_e(x) = \fn{Un}(e,x)$ વડે વ્યાખ્યાયિત
વિધેયને $f_e$ લખીએ, તો અનુક્રમ $f_0$, $f_1$, $f_2$, \dots માં બધાં
આંશિક સંગણનીય વિધેયો આવે છે. વધુમાં $f_e(x)$ની ગણતરી $e$ અને~$x$
બંનેની દૃષ્ટિએ ``એકસરખી રીતે'' થઈ શકે છે. સરળતા ખાતર અહીં એક-આર્ગ્યુમેન્ટવાળાં
વિધેયો માટે સર્વવ્યાપી દ્વિ-આર્ગ્યુમેન્ટવાળું વિધેય લીધું છે; પરંતુ
સંખ્યાઓના અનુક્રમોને કોડ કરીને આ વાત વધુ આર્ગ્યુમેન્ટો માટે સહેલાઈથી
વિસ્તારી શકાય છે. ઉદાહરણ તરીકે, $f(x,y,z)$ જો $3$-સ્થાનીય આંશિક
પુનરાવર્તી વિધેય હોય, તો $g(x) \simeq f((x)_0, (x)_1, (x)_2)$
એક-આર્ગ્યુમેન્ટવાળું પુનરાવર્તી વિધેય છે.
\end{explain}
% END OLP-0233
\endgroup


\begingroup
\makeatletter\def\ol@sectioncs{section}\makeatother

% BEGIN OLP-0234
\olfileid{cmp}{thy}{nou}
\olsection{કોઈ સર્વવ્યાપી સંપૂર્ણ સંગણનીય વિધેય નથી}
\phantomsection\label{guunit:OLP-0234}


આંશિક સંગણનીય વિધેયો માટે સર્વવ્યાપી આંશિક સંગણનીય વિધેય તો છે,
પરંતુ સંપૂર્ણ સંગણનીય વિધેયો માટે સર્વવ્યાપી હોય એવું સંપૂર્ણ
સંગણનીય વિધેય નથી.

\begin{thm}
\ollabel{thm:no-univ}
કોઈ સર્વવ્યાપી સંગણનીય વિધેય નથી. બીજા શબ્દોમાં, જો
$\fn{Un}'(k, x)$ એવું વિધેય હોય કે દરેક સંપૂર્ણ સંગણનીય વિધેય
$f(x)$ માટે એવી પ્રાકૃતિક સંખ્યા~$k$ હોય કે દરેક~$x$ માટે
$f(x) = \fn{Un}'(k,x)$, તો તે સંગણનીય નથી.
\end{thm}

\begin{proof}
આ સાદી વિકર્ણીકરણ સાબિતી છે: $\fn{Un}'(k,x)$ સંપૂર્ણ અને સંગણનીય
હોય, તો
\[
d(x) = \fn{Un}'(x, x) + 1
\]
પણ સંપૂર્ણ અને સંગણનીય થાય. પરંતુ વ્યાખ્યા મુજબ $d(k)$ અને
$\fn{Un}'(k,k)$ સમાન નથી. તેથી દરેક $k$ માટે $d(x)$ અને
$\fn{Un}'(k, x)$નાં મૂલ્યો ઓછામાં ઓછા એક~$x$ પર, એટલે $x = k$
પર, જુદાં પડે છે.
\end{proof}

\begin{explain}
ઉપરનું \olref[uni]{thm:univ-comp} બતાવે છે કે સર્વવ્યાપી વિધેયને
આંશિક રહેવા દઈએ તો આ વિકર્ણીકરણ દલીલથી બચી શકાય છે. એટલે
$\fn{Un}$ સંપૂર્ણ સંગણનીય વિધેયો માટે પણ સર્વવ્યાપી છે, પરંતુ તે
પોતે સંપૂર્ણ નથી. આંશિક સ્થિતિમાં વિકર્ણીકરણ દલીલ ચાલતી નથી.
\end{explain}

\begin{prob}
  \olref{thm:no-univ}ની સાબિતીની વિકર્ણીકરણ દલીલ આંશિક
  સ્થિતિમાં કેમ ચાલતી નથી તે સમજવા, વિધેય
  $f(x) \simeq \fn{Un}(x,x)+1$ લો. શું તે આંશિક સંગણનીય છે?
  જો હા, તો તેનો કોઈ સૂચક~$e$ હશે, એટલે $f(x) \simeq
  \fn{Un}(e,x)$. તો~$f(e)$ વિશે શું કહી શકાય?
\end{prob}
% END OLP-0234
\endgroup


\begingroup
\makeatletter\def\ol@sectioncs{section}\makeatother

% BEGIN OLP-0235
\olfileid{cmp}{thy}{hlt}
\olsection{અટકવાની સમસ્યા}
\phantomsection\label{guunit:OLP-0235}


રચના મુજબ, સર્વવ્યાપી આંશિક સંગણનીય વિધેય~$\fn{Un}(e,x)$ ત્યારે અને
માત્ર ત્યારે વ્યાખ્યાયિત છે જ્યારે~$e$ વડે કોડ કરેલા વિધેયની
સંગણના ઇનપુટ~$x$ માટે મૂલ્ય આપે. આવું થાય છે કે નહીં તે નક્કી
કરી શકીએ કે કેમ એવો પ્રશ્ન સ્વાભાવિક છે. હકીકતમાં તે શક્ય નથી.
સંગણનાના ટ્યુરિંગ મશીન નિદર્શમાં તેનો અર્થ એ થાય કે આપેલું
ટ્યુરિંગ મશીન આપેલા ઇનપુટ પર થંભશે કે નહીં તે સંગણકીય રીતે
અનિર્ણેય છે. આથી આગળનું પ્રમેય ``અટકવાની સમસ્યાની અનિર્ણેયતા''
તરીકે ઓળખાય છે. નીચે બે સાબિતીઓ આપીશું. પહેલી અગાઉની ચર્ચાને
આગળ ધપાવે છે અને બીજી વધુ સીધી છે.

\begin{thm}
\ollabel{thm:halting-problem}
માનો કે
\[
h(e, x)  =
\begin{cases}
1 & \text{જો\/ $\fn{Un}(e, x)$ વ્યાખ્યાયિત હોય} \\
0 & \text{અન્યથા.}
\end{cases}
\]
તો $h$ સંગણનીય નથી.
\end{thm}

\begin{proof}
માનો કે $h$ સંગણનીય છે. તો તેનાથી સર્વવ્યાપી સંપૂર્ણ સંગણનીય
વિધેય વ્યાખ્યાયિત કરી શકાય એવું બતાવીશું. વ્યાખ્યાયિત કરો:
\[
\fn{Un'}(e,x) =
\begin{cases}
\fn{Un}(e,x) & \text{જો $h(e,x) = 1$} \\
0 & \text{અન્યથા.}
\end{cases}
\]
હવે $\fn{Un'}(e, x)$ સંપૂર્ણ વિધેય છે અને $h$ સંગણનીય હોય તો
તે પણ સંગણનીય છે. ઉદાહરણ તરીકે, નીચે પ્રમાણે આદિમ પુનરાવર્તનથી
$g$ વ્યાખ્યાયિત કરી શકાય:
\begin{align*}
g(0, e, x) & \simeq 0 \\
g(y+1, e, x) & \simeq \fn{Un}(e,x);
\end{align*}
ત્યારે
\[
\fn{Un'}(e,x) \simeq g(h(e,x),e,x).
\]
જ્યાં $\fn{Un}(e,x)$ વ્યાખ્યાયિત છે ત્યાં $\fn{Un'}(e,x)$ તેની
સાથે મેળ ખાય છે. આથી જે આંશિક સંગણનીય વિધેયો હકીકતમાં સંપૂર્ણ
છે, તેમના માટે $\fn{Un'}$ સર્વવ્યાપી છે. પરંતુ આ
\olref[nou]{thm:no-univ} સાથે વિરોધાભાસ છે.
\end{proof}

\begin{proof}
માનો કે $h(e,x)$ સંગણનીય છે. વિધેય $g$ આ રીતે વ્યાખ્યાયિત કરો:
\[
g(x) =
\begin{cases}
  0                & \text{જો $h(x,x) = 0$} \\
  \fundefined & \text{અન્યથા.}
\end{cases}
\]
$g$ આંશિક સંગણનીય છે. દાખલા તરીકે, તેને
$\umin{y}{h(x,x) = 0}$ તરીકે વ્યાખ્યાયિત કરી શકાય. તેથી કોઈ~$e$
માટે દરેક~$x$ પર $g(x) \simeq \fn{Un}(e,x)$ છે. શું $g$, $e$ પર
વ્યાખ્યાયિત છે? જો હોય, તો $g$ની વ્યાખ્યા મુજબ $h(e,e) = 0$ છે
($h$ વ્યાખ્યાયિત હોય ત્યારે જ~$0$ મૂલ્ય લઈ શકે છે).
$h$ની વ્યાખ્યા મુજબ તેથી $\fn{Un}(e, e)$ અનિર્ધારિત છે.
દરેક~$x$ માટે $g(x) \simeq \fn{Un}(e, x)$ હોવાની આપણી ધારણા
મુજબ $g(e)$ અનિર્ધારિત થાય, જે વિરોધાભાસ છે. બીજી તરફ, જો
$g(e)$ અનિર્ધારિત હોય, તો $h(e,e) \neq 0$, એટલે
$h(e,e) = 1$. તેથી $\fn{Un}(e, e)$ વ્યાખ્યાયિત છે. પરંતુ
$g(x) \simeq \fn{Un}(e, x)$ હોવાથી $g(e)$ પણ વ્યાખ્યાયિત
હોવું જોઈએ. આ પણ વિરોધાભાસ છે.
\end{proof}

\begin{explain}
આ દલીલ ટ્યુરિંગ મશીનોની ભાષામાં પણ કહી શકાય. માનો કે એવું
ટ્યુરિંગ મશીન~$H$ છે જે ટ્યુરિંગ મશીન~$E$નું વર્ણન અને ઇનપુટ~$x$
લઈને, $E$ ઇનપુટ~$x$ પર થંભે છે કે નહીં તે નક્કી કરે છે.
તો એક ઇનપુટ~$x$ લેતું બીજું ટ્યુરિંગ મશીન~$G$ બનાવી શકાય:
તે $H$ ચલાવીને તપાસે કે સૂચક~$x$ ધરાવતું મશીન $M_x$
ઇનપુટ~$x$ પર થંભે છે કે નહીં, અને પછી વિરુદ્ધ કરે. એટલે કે,
જો $H$ કહે કે $M_x$ ઇનપુટ~$x$ પર થંભે છે, તો $G$ અનંત લૂપમાં
જાય છે; અને જો $H$ કહે કે $M_x$ ઇનપુટ~$x$ પર થંભતું નથી,
તો $G$ થંભી જાય છે. પોતાના સૂચકને ઇનપુટ તરીકે આપીએ ત્યારે
$G$ થંભે છે? ઉપરની દલીલ કહે છે કે તે ત્યારે અને માત્ર ત્યારે
થંભે છે જ્યારે તે થંભતું નથી---વિરોધાભાસ. તેથી એવું ટ્યુરિંગ
મશીન~$H$ છે એવી આપણી ધારણા ખોટી છે.
\end{explain}
% END OLP-0235
\endgroup


\begingroup
\makeatletter\def\ol@sectioncs{section}\makeatother

% BEGIN OLP-0236
\olfileid{cmp}{thy}{rus}
\olsection{રસેલના વિરોધાભાસ સાથે સરખામણી}
\phantomsection\label{guunit:OLP-0236}


આ વિભાગની દલીલોને રસેલના વિરોધાભાસ સાથે સરખાવવાથી વાત વધુ
સ્પષ્ટ થાય છે:
\begin{enumerate}
\item રસેલનો વિરોધાભાસ: $S = \Setabs{x}{x \notin x}$ રાખો. તો
  $S \in S$ ત્યારે અને માત્ર ત્યારે જ્યારે $S \notin S$,
  જે વિરોધાભાસ છે.
  \sourcecorrection{OLCT-003}{સ્થિર અંગ્રેજી મૂળ અહીં
  $X \notin S$ લખે છે; અગાઉ વ્યાખ્યાયિત કરેલો ગણ $S$ છે અને
  સ્વ-સભ્યતાના આ વિરોધાભાસમાં $S \notin S$ જ જોઈએ.}

  \emph{નિષ્કર્ષ:} આવો ગણ~$S$ નથી. ``બધા ગણોનો ગણ'' અસ્તિત્વમાં
  છે એમ માનવું ગણસિદ્ધાંતના બીજા સ્વયંસિદ્ધાંતો સાથે અસંગત છે.

\item રસેલના વિરોધાભાસનું એક રૂપાંતર: બધા વિધેયોના ગણમાંથી
  $\{ 0, 1 \}$માં જતું ``વિધેય'' $F$ નીચે પ્રમાણે વ્યાખ્યાયિત કરો:
  \[
  F(f) =
  \begin{cases}
    1 & \text{જો $f$ પોતે $f$ના પ્રદેશમાં હોય અને $f(f) = 0$} \\
    0 & \text{અન્યથા}
  \end{cases}
  \]
  સમાન દલીલ બતાવે છે કે $F(F) = 0$ ત્યારે અને માત્ર ત્યારે જ્યારે
  $F(F) = 1$, જે વિરોધાભાસ છે.

  \emph{નિષ્કર્ષ:} $F$ વિધેય નથી. ``બધા વિધેયોનો ગણ'' કોઈ વિધેયનો
  પ્રદેશ બની શકે એટલો નાનો નથી.

\item વિકર્ણીકરણ દલીલ: આંશિક સંગણનીય વિધેયોની પરિગણના
  $f_0$, $f_1$, \dots લો અને $G \colon \Nat \to
  \{ 0, 1 \}$ને નીચે પ્રમાણે વ્યાખ્યાયિત કરો:
  \[
  G(x) =
  \begin{cases}
    1 & \text{જો $f_x(x)\downarrow = 0$} \\
    0 & \text{અન્યથા}
  \end{cases}
  \]
  $G$ સંગણનીય હોય, તો કોઈ $k$ માટે તે વિધેય $f_k$ જ હશે.
  પરંતુ પછી $G(k) = 1$ ત્યારે અને માત્ર ત્યારે જ્યારે $G(k) = 0$,
  જે વિરોધાભાસ છે.

  \emph{નિષ્કર્ષ:} $G$ સંગણનીય નથી. નોંધો કે ગણસિદ્ધાંતના
  સ્વયંસિદ્ધાંતો મુજબ $G$ હજી પણ વિધેય છે; અહીં વિરોધાભાસ
  નથી, માત્ર ભેદ વધુ સ્પષ્ટ થાય છે.
\end{enumerate}

આંશિક વિધેયો, સંગણનીય વિધેયો, આંશિક સંગણનીય વિધેયો વગેરેની
ચર્ચામાં ગૂંચવણ થઈ શકે છે. $\Nat$માંથી $\Nat$માં જતાં બધાં
આંશિક વિધેયોનો ગણ ઘણાં પદાર્થો ધરાવે છે. તેમાંના કેટલાંક સંપૂર્ણ
છે, કેટલાંક સંગણનીય છે, કેટલાંક સંપૂર્ણ અને સંગણનીય બંને છે અને
કેટલાંક બેમાંથી એકેય નથી. યાદ રાખો કે માત્ર ``વિધેય'' કહીએ ત્યારે
સામાન્ય રીતે સંપૂર્ણ વિધેયનો અર્થ લઈએ છીએ. આમ, નીચેના પ્રકારો છે:
\begin{enumerate}
\item સંગણનીય વિધેયો
\item સંપૂર્ણ ન હોય એવા આંશિક સંગણનીય વિધેયો
\item સંગણનીય ન હોય એવા વિધેયો
\item સંપૂર્ણ પણ ન હોય અને સંગણનીય પણ ન હોય એવા આંશિક વિધેયો
\end{enumerate}
આ ભેદ સમજવા માટે $\Nat$માંથી $\Nat$માં જતાં બધા આંશિક વિધેયો
દર્શાવતો મોટો ચોરસ દોરો અને તેમાં અનુક્રમે સંપૂર્ણ વિધેયો તથા
સંગણનીય આંશિક વિધેયો દર્શાવતા બે આચ્છાદિત પ્રદેશો ચિહ્નિત કરો.
આ રેખાંકનમાં મળતા દરેક પ્રદેશમાં આવતો એક પદાર્થ વર્ણવી શકો કે
નહીં તે તપાસવું સારી કસરત છે.
% END OLP-0236
\endgroup


\begingroup
\makeatletter\def\ol@sectioncs{section}\makeatother

% BEGIN OLP-0237
\olfileid{cmp}{thy}{cps}
\olsection{સંગણનીય ગણો}
\phantomsection\label{guunit:OLP-0237}


સંગણનીયતાની કલ્પનાને સંગણનીય વિધેયોથી સંગણનીય ગણો સુધી
વિસ્તારી શકાય છે:

\begin{defn}
  માનો કે $S$ પ્રાકૃતિક સંખ્યાઓનો ગણ છે. તો $S$ ત્યારે અને માત્ર
  ત્યારે \emph{સંગણનીય} છે જ્યારે તેનું લક્ષણ વિધેય~$\Char{S}$
  સંગણનીય હોય. બીજા શબ્દોમાં, $S$ ત્યારે અને માત્ર ત્યારે
  સંગણનીય છે જ્યારે
\[
\Char{S}(x) =
\begin{cases}
1 & \text{જો $x \in S$} \\
0 & \text{અન્યથા}
\end{cases}
\]
વિધેય સંગણનીય હોય. એ જ રીતે, સંબંધ $R(x_0, \dots, x_{k-1})$
ત્યારે અને માત્ર ત્યારે સંગણનીય છે જ્યારે તેનું લક્ષણ વિધેય
સંગણનીય હોય.

સંગણનીય ગણો અને સંબંધોને \emph{નિર્ણેય} પણ કહે છે.
\end{defn}

\begin{explain}
હવે સંગણનીયતાની ઘણી કલ્પનાઓ છે: આંશિક વિધેયો માટે, વિધેયો
માટે અને ગણો માટે. તેમને ગૂંચવશો નહીં!{} આંશિક વિધેય
  ગણતું ટ્યુરિંગ મશીન, વિધેય જ્યાં વ્યાખ્યાયિત છે એવા ઇનપુટ માટે,
  વિધેયનું આઉટપુટ આપે છે; ગણ નક્કી કરતું ટ્યુરિંગ મશીન ઇનપુટ
  તે ગણમાં છે કે નહીં તે નક્કી કરીને $1$ અથવા~$0$ આપે છે.
\end{explain}
% END OLP-0237
\endgroup


\begingroup
\makeatletter\def\ol@sectioncs{section}\makeatother

% BEGIN OLP-0238
\olfileid{cmp}{thy}{ces}
\olsection{સંગણકીય રીતે પરિગણનીય ગણો}
\phantomsection\label{guunit:OLP-0238}


\begin{defn}
કોઈ ગણ ખાલી હોય અથવા કોઈ સંગણનીય વિધેયનો મૂલ્યગણ હોય તો તે
\emph{સંગણકીય રીતે પરિગણનીય} કહેવાય.
\end{defn}

\begin{history}
સંગણકીય રીતે પરિગણનીય ગણોને \emph{પુનરાવર્તી રીતે પરિગણનીય}
પણ કહે છે. આ મૂળ પરિભાષા છે. આજે બંને નામો તથા ``c.e.'' અને
``r.e.'' સંક્ષેપો સામાન્ય રીતે વપરાય છે.
\end{history}

\begin{explain}
આ વ્યાખ્યાનો અર્થ શું થાય છે અને આ પરિભાષા કેમ યોગ્ય છે તે
વિચારો. વિચાર એવો છે કે $S$ સંગણનીય વિધેય~$f$નો મૂલ્યગણ હોય, તો
\[
S = \{ f(0), f(1), f(2), \dots \},
\]
અને તેથી $f$, $S$ના તત્ત્વોની ``પરિગણના'' કરતું વિધેય છે.
નોંધો કે વ્યાખ્યા મુજબ $f$ વધતું વિધેય હોવું જરૂરી નથી; એટલે
પરિગણના ચડતા ક્રમમાં હોવી જરૂરી નથી. હકીકતમાં, $f$ એક-એક
હોવું પણ જરૂરી નથી; એટલે $S$ની પરિગણનામાં $f(0)$, $f(1)$,
$f(2)$, \dots{} નું પુનરાવર્તન ચાલે છે. દાખલા તરીકે,
અચળ વિધેય $f(x) = 0$, ગણ $\{ 0 \}$ની પરિગણના કરે છે.
\end{explain}

દરેક સંગણનીય ગણ સંગણકીય રીતે પરિગણનીય છે. આ જોવા માટે
માનો કે $S$ સંગણનીય છે. $S$ ખાલી હોય, તો વ્યાખ્યા મુજબ તે
સંગણકીય રીતે પરિગણનીય છે. અન્યથા $S$નું કોઈ તત્ત્વ $a$ લો.
$f$ને નીચે પ્રમાણે વ્યાખ્યાયિત કરો:
\[
f(x) =
\begin{cases}
x & \text{જો $\Char{S}(x) = 1$} \\
a & \text{અન્યથા.}
\end{cases}
\]
ત્યારે $f$ સંગણનીય વિધેય છે અને $S$ એ $f$નો મૂલ્યગણ છે.
% END OLP-0238
\endgroup


\begingroup
\makeatletter\def\ol@sectioncs{section}\makeatother

% BEGIN OLP-0239
\olfileid{cmp}{thy}{eqc}

\olsection[પરિગણનીય ગણોની વ્યાખ્યાઓ]{સંગણકીય રીતે પરિગણનીય
  ગણોની સમકક્ષ વ્યાખ્યાઓ}
\phantomsection\label{guunit:OLP-0239}



આગળની વાત સંગણકીય રીતે પરિગણનીય હોવાના અર્થનાં કેટલાંક
મહત્ત્વનાં સમકક્ષ વિધાનો આપે છે.

\begin{thm}
\ollabel{thm:ce-equiv}
માનો કે $S$ પ્રાકૃતિક સંખ્યાઓનો ગણ છે. તો નીચેના વિધાનો
સમકક્ષ છે:
\begin{enumerate}
\item\ollabel{case:ce} $S$ સંગણકીય રીતે પરિગણનીય છે.
\item\ollabel{case:ran-pc} $S$ કોઈ \emph{આંશિક} સંગણનીય વિધેયનો મૂલ્યગણ છે.
\item\ollabel{case:ran-prim} $S$ ખાલી છે અથવા આદિમ પુનરાવર્તી વિધેયનો મૂલ્યગણ છે.
\item\ollabel{case:ce-domain} $S$ કોઈ આંશિક સંગણનીય વિધેયનો \emph{પ્રદેશ} છે.
\end{enumerate}
\end{thm}

\begin{explain}
પહેલાં ત્રણ ખંડો કહે છે કે દરેક ખાલી ન હોય એવા સંગણકીય રીતે
પરિગણનીય ગણની પરિગણના સંગણનીય વિધેય, આંશિક સંગણનીય વિધેય
અથવા આદિમ પુનરાવર્તી વિધેય વડે થઈ શકે છે. ચોથો ખંડ કહે છે કે
$S$ સંગણકીય રીતે પરિગણનીય હોય, તો કોઈ સૂચક~$e$ માટે
\[
S = \Setabs{x}{\cfind{e}(x) \fdefined}.
\]
બીજા શબ્દોમાં, $S$ એ એવા ઇનપુટોનો ગણ છે જેના માટે
$\cfind{e}$ની સંગણના થંભે છે. તેથી સંગણકીય રીતે પરિગણનીય
ગણોને ક્યારેક \emph{અર્ધનિર્ણેય} પણ કહે છે: સંખ્યા ગણમાં હોય
તો છેવટે ``હા'' મળે છે, પરંતુ ન હોય તો ``ના'' કદી મળતું નથી!{}
\end{explain}

\begin{proof}
દરેક આદિમ પુનરાવર્તી વિધેય સંગણનીય છે અને દરેક સંગણનીય વિધેય
આંશિક સંગણનીય પણ છે. તેથી \olref{case:ran-prim} પરથી
\olref{case:ce} અને \olref{case:ce} પરથી
\olref{case:ran-pc} મળે છે. (જો $S$ ખાલી હોય, તો $S$ ક્યાંય
વ્યાખ્યાયિત ન હોય એવા આંશિક સંગણનીય વિધેયનો મૂલ્યગણ છે.)
\olref{case:ran-pc} પરથી \olref{case:ran-prim} બતાવીએ,
તો પહેલાં ત્રણ ખંડોની સમકક્ષતા સાબિત થશે.

હવે માનો કે $S$ આંશિક સંગણનીય વિધેય $\cfind{e}$નો મૂલ્યગણ
છે. $S$ ખાલી હોય તો કામ પૂરું. અન્યથા $S$નું કોઈ તત્ત્વ $a$
લો. ક્લીનીના સામાન્ય સ્વરૂપના પ્રમેયથી
\[
\cfind{e}(x) = U(\umin{s}{T(e, x, s)}).
\]
ખાસ કરીને, $\cfind{e}(x) \fdefined$ અને તેનું મૂલ્ય $= y$
ત્યારે અને માત્ર ત્યારે જ્યારે કોઈ $s$ માટે $T(e, x, s)$
અને $U(s) = y$ હોય. $f(z)$ને નીચે પ્રમાણે વ્યાખ્યાયિત કરો:
\[
f(z) = \begin{cases}
  U((z)_1) & \text{જો $T(e, (z)_0, (z)_1)$} \\
  a        & \text{અન્યથા.}
\end{cases}
\]
$T$ અને $U$ આદિમ પુનરાવર્તી હોવાથી $f$ પણ આદિમ પુનરાવર્તી
છે. ટ્યુરિંગ મશીનની ભાષામાં, જો $z$ એવી
  જોડી $\tuple{(z)_0, (z)_1}$નો કોડ હોય કે $(z)_1$,
  મશીન~$M_e$ની ઇનપુટ $(z)_0$ પરની થંભતી સંગણના હોય,
  તો $f$ તે સંગણનાનું આઉટપુટ આપે છે; અન્યથા~$a$ આપે છે.

હવે $S$, $f$નો મૂલ્યગણ છે તે બતાવવાનું છે: દરેક પ્રાકૃતિક
સંખ્યા~$y$ માટે $y \in S$ ત્યારે અને માત્ર ત્યારે જ્યારે
તે $f$ના મૂલ્યગણમાં હોય. એક દિશામાં, માનો કે $y \in S$.
તો $y$, $\cfind{e}$ના મૂલ્યગણમાં છે; તેથી કોઈ $x$ અને~$s$
માટે $T(e,x,s)$ અને $U(s) = y$ છે. પરંતુ પછી
$y = f(\tuple{x,s})$. ઊલટી દિશામાં, માનો કે $y$, $f$ના
મૂલ્યગણમાં છે. તો $y = a$ છે અથવા કોઈ~$z$ માટે
$T(e,(z)_0,(z)_1)$ અને $U((z)_1) = y$ છે. પછીના કિસ્સામાં
$\cfind{e}((z)_0) \fdefined = y$ હોવાથી બંને રીતે $y$,
$S$માં છે.
\sourcecorrection{OLCT-004}{સ્થિર અંગ્રેજી મૂળ છેલ્લી
અનુમાન-કડીમાં $\cfind{e}(x) \fdefined = y$ લખે છે,
પરંતુ આ દિશામાં ઉપલબ્ધ સાક્ષી $z$ છે; તેનું પ્રથમ ઘટક
$(z)_0$ ઇનપુટ છે. તેથી અહીં $\cfind{e}((z)_0)$ રાખ્યું છે.}

(સંકેત $\cfind{e}(x) \fdefined = y$નો અર્થ
``$\cfind{e}(x)$ વ્યાખ્યાયિત છે અને તેનું મૂલ્ય $y$ છે.''
અહીં $\cfind{e}(x) = y$ પણ લખી શકાય; વધારાનું તીર
આંશિક વિધેય સાથે કામ કરીએ છીએ એ યાદ અપાવે છે.)

\olref{thm:ce-equiv}ની સાબિતી પૂરી કરવા
\olref{case:ce} અને~\olref{case:ce-domain} સમકક્ષ છે એ
બતાવવું પૂરતું છે. પહેલાં \olref{case:ce} પરથી
\olref{case:ce-domain} બતાવીએ. ખાલી ગણ ક્યાંય વ્યાખ્યાયિત
ન હોય એવા આંશિક સંગણનીય વિધેયનો પ્રદેશ છે; હવે બાકીના
કિસ્સામાં માનો કે $S$ સંગણનીય વિધેય~$f$નો મૂલ્યગણ છે, એટલે
\[
S = \Setabs{y}{\text{કોઈ $x$ માટે, } f(x) = y}.
\]
\sourcecorrection{OLCT-005}{સ્થિર અંગ્રેજી સાબિતી આ
દિશામાં ખાલી ગણને અલગથી લેતી નથી, જ્યારે કુલ વ્યાખ્યા
તેનો સમાવેશ કરે છે. અહીં ખાલી ગણ માટે ક્યાંય વ્યાખ્યાયિત
ન હોય એવા આંશિક વિધેયનો પ્રદેશ સ્પષ્ટ કર્યો છે.}
માનો કે
\[
g(y) = \umin{x}{(f(x) = y)}.
\]
ત્યારે $g$ આંશિક સંગણનીય છે અને $g(y)$ ત્યારે અને માત્ર
ત્યારે વ્યાખ્યાયિત છે જ્યારે કોઈ~$x$ માટે $f(x) = y$.
બીજા શબ્દોમાં, $g$નો પ્રદેશ $f$નો મૂલ્યગણ છે.
ટ્યુરિંગ મશીનની ભાષામાં: $S$નાં તત્ત્વોની
પરિગણના કરતું ટ્યુરિંગ મશીન~$F$ હોય, તો $G$ એવું મશીન
બનાવો કે જે $F$નાં આઉટપુટોમાં આપેલું તત્ત્વ શોધીને $S$નું
અર્ધનિર્ણય કરે; તત્ત્વ મળે તો થંભે અને ન મળે તો અનંતકાળ સુધી
શોધ ચાલુ રાખે.

છેવટે \olref{case:ce-domain} પરથી~\olref{case:ce}
બતાવવા, માનો કે $S$ આંશિક સંગણનીય વિધેય~$\cfind{e}$નો
પ્રદેશ છે, એટલે
\[
S = \Setabs{x}{\cfind{e}(x) \fdefined}.
\]
$S$ ખાલી હોય તો કામ પૂરું; અન્યથા $S$નું કોઈ તત્ત્વ $a$
લો. $f$ને નીચે પ્રમાણે વ્યાખ્યાયિત કરો:
\[
f(z) = \begin{cases}
(z)_0 & \text{જો $T(e,(z)_0,(z)_1)$} \\
a & \text{અન્યથા.}
\end{cases}
\]
તો, ઉપરની જેમ, સંખ્યા $x$, $f$ના મૂલ્યગણમાં ત્યારે અને
માત્ર ત્યારે છે જ્યારે $\cfind{e}(x) \fdefined$, એટલે જ્યારે
$x \in S$. ટ્યુરિંગ મશીનની ભાષામાં:
$S$નું અર્ધનિર્ણય કરતું મશીન $M_e$ હોય, તો ટ્યુરિંગ મશીનની
બધી શક્ય સંગણનાઓ તપાસીને, થંભતી સંગણનાઓને અનુરૂપ
ઇનપુટો આપતાં જવાથી $S$નાં તત્ત્વોની પરિગણના થાય છે.
\end{proof}

\olref{thm:ce-equiv}ના ખંડ~\olref{case:ce-domain} વડે
સંગણકીય રીતે પરિગણનીય ગણોની પરિગણના કરવાની સગવડભરી રીત
મળે છે: દરેક~$e$ માટે $W_e$, $\cfind{e}$નો પ્રદેશ હોય, એટલે
\[
W_e = \Setabs{x}{\cfind{e}(x) \fdefined}.
\]
તો $A$ કોઈ પણ સંગણકીય રીતે પરિગણનીય ગણ હોય, તો કોઈ~$e$
માટે $A = W_e$.

આગળનું પ્રમેય સંગણકીય રીતે પરિગણનીય ગણોનું વધુ એક લક્ષણ આપે છે.

\begin{thm}
\ollabel{thm:exists-char}
ગણ $S$ ત્યારે અને માત્ર ત્યારે સંગણકીય રીતે પરિગણનીય છે
જ્યારે એવો સંગણનીય સંબંધ $R(x,y)$ હોય કે
\[
S = \Setabs{ x }{ \lexists[y][R(x,y)] }.
\]
\end{thm}

\begin{proof}
એક દિશામાં, માનો કે $S$ સંગણકીય રીતે પરિગણનીય છે.
તો કોઈ $e$ માટે $S = W_e$. આ~$e$ માટે $S$ને આમ લખી શકાય:
\[
S = \Setabs{ x }{ \lexists[y][T(e, x, y)] }.
\]
ઊલટી દિશામાં, માનો કે $S = \Setabs{ x }{
  \lexists[y][R(x, y)] }$. $f$ને આ રીતે વ્યાખ્યાયિત કરો:
\[
f(x) \simeq \umin{y}{R(x, y)}.
\]
તો $f$ આંશિક સંગણનીય છે અને $S$, $f$નો પ્રદેશ છે.
\end{proof}
% END OLP-0239
\endgroup


\begingroup
\makeatletter\def\ol@sectioncs{section}\makeatother

% BEGIN OLP-0240
\olfileid{cmp}{thy}{ncp}
\olsection{અસંગણનીય ગણો અસ્તિત્વમાં છે}
\phantomsection\label{guunit:OLP-0240}



ઉપર જોયું કે દરેક સંગણનીય ગણ સંગણકીય રીતે પરિગણનીય છે. શું
વિરુદ્ધ દાવો પણ સાચો છે? આગળનું પ્રમેય બતાવે છે કે સામાન્ય રીતે
તે સાચો નથી.

\begin{thm}\ollabel{thm:K-0}
$K_0$ એ ગણ $\Setabs{\tuple{e, x}}{\cfind{e}(x) \fdefined}$ હોય.
તો $K_0$ સંગણકીય રીતે પરિગણનીય છે, પરંતુ સંગણનીય નથી.
\end{thm}

\begin{proof}
$K_0$ સંગણકીય રીતે પરિગણનીય છે તે જોવા, નોંધો કે તે નીચેના
વિધેય~$f$નો પ્રદેશ છે:
\[
f(z) = \umin{y}{(\len{z} = 2 \land T((z)_0, (z)_1, y))}.
\]
કારણ કે $\cfind{e}(x)$ વ્યાખ્યાયિત હોય તો
$f(\tuple{e, x})$ થંભતી સંગણનાનો અનુક્રમ શોધે છે; જો
$\cfind{e}(x)$ અનિર્ધારિત હોય, તો $f(\tuple{e, x})$ પણ
અનિર્ધારિત છે; અને $z$ જોડીનો કોડ પણ ન હોય, તો $f(z)$ પણ
અનિર્ધારિત છે.

$K_0$ સંગણનીય નથી એ હકીકત અટકવાની સમસ્યાની અનિર્ણેયતા,
\olref[hlt]{thm:halting-problem}, જ છે.
\end{proof}

ગણ $K_0$ એ એવી જોડીઓ $\tuple{e,x}$નો ગણ છે કે
$\cfind{e}(x) \fdefined$; એટલે $\tuple{e,x} \in K_0$ ત્યારે
અને માત્ર ત્યારે જ્યારે $\cfind{e}$ ઇનપુટ~$x$ પર વ્યાખ્યાયિત
હોય, અર્થાત્ અટકે. તેથી તેને ``અટકવાનો ગણ'' પણ કહે છે.
ગણ $K = \Setabs{e}{\cfind{e}(e) \fdefined}$ એ
``સ્વ-અટકવાનો ગણ'' છે. તેને ઘણી વાર પ્રમાણભૂત અનિર્ણેય ગણ
તરીકે વાપરાય છે.

\begin{thm}\ollabel{thm:K}
સ્વ-અટકવાનો ગણ $K = \Setabs{e}{\cfind{e}(e) \fdefined}$
સંગણકીય રીતે પરિગણનીય છે, પરંતુ નિર્ણેય નથી.
\end{thm}

\begin{proof}
  માનો કે $K$ નિર્ણેય છે, એટલે તેનું લક્ષણ વિધેય
  $\Char{K}$ સંગણનીય છે. માનો કે
  \[d(e) = \begin{cases}
  1 & \text{જો\/ $\Char{K}(e) = 0$}\\
  \fundefined & \text{અન્યથા.}
  \end{cases}
  \]
  $k$, $d$નો સૂચક હોય, એટલે $d \simeq \cfind{k}$.
  તો $d(k) \simeq \cfind{k}(k)$. પરંતુ $d$ની વ્યાખ્યા મુજબ
  $d(k) \fdefined$ ત્યારે અને માત્ર ત્યારે જ્યારે
  $\cfind{k}(k) \fundefined$; આ વિરોધાભાસ છે.

  $K$, $f(x) = \umin{y}{T(x,x,y)}$નો પ્રદેશ છે અને તેથી
  સંગણકીય રીતે પરિગણનીય છે.
\end{proof}
% END OLP-0240
\endgroup


\begingroup
\makeatletter\def\ol@sectioncs{section}\makeatother

% BEGIN OLP-0241
\olfileid{cmp}{thy}{clo}

\olsection[પરિગણનીય ગણોનો યોગગણ અને છેદગણ]{સંગણકીય રીતે પરિગણનીય
  ગણો યોગગણ અને છેદગણ હેઠળ બંધ છે}
\phantomsection\label{guunit:OLP-0241}


આગળનું પ્રમેય સંગણકીય રીતે પરિગણનીય ગણોના કેટલાક બંધત્વના
ગુણધર્મો આપે છે.

\begin{thm}
માનો કે $A$ અને $B$ સંગણકીય રીતે પરિગણનીય છે. તો $A \cap B$
અને $A \cup B$ પણ એવા જ છે.
\end{thm}

\begin{proof}
\olref[eqc]{thm:ce-equiv} વડે સંગણકીય રીતે પરિગણનીય ગણોના
વિવિધ લક્ષણો વાપરી શકીએ છીએ. વાત સ્પષ્ટ કરવા કેટલીક જુદી
સાબિતીઓ આપીશું.

પહેલી સાબિતી માટે માનો કે સંગણનીય વિધેય~$f$, $A$ની અને
સંગણનીય વિધેય~$g$, $B$ની પરિગણના કરે છે. માનો કે
\begin{align*}
h(x) & = \umin{y}{(f(y) = x \lor g(y) = x)} \text{ અને}\\
j(x) & = \umin{y}{(f((y)_0) = x \land g((y)_1) = x)}.
\end{align*}
ત્યારે $A \cup B$, $h$નો પ્રદેશ છે અને $A \cap B$,
$j$નો પ્રદેશ છે.

\begin{explain}
સંગણનાની દૃષ્ટિએ અહીં આવું બને છે: $A$ અને $B$ની પરિગણના
કરતી પ્રક્રિયાઓ હોય, તો કોઈ તત્ત્વ $x$, $A \cup B$માં છે
કે નહીં તેનો અર્ધનિર્ણય બંનેમાંથી એક પરિગણનામાં $x$ને
શોધીને કરી શકાય; અને $x$, $A \cap B$માં છે કે નહીં તેનો
અર્ધનિર્ણય બંને પરિગણનાઓમાં એકસાથે $x$ને શોધીને કરી શકાય.
\end{explain}

બીજી સાબિતી માટે ફરી માનો કે $f$, $A$ની અને $g$, $B$ની
પરિગણના કરે છે. માનો કે
\[
k(x) = \begin{cases}
f(x/2) & \text{જો $x$ સમ હોય} \\
g((x-1)/2) & \text{જો $x$ વિષમ હોય.}
\end{cases}
\]
ત્યારે $k$, $A \cup B$ની પરિગણના કરે છે: $k$ વારાફરતી
$f$ અને~$g$ની પરિગણનાઓમાંથી મૂલ્યો લે છે. $A \cap B$ની
પરિગણના કરવી વધુ મુશ્કેલ છે. $A \cap B$ ખાલી હોય, તો તે
સહેલાઈથી સંગણકીય રીતે પરિગણનીય છે. અન્યથા $A \cap B$નું
કોઈ તત્ત્વ $c$ લો અને $l$ને આ રીતે વ્યાખ્યાયિત કરો:
\[
l(x) = \begin{cases}
f((x)_0) & \text{જો $f((x)_0) = g((x)_1)$} \\
c & \text{અન્યથા.}
\end{cases}
\]
સંગણનાની દૃષ્ટિએ, $l$, $f$ અને $g$ની પરિગણનાઓમાં આવતાં
તત્ત્વોની જોડીઓ તપાસે છે અને મળતો દરેક મેળ આપે છે; અન્યથા
$c$ ફરી આપવાથી પરિગણનામાં કોઈ નવું તત્ત્વ ઉમેરાતું નથી.

છેલ્લી સાબિતી માટે માનો કે $A$, આંશિક વિધેય $m(x)$નો
\emph{પ્રદેશ} છે અને $B$, આંશિક વિધેય $n(x)$નો પ્રદેશ છે.
ત્યારે $A \cap B$, આંશિક વિધેય $m(x) + n(x)$નો પ્રદેશ છે.

\begin{explain}
સંગણનાની દૃષ્ટિએ, $A$ એ એવા મૂલ્યોનો ગણ છે જેના માટે
$m$ થંભે છે અને $B$ એ એવા મૂલ્યોનો ગણ છે જેના માટે
$n$ થંભે છે; $A \cap B$ એવા મૂલ્યોનો ગણ છે જેના માટે
બંને પ્રક્રિયાઓ થંભે છે.
\end{explain}

$A \cup B$ને થંભવાના મૂલ્યોના ગણ તરીકે દર્શાવવું વધુ
મુશ્કેલ છે, કારણ કે $m$ અને $n$ને સમાંતરે અનુકરવા પડે છે.
$d$, $m$નો અને $e$, $n$નો સૂચક હોય; બીજા શબ્દોમાં,
$m = \cfind{d}$ અને $n = \cfind{e}$. ત્યારે $A \cup B$
આ વિધેયનો પ્રદેશ છે:
\[
p(x) = \umin{y}{(T(d,x,y) \lor T(e,x,y))}.
\]
\begin{explain}
સંગણનાની દૃષ્ટિએ, ઇનપુટ $x$ પર $p$, $m$ અથવા $n$ની
થંભતી સંગણના શોધે છે અને બેમાંથી કોઈ એક મળે તો થંભે છે.
\end{explain}
\end{proof}
% END OLP-0241
\endgroup


\begingroup
\makeatletter\def\ol@sectioncs{section}\makeatother

% BEGIN OLP-0242
\olfileid{cmp}{thy}{cmp}
\olsection{સંગણકીય રીતે પરિગણનીય ગણો પૂરક હેઠળ બંધ નથી}
\phantomsection\label{guunit:OLP-0242}


માનો કે $A$ સંગણકીય રીતે પરિગણનીય છે. શું $A$નો પૂરક,
$\Complement{A} = \Nat \setminus A$, હંમેશાં સંગણકીય રીતે
પરિગણનીય હોય છે? આગળનું પ્રમેય અને પરિણામ બતાવે છે કે
જવાબ ``ના'' છે.

\begin{thm}
\ollabel{thm:ce-comp}
$A$ પ્રાકૃતિક સંખ્યાઓનો કોઈ પણ ગણ હોય. ત્યારે $A$ સંગણનીય
ત્યારે અને માત્ર ત્યારે છે જ્યારે $A$ અને $\Complement{A}$
બંને સંગણકીય રીતે પરિગણનીય હોય.
\end{thm}

\begin{proof}
એક દિશા સહેલી છે: $A$ સંગણનીય હોય તો $\Complement{A}$ પણ
સંગણનીય છે ($\Char{A} = 1 \tsub
\Char{\Complement{A}}$), તેથી બંને સંગણકીય રીતે પરિગણનીય છે.

ઊલટી દિશામાં, માનો કે $A$ અને~$\Complement{A}$ બંને
સંગણકીય રીતે પરિગણનીય છે. $A$, $\cfind{d}$નો પ્રદેશ
અને $\Complement{A}$, $\cfind{e}$નો પ્રદેશ હોય. $h$ને
આ રીતે વ્યાખ્યાયિત કરો:
\[
h(x) = \umin{s}{(T(d,x,s) \lor T(e,x,s))}.
\]
બીજા શબ્દોમાં, ઇનપુટ~$x$ પર $h$, $\cfind{d}$ અથવા
$\cfind{e}$ની થંભતી સંગણના શોધે છે. હવે $x \in A$ હોય
તો પહેલી સ્થિતિમાં અને $x \in \Complement{A}$ હોય તો
બીજી સ્થિતિમાં શોધ સફળ થાય છે. તેથી $h$ સંપૂર્ણ સંગણનીય
વિધેય છે. પરંતુ હવે દરેક~$x$ માટે $x \in A$ ત્યારે અને
માત્ર ત્યારે જ્યારે $T(d, x, h(x))$, એટલે વ્યાખ્યાયિત
થનાર વિધેય $\cfind{d}$ હોય. $T(d, x, h(x))$ સંગણનીય
સંબંધ હોવાથી $A$ સંગણનીય છે.
\sourcecorrection{OLCT-006}{સ્થિર અંગ્રેજી મૂળ આ અંતિમ
પરિક્ષણમાં $T(e, x, h(x))$ અને $\cfind{e}$ લખે છે,
જ્યારે ઉપર $d$નો પ્રદેશ $A$ અને $e$નો પ્રદેશ તેનો પૂરક
રાખ્યો છે. તેથી $A$ માટે અહીં બંને જગ્યાએ $d$ જોઈએ.}
\end{proof}

\begin{explain}
અનૌપચારિક સંગણનાની ભાષામાં વાત વધુ સહેલી છે: $A$નો નિર્ણય
કરવા ઇનપુટ $x$ પર $\cfind{d}$ અને $\cfind{e}$ની થંભતી
સંગણનાઓ શોધો. બેમાંથી એક તો જરૂર થંભશે. જો તે
$\cfind{d}$ હોય, તો $x$ એ~$A$માં છે; અન્યથા
$x$, $\Complement{A}$માં છે.
\sourcecorrection{OLCT-007}{સ્થિર અંગ્રેજી મૂળ આ
સમજૂતીમાં $\cfind{e}$ અને $\cfind{f}$ લખીને પહેલા
વિધેયને $A$ માટે ગણે છે. ઉપર નિર્ધારિત સૂચકો $d$ અને
$e$ છે; $A$નો પ્રદેશ $\cfind{d}$નો છે.}
\end{explain}

\begin{cor}
\ollabel{cor:comp-k}
$\Complement{K_0}$ સંગણકીય રીતે પરિગણનીય નથી.
\end{cor}

\begin{proof}
આપણે જાણીએ છીએ કે $K_0$ સંગણકીય રીતે પરિગણનીય છે,
પરંતુ સંગણનીય નથી. જો $\Complement{K_0}$ સંગણકીય રીતે
પરિગણનીય હોત, તો \olref{thm:ce-comp} મુજબ $K_0$
સંગણનીય હોત, જે \olref[ncp]{thm:K-0} સાથે વિરોધાભાસ છે.
\end{proof}
% END OLP-0242
\endgroup


\begingroup
\makeatletter\def\ol@sectioncs{section}\makeatother

% BEGIN OLP-0243
\olfileid{cmp}{thy}{red}
\olsection{ન્યૂનીકરણક્ષમતા}
\phantomsection\label{guunit:OLP-0243}


\begin{explain}
હવે આપણે જાણીએ છીએ કે ઓછામાં ઓછો એક ગણ $K_0$ એવો છે જે
સંગણકીય રીતે પરિગણનીય છે, પરંતુ સંગણનીય નથી. બીજા ગણો પણ
એવા જ હશે તે સ્પષ્ટ છે. ન્યૂનીકરણની પદ્ધતિ વારંવાર મૂળ
સિદ્ધાંતો તરફ પાછા જવું ન પડે તેમ બીજા ગણોમાં આ ગુણધર્મો
બતાવવાની શક્તિશાળી રીત આપે છે.

સામાન્ય રીતે, ગણ $A$નું ગણ~$B$માં ``ન્યૂનીકરણ'' એટલે
ઘટકો~$B$માં છે કે નહીં તે અંગેના જવાબોમાંથી
ઘટકો~$A$માં છે કે નહીં તે અંગેના જવાબો મેળવવાની
રીત. આપણે ``બહુ-એક ન્યૂનીકરણક્ષમતા'' પર ધ્યાન આપીશું,
જોકે જુદા ગુણધર્મો ધરાવતી ન્યૂનીકરણક્ષમતાની બીજી ઘણી
કલ્પનાઓ છે. સંગણકીય જટિલતાના અભ્યાસમાં પણ ન્યૂનીકરણની
કલ્પનાઓ કેન્દ્રસ્થાને છે; ત્યાં કાર્યક્ષમતાનો વિચાર પણ
કરવો પડે છે. દાખલા તરીકે, કોઈ ગણ NPમાં હોય અને દરેક NP
સમસ્યાને તેમાં એવી રીતથી ઘટાડી શકાય કે ન્યૂનીકરણ બહુ-એક
ન્યૂનીકરણ જેવું હોય, પરંતુ બહુપદી સમયમાં ગણી શકાય તેવી
વધારાની શરત પણ પૂરી કરે, તો તેને ``NP-પૂર્ણ'' કહે છે.

ન્યૂનીકરણની કલ્પના આપણે ગર્ભિત રીતે વાપરી જ છે. ગણ~$K$
આ રીતે વ્યાખ્યાયિત કરો:
\[
K = \Setabs{x}{\cfind{x}(x) \fdefined},
\]
એટલે $K = \Setabs{x}{x \in W_x}$. અટકવાની સમસ્યા
ઉકેલી શકાતી નથી તેની આપણી સાબિતી
(\olref[hlt]{thm:halting-problem}) સૌથી સીધે બતાવે છે કે
$K$ સંગણનીય નથી. યાદ કરો કે $K_0$ આ ગણ છે:
\[
K_0 = \Setabs{\tuple{e, x}}{\cfind{e}(x) \fdefined },
\]
એટલે $K_0 = \Setabs{\tuple{e,x}}{x \in W_e}$.
\sourcecorrection{OLCT-008}{સ્થિર અંગ્રેજી મૂળ બીજી
લખાવટમાં જોડીનો ક્રમ $\tuple{x,e}$ કરે છે. ઉપરની વ્યાખ્યા
અને $x \in W_e$ની શરત મુજબ ક્રમ $\tuple{e,x}$ જ છે.}
$K$ની અસંગણનીયતાની કોઈ પણ સાબિતીને $K_0$ની
અસંગણનીયતાની સાબિતી સુધી વિસ્તારી શકાય છે: જો $K_0$
સંગણનીય હોય, તો કોઈ ઘટક~$x$, $K$માં છે કે નહીં
તે જોડી $\tuple{x, x}$, $K_0$માં છે કે નહીં એ પૂછીને નક્કી
કરી શકીએ. $x$ને $\tuple{x, x}$માં મોકલતું વિધેય~$f$,
$K$ના $K_0$માં \emph{ન્યૂનીકરણ}નું ઉદાહરણ છે.
\end{explain}

\begin{defn}
માનો કે $A$ અને $B$ પ્રાકૃતિક સંખ્યાઓના ગણો છે.
સંગણનીય વિધેય~$f\colon \Nat \to \Nat$, $A$નું~$B$માં
\emph{બહુ-એક ન્યૂનીકરણ} ત્યારે અને માત્ર ત્યારે છે જ્યારે
દરેક પ્રાકૃતિક સંખ્યા~$x$ માટે
\[
x \in A \quad \text{ત્યારે અને માત્ર ત્યારે} \quad f(x) \in B.
\]
આવું ન્યૂનીકરણ $f$ હોય, તો $A$ને~$B$માં \emph{બહુ-એક
ન્યૂનીકરણક્ષમ} કહીએ; તેને $A \leq_m B$ લખીએ. જો $A$
ના $B$માં અને વિરુદ્ધ દિશામાં પણ બહુ-એક ન્યૂનીકરણ હોય,
તો $A$ અને $B$ને \emph{બહુ-એક સમકક્ષ} કહીએ; તેને
$A \equiv_m B$ લખીએ.
\end{defn}

ઉપરની વ્યાખ્યાનું વિધેય $f$ એક-એક પણ હોય, તો $A$ને~$B$માં
\emph{એક-એક ન્યૂનીકરણક્ષમ} કહેવાય છે. નીચેના મોટા ભાગના
ન્યૂનીકરણો આ વધુ મજબૂત શરત પૂરી કરે છે, પરંતુ આ હકીકતનો
અહીં ઉપયોગ કરીશું નહીં.

\begin{digress}
એક-એક ન્યૂનીકરણક્ષમતા ખરેખર બહુ-એક ન્યૂનીકરણક્ષમતા કરતાં
વધુ મજબૂત શરત છે; આ સાચું છે, પણ બિલકુલ સ્પષ્ટ નથી.
બીજા શબ્દોમાં, એવા અનંત ગણો $A$ અને~$B$ છે કે $A$ને~$B$માં
બહુ-એક રીતે ઘટાડી શકાય છે, પરંતુ $B$માં એક-એક રીતે નહીં.
\end{digress}
% END OLP-0243
\endgroup


\begingroup
\makeatletter\def\ol@sectioncs{section}\makeatother

% BEGIN OLP-0244
\olfileid{cmp}{thy}{ppr}
\olsection{ન્યૂનીકરણક્ષમતાના ગુણધર્મો}
\phantomsection\label{guunit:OLP-0244}


$A$ને~$B$માં ઘટાડી શકાય ત્યારે $A \leq_m B$ લખીએ છીએ. આ
સંકેત સૂચવે છે કે $A \leq_m B$ હોય તો $A$ એ~$B$ કરતાં
``વધુ કઠિન નથી'' અને $B$ એ~$A$ જેટલો અથવા વધુ કઠિન છે.
સંખ્યાઓ પરના સામાન્ય $\le$ની જેમ $\le_m$ ગણો (અથવા
નિર્ણય સમસ્યાઓ)ને તેમની જટિલતા મુજબ ક્રમિત કરે છે.
આગળના બે વિધાનો આ સમજને ટેકો આપે છે. પહેલું કહે છે કે
$\le_m$ સંક્રામક છે.

\begin{prop}
\ollabel{prop:trans-red}
$A \leq_m B$ અને $B \leq_m C$ હોય, તો $A \leq_m C$.
\end{prop}

\begin{proof}
$A$નું~$B$માં ન્યૂનીકરણ અને $B$નું $C$માં ન્યૂનીકરણ
સંયોજવાથી $A$નું~$C$માં ન્યૂનીકરણ મળે છે.
\end{proof}

\begin{prob}
જો $f$ અને~$g$ અનુક્રમે $A$નું~$B$માં અને $B$નું~$C$માં
બહુ-એક ન્યૂનીકરણ હોય, તો $\comp{f}{g}$ એ $A$નું~$C$માં
બહુ-એક ન્યૂનીકરણ છે એવું બતાવી \olref{prop:trans-red}
સાબિત કરો.
\end{prob}

\begin{prop}
\ollabel{prop:reduce}
$A$ અને $B$ કોઈ પણ ગણો હોય અને $A \leq_m B$ માનો.
\begin{enumerate}
\item $B$ સંગણકીય રીતે પરિગણનીય હોય, તો $A$ પણ છે.
\item $B$ સંગણનીય હોય, તો $A$ પણ છે.
\end{enumerate}
\end{prop}

\begin{proof}
$f$, $A$નું~$B$માં બહુ-એક ન્યૂનીકરણ હોય. પહેલા દાવા
માટે તપાસો કે $B$ આંશિક વિધેય~$g$નો પ્રદેશ હોય, તો
$A$ એ~$\comp{f}{g}$નો પ્રદેશ છે:
\begin{align*}
x \in A & \text{ ત્યારે અને માત્ર ત્યારે } f(x) \in B \\
& \text{ ત્યારે અને માત્ર ત્યારે }  g(f(x)) \fdefined.
\end{align*}

બીજા દાવા માટે યાદ કરો કે $B$ સંગણનીય હોય, તો $B$
અને~$\Complement{B}$ બંને સંગણકીય રીતે પરિગણનીય છે
(\olref[cmp]{thm:ce-comp}). $f$, $\Complement{A}$નું
$\Complement{B}$માં પણ બહુ-એક ન્યૂનીકરણ છે તે તપાસવું
મુશ્કેલ નથી. તેથી સાબિતીના પહેલા ભાગથી $A$
અને~$\Complement{A}$ બંને સંગણકીય રીતે પરિગણનીય છે. આથી \olref[cmp]{thm:ce-comp} મુજબ $A$ પણ
સંગણનીય છે. (વૈકલ્પિક રીતે, $\Char{A} =
\comp{f}{\Char{B}}$ છે તે તપાસી શકાય; તેથી $\Char{B}$
સંગણનીય હોય, તો~$\Char{A}$ પણ છે.)
\end{proof}

\begin{prob}
માનો કે $f$, $A$નું~$B$માં બહુ-એક ન્યૂનીકરણ છે. બતાવો કે
$f$, $\Complement{A}$નું $\Complement{B}$માં પણ બહુ-એક
ન્યૂનીકરણ છે.
\end{prob}

\begin{prob}
$f\colon A \to B$ બહુ-એક ન્યૂનીકરણ હોય, તો
$\Char{A} = \comp{f}{\Char{B}}$ છે એવું બતાવો.
\end{prob}

\begin{digress}
\emph{ટ્યુરિંગ ન્યૂનીકરણક્ષમતા} કહેવાતી ન્યૂનીકરણક્ષમતાની
વધુ સામાન્ય કલ્પના બીજા સંદર્ભોમાં, ખાસ કરીને
અનિર્ણેયતાનાં પરિણામો સાબિત કરવા, ઉપયોગી છે. નોંધો કે
\olref[cmp]{cor:comp-k} મુજબ $K_0$નો પૂરક~$K_0$માં
ઘટાડી શકાતો નથી, કારણ કે તે સંગણકીય રીતે પરિગણનીય નથી.
પરંતુ સહજ રીતે, $K_0$ અંગેના પ્રશ્નોના જવાબો ખબર હોય, તો
તેના પૂરક અંગેના પ્રશ્નોના જવાબો પણ ખબર પડે. ગણ $A$ને~$B$માં
ટ્યુરિંગ-ન્યૂનીકરણક્ષમ ત્યારે કહીએ જ્યારે~$B$ વિશે પ્રશ્નો
પૂછી શકતી સંગણનીય પ્રક્રિયા વડે~$A$ અંગેના પ્રશ્નોના જવાબો
નક્કી કરી શકાય. આ બહુ-એક ન્યૂનીકરણક્ષમતા કરતાં ઉદાર છે:
તેમાં (1)~$B$ વિશે માત્ર એક જ પ્રશ્ન પૂછી શકાય છે અને
(2)~``હા'' જવાબનો $A$ વિશેના પ્રશ્ન માટે ``હા''માં,
અને ``ના''નો ``ના''માં અનુવાદ થવો જોઈએ. તેમ છતાં,
$A$ને~$B$માં ટ્યુરિંગ રીતે ઘટાડી શકાય અને $B$ સંગણનીય
હોય, તો $A$ પણ સંગણનીય છે; જોકે, આપણે જોયું તેમ,
સંગણકીય પરિગણનીયતા માટે સમાન વિધાન સાચું નથી.

આપણે ચર્ચેલી ન્યૂનીકરણક્ષમતાની જુદી કલ્પનાઓ અને તેમના
ભેદો વિશે વિચારો. આ પ્રકરણમાં, જોકે, માત્ર બહુ-એક
ન્યૂનીકરણક્ષમતા જ લઈશું. ઉપરના ફકરામાં ચર્ચેલા બંને પ્રકારનાં
ન્યૂનીકરણને સંગણકીય જટિલતામાં પણ અનુરૂપ પ્રકારો છે, જેમાં
ટ્યુરિંગ મશીનો બહુપદી સમયમાં ચાલે એવી વધારાની શરત છે:
બહુ-એક ન્યૂનીકરણક્ષમતાના જટિલતા-રૂપને \emph{કાર્પ
ન્યૂનીકરણક્ષમતા} અને ટ્યુરિંગ ન્યૂનીકરણક્ષમતાના રૂપને
\emph{કુક ન્યૂનીકરણક્ષમતા} કહે છે.
\end{digress}
% END OLP-0244
\endgroup


\begingroup
\makeatletter\def\ol@sectioncs{section}\makeatother

% BEGIN OLP-0245
\olfileid{cmp}{thy}{cce}
\olsection{પૂર્ણ સંગણકીય રીતે પરિગણનીય ગણો}
\phantomsection\label{guunit:OLP-0245}


\begin{defn}
ગણ $A$ને (બહુ-એક ન્યૂનીકરણક્ષમતાની દૃષ્ટિએ)
\emph{પૂર્ણ સંગણકીય રીતે પરિગણનીય ગણ} ત્યારે કહીએ જ્યારે
\begin{enumerate}
\item $A$ સંગણકીય રીતે પરિગણનીય હોય અને
\item બીજા દરેક સંગણકીય રીતે પરિગણનીય ગણ $B$ માટે $B \leq_m A$ હોય.
\end{enumerate}
\end{defn}

બીજા શબ્દોમાં, પૂર્ણ સંગણકીય રીતે પરિગણનીય ગણો શક્ય
સૌથી ``કઠિન'' એવા પરિગણનીય ગણો છે. તેઓ \emph{કોઈ પણ}
સંગણકીય રીતે પરિગણનીય ગણ વિશેના પ્રશ્નોના જવાબ મેળવવા દે છે.

\begin{thm}
$K$, $K_0$ અને આગળ વ્યાખ્યાયિત થનારો $K_1$ બધા પૂર્ણ
સંગણકીય રીતે પરિગણનીય ગણો છે.
\end{thm}

\begin{proof}
$K_0$ પૂર્ણ છે તે જોવા કોઈ પણ સંગણકીય રીતે પરિગણનીય ગણ
$B$ લો. ત્યારે કોઈ સૂચક $e$ માટે
\[
B = W_e = \Setabs{x}{\cfind{e}(x) \fdefined}.
\]
$f$ એ $f(x) = \tuple{e, x}$ વિધેય હોય. ત્યારે દરેક પ્રાકૃતિક
સંખ્યા $x$ માટે $x \in B$ ત્યારે અને માત્ર ત્યારે જ્યારે
$f(x) \in K_0$. બીજા શબ્દોમાં, $f$, $B$નું~$K_0$માં
ન્યૂનીકરણ કરે છે.

$K_1$ પૂર્ણ છે તે જોવા નોંધો કે \olref[k1]{prop:k1}ની
સાબિતીમાં આપણે $K_0$નું તેમાં ન્યૂનીકરણ કર્યું છે. તેથી
\olref[ppr]{prop:trans-red} મુજબ કોઈ પણ સંગણકીય રીતે
પરિગણનીય ગણને~$K_1$માં ઘટાડી શકાય છે.

$K_0$નું $K$માં પણ લગભગ આ જ રીતે ન્યૂનીકરણ થાય છે:
ઉપરના કાર્યક્રમ-સૂચકના બાંધકામની જેમ, ઇનપુટને અવગણીને
આપેલી સંગણના ચલાવતો સૂચક બનાવો.
\sourcecorrection{OLCT-009}{સ્થિર અંગ્રેજી મૂળ છેલ્લે
$K$નું $K_0$માં ન્યૂનીકરણ કહે છે. તે વિરુદ્ધ દિશા છે અને
$K$ પૂર્ણ છે એ દાવાને સાબિત કરતી નથી. જરૂરી દિશા
$K_0$માંથી $K$ તરફ છે.}
\end{proof}

\begin{prob}
$K$નું $K_0$માં ન્યૂનીકરણ આપો.
\end{prob}

\begin{digress}
આમ, અત્યાર સુધી આપણે જોયેલા સંગણકીય રીતે પરિગણનીય ગણોના
બધાં ઉદાહરણો કાં તો સંગણનીય છે અથવા પૂર્ણ છે. આ વિચિત્ર
લાગવું જોઈએ!{} શું એવા સંગણકીય રીતે પરિગણનીય ગણો છે જે
સંગણનીય પણ નથી અને પૂર્ણ પણ નથી? જવાબ હા છે, પરંતુ
ફ્રિડબર્ગ અને મુચનિકે સ્વતંત્ર રીતે ૧૯૫૦ના દાયકાના મધ્યમાં
આ સ્થાપ્યું ત્યાં સુધી આ જાણીતું નહોતું.
\end{digress}
% END OLP-0245
\endgroup


\begingroup
\makeatletter\def\ol@sectioncs{section}\makeatother

% BEGIN OLP-0246
\olfileid{cmp}{thy}{k1}
\olsection{ન્યૂનીકરણક્ષમતાનું એક ઉદાહરણ}
\phantomsection\label{guunit:OLP-0246}


\olref[ppr]{prop:reduce}નો એક ઉપયોગ જોઈએ.

\begin{prop}
\ollabel{prop:k1}
માનો કે
\[
K_1 = \Setabs{e}{\cfind{e}(0) \fdefined}.
\]
તો $K_1$ સંગણકીય રીતે પરિગણનીય છે, પરંતુ સંગણનીય નથી.
\end{prop}

\begin{proof}
$K_1 = \Setabs{e}{\lexists[s][T(e,0,s)]}$ હોવાથી,
\olref[eqc]{thm:exists-char} મુજબ $K_1$ સંગણકીય રીતે
પરિગણનીય છે.

$K_1$ સંગણનીય નથી તે બતાવવા, $K_0$ને તેમાં ઘટાડી શકાય
છે એવું બતાવીશું.

\begin{explain}
આ થોડું મુશ્કેલ છે, કારણ કે $K_1$ વડે આપણે ફક્ત એક
ચોક્કસ ઇનપુટ $0$થી શરૂ થતી સંગણનાઓ વિશે જ પૂછી શકીએ છીએ.
માનો કે આ પ્રકારના પ્રશ્નોના જવાબ આપી શકે એવો તમારો હોશિયાર
મિત્ર છે (આવા મિત્રોને ``ઓરેકલ'' કહે છે). પછી માનો કે
કોઈ તમને $\tuple{e, x}$ એ $K_0$માં છે કે નહીં,
એટલે મશીન $e$ ઇનપુટ $x$ પર થંભે છે કે નહીં, એવો પ્રશ્ન
પૂછે છે. તમે બીજું મશીન $e_x$ બનાવી શકો છો: તે
\emph{કોઈ પણ} ઇનપુટ લે, તેને અવગણે અને તેના બદલે
ઇનપુટ $x$ પર~$e$ ચલાવે. તેથી મશીન $e$ ઇનપુટ $x$ પર
થંભે છે કે નહીં તે પ્રશ્ન, મશીન $e_x$ ઇનપુટ $0$ પર
(અથવા બીજા કોઈ પણ ઇનપુટ પર) થંભે છે કે નહીં તે પ્રશ્નને
સમકક્ષ છે. એટલે તમે મિત્રને પૂછો કે આ નવું મશીન $e_x$
ઇનપુટ $0$ પર થંભે છે કે નહીં; સુધારેલા પ્રશ્નનો જવાબ
મૂળ પ્રશ્નનો જવાબ આપે છે. આમ $K_0$નું~$K_1$માં ઇચ્છિત
ન્યૂનીકરણ મળે છે.
\end{explain}

સર્વવ્યાપી આંશિક સંગણનીય વિધેયનો ઉપયોગ કરીને, $f$ આ રીતે
વ્યાખ્યાયિત ત્રિ-આર્ગ્યુમેન્ટવાળું વિધેય હોય:
\[
f(x,y,z) \simeq \cfind{x}(y).
\]
નોંધો કે $f$ તેના ત્રીજા ઇનપુટને સંપૂર્ણપણે અવગણે છે.
એવો સૂચક $e$ લો કે $f = \cfind{e}[3]$; એટલે
\[
\cfind{e}[3](x,y,z) \simeq \cfind{x}(y).
\]
$s$-$m$-$n$ પ્રમેય મુજબ એવું વિધેય $s(e,x,y)$ છે કે
દરેક $z$ માટે
\begin{align*}
\cfind{s(e,x,y)}(z) & \simeq \cfind{e}[3](x,y,z) \\
& \simeq \cfind{x}(y).
\end{align*}

\begin{explain}
ઉપરની અનૌપચારિક દલીલની ભાષામાં $s(e,x,y)$ એ એવા મશીનનો
સૂચક છે જે કોઈ પણ ઇનપુટ $z$ને અવગણીને $\cfind{x}(y)$
ગણે છે.
\end{explain}

ખાસ કરીને,
\[
\cfind{s(e,x,y)}(0) \fdefined \quad \text{ત્યારે અને માત્ર ત્યારે} \quad
\cfind{x}(y) \fdefined.
\]
બીજા શબ્દોમાં, $\tuple{x, y} \in K_0$ ત્યારે અને માત્ર
ત્યારે જ્યારે $s(e,x,y) \in K_1$. તેથી નીચે પ્રમાણે
વ્યાખ્યાયિત વિધેય $g$,
\[
g(w) = s(e,(w)_0,(w)_1)
\]
$K_0$નું~$K_1$માં ન્યૂનીકરણ છે.
\end{proof}
% END OLP-0246
\endgroup


\begingroup
\makeatletter\def\ol@sectioncs{section}\makeatother

% BEGIN OLP-0247
\olfileid{cmp}{thy}{tot}
\olsection{સંપૂર્ણતા અનિર્ણેય છે}
\phantomsection\label{guunit:OLP-0247}


કોઈ વસ્તુ અસંગણનીય છે તે બતાવવા $s$-$m$-$n$ પ્રમેયનો
વધુ એક ઉપયોગ જોઈએ. $\fn{Tot}$ એ સંપૂર્ણ સંગણનીય વિધેયોના
સૂચકોનો ગણ હોય, એટલે
\[
\fn{Tot} = \Setabs{x}{\text{દરેક $y$ માટે, $\cfind{x}(y)\fdefined$}}.
\]

\begin{prop}
\ollabel{prop:total}
$\fn{Tot}$ સંગણનીય નથી.
\end{prop}

\begin{proof}
$\fn{Tot}$ સંગણનીય નથી તે જોવા $K$ને તેમાં ઘટાડી
શકાય છે એવું બતાવવું પૂરતું છે. $h(x,y)$ને આમ વ્યાખ્યાયિત કરો:
\[
h(x,y) \simeq
\begin{cases}
0 & \text{જો $x \in K$} \\
\fundefined & \text{અન્યથા}
\end{cases}
\]
નોંધો કે $h(x,y)$, $y$ પર જરાય આધાર રાખતું નથી.
$h$ આંશિક સંગણનીય છે તે સમજવું મુશ્કેલ નથી: ઇનપુટ $x, y$
પર~$h$ ગણવા પહેલાં ઇનપુટ~$x$ પર વિધેય~$\cfind{x}$નું
અનુકરણ કરીએ; તે સંગણના થંભે, તો $h(x,y)$, $0$ આપે અને
થંભે. આમ $h(x,y)$ એ $\Zero(\umin{s}{T(x,x,s)})$ જ છે,
જ્યાં $\Zero$ અચળ શૂન્ય વિધેય છે.

$s$-$m$-$n$ પ્રમેયથી એવું આદિમ પુનરાવર્તી વિધેય~$k(x)$
છે કે દરેક $x$ અને~$y$ માટે
\[
\cfind{k(x)}(y) =
\begin{cases}
0 & \text{જો $x \in K$} \\
\fundefined & \text{અન્યથા}
\end{cases}
\]
તેથી $x \in K$ હોય તો $\cfind{k(x)}$ સંપૂર્ણ છે અને
અન્યથા અનિર્ધારિત છે. આમ $k$, $K$નું~$\fn{Tot}$માં
ન્યૂનીકરણ છે.
\end{proof}

\begin{digress}
હકીકતમાં $\fn{Tot}$ સંગણકીય રીતે પરિગણનીય પણ નથી---
તેની જટિલતા ``અંકગણિતીય પદાનુક્રમ''માં વધુ ઊંચે આવે છે.
પરંતુ અહીં આ વધુ મજબૂત પરિણામની ચિંતા કરીશું નહીં.
\end{digress}
% END OLP-0247
\endgroup


\begingroup
\makeatletter\def\ol@sectioncs{section}\makeatother

% BEGIN OLP-0248
\olfileid{cmp}{thy}{rce}
\olsection{રાઇસનું પ્રમેય}
\phantomsection\label{guunit:OLP-0248}


વિચારો તો દેખાશે કે \olref[tot]{prop:total}ની સાબિતીમાં
$\fn{Tot}$ની કોઈ ખાસ વિશેષતાનો ઉપયોગ થયો નથી. $x$, $K$માં
હોય ત્યારે અને માત્ર ત્યારે $h(x,y)$ અચળ વિધેય $j(y) = 0$
ની જેમ વર્તે એવું આપણે ઘડ્યું; પરંતુ એ સ્થિતિમાં તેને
બીજા કોઈ પણ આંશિક સંગણનીય વિધેયની જેમ વર્તે એવું પણ
ઘડી શકતા. આ અવલોકનથી વધુ સામાન્ય પ્રમેય મળે છે:
આશરે કહીએ તો સંગણનીય વિધેયોનો કોઈ અતુચ્છ ગુણધર્મ
નિર્ણેય નથી.

યાદ રાખો કે $\cfind{0}$, $\cfind{1}$, $\cfind{2}$,~\dots
આંશિક સંગણનીય વિધેયોની આપણી પ્રમાણભૂત પરિગણના છે.

\begin{thm}[રાઇસનું પ્રમેય]
  $C$ આંશિક સંગણનીય વિધેયોનો કોઈ પણ ગણ હોય અને
  $A = \Setabs{n}{\cfind{n} \in C}$ રાખો. જો $A$ સંગણનીય
  હોય, તો $C$ ખાલી છે અથવા $C$ બધા આંશિક સંગણનીય
  વિધેયોનો ગણ છે.
\end{thm}

\emph{સૂચક ગણ} એટલે એવો ગણ $A$ કે જો સૂચકો $n$ અને $m$
એક જ વિધેય ``ગણતા'' હોય, તો બંને $n$ અને $m$ એ $A$માં
હોય, અથવા બંને ન હોય. પ્રમેયમાં આવેલો ગણ~$A$ આ
ગુણધર્મ ધરાવે છે તે જોવું મુશ્કેલ નથી. ઊલટું, જો $A$
સૂચક ગણ હોય અને $C$ આ સૂચકો વડે ગણાતા વિધેયોનો ગણ
હોય, તો $A = \Setabs{n}{\cfind{n} \in C}$.

\begin{explain}
આ પરિભાષામાં રાઇસના પ્રમેયનો અર્થ એવો થાય કે કોઈ અતુચ્છ
સૂચક ગણ નિર્ણેય નથી. પ્રમેય સમજવા \emph{પ્રોગ્રામો}
(માનો તમારી મનપસંદ પ્રોગ્રામિંગ ભાષામાં) અને તેઓ ગણતા
વિધેયો વચ્ચેનો ભેદ ધ્યાનમાં રાખો. પ્રોગ્રામો (સૂચકો)
વાક્યરચનાત્મક પદાર્થો છે; તેમની કેટલીક બાબતો નિશ્ચિતપણે
સંગણનીય છે: આ પ્રોગ્રામમાં ૧૫૦થી વધુ પ્રતીકો છે? તેમાં
૨૨થી વધુ લીટીઓ છે? તેમાં ``while'' વિધાન છે? ``print''
વિધાનના આર્ગ્યુમેન્ટ તરીકે ``hello world'' શબ્દમાળા આવે
છે? રાઇસનું પ્રમેય કહે છે કે પ્રોગ્રામના \emph{વર્તન}
વિશેનો કોઈ અતુચ્છ પ્રશ્ન સંગણનીય નથી. જેમ કે: પ્રોગ્રામ
ઇનપુટ $0$ પર થંભે છે? તે કદી થંભે છે? તે કદી સમ સંખ્યા
આઉટપુટ આપે છે?
\end{explain}

\begin{proof}[રાઇસના પ્રમેયની સાબિતી]
માનો કે $C$ ન તો ખાલી છે, ન બધા આંશિક સંગણનીય વિધેયોનો
ગણ; અને $A$, $C$માંના વિધેયોના સૂચકોનો ગણ છે. $A$
સંગણનીય હોય તો અટકવાની સમસ્યા ઉકેલી શકાય એવું બતાવીશું;
તેથી $A$ સંગણનીય નથી.

સામાન્યતા ગુમાવ્યા વિના, ક્યાંય વ્યાખ્યાયિત ન હોય એવું
વિધેય $f$, $C$માં નથી એમ માની શકીએ (અન્યથા નીચેની
દલીલમાં $C$ અને તેનો પૂરક અદલબદલ કરો). $g$, $C$માંનું
કોઈ પણ વિધેય હોય. વિચાર એવો છે કે $A$નો નિર્ણય કરી
શકીએ, તો $f$ ગણતા સૂચકો અને $g$ ગણતા સૂચકો વચ્ચે
ભેદ કરી શકીએ; પછી આ ક્ષમતાથી અટકવાની સમસ્યા ઉકેલી શકીએ.

રીત આ છે. સર્વવ્યાપી આંશિક સંગણનીય વિધેયોની મદદથી
આવું વિધેય વ્યાખ્યાયિત કરીએ:
\[
h(x,y) \simeq
\begin{cases}
\text{અનિર્ધારિત} & \text{જો $\cfind{x}(x) \fundefined$} \\
g(y) & \text{અન્યથા.}
\end{cases}
\]
$h$ ગણવા પહેલાં $\cfind{x}(x)$ ગણવાનો પ્રયત્ન કરીએ.
તે સંગણના થંભે, તો પછી~$g(y)$ ગણીએ; અને \emph{તે}
સંગણના પણ થંભે, તો આઉટપુટ પાછું આપીએ. વધુ ઔપચારિક
રીતે આમ લખી શકાય:
\[
h(x,y) \simeq \Proj{2}{0}(g(y),\fn{Un}(x,x)).
\]
અહીં $\Proj{2}{0}(z_0, z_1) = z_0$ એ $2$-સ્થાનીય પ્રક્ષેપ
વિધેય છે, જે $0$મા આર્ગ્યુમેન્ટ આપે છે અને
સંગણનીય છે.

તેથી $h$ આંશિક સંગણનીય વિધેયોના સંયોજનથી બને છે અને
જમણી બાજુ ત્યારે અને માત્ર ત્યારે વ્યાખ્યાયિત થઈને~$g(y)$
જેટલી થાય છે જ્યારે $\fn{Un}(x,x)$ અને $g(y)$ બંને
વ્યાખ્યાયિત હોય.

નોંધો કે કોઈ નિશ્ચિત~$x$ માટે $\cfind{x}(x)$ અનિર્ધારિત
હોય, તો દરેક~$y$ માટે $h(x,y)$ અનિર્ધારિત છે; અને
$\cfind{x}(x)$ વ્યાખ્યાયિત હોય, તો $h(x,y) \simeq g(y)$.
આમ~$x$નું કોઈ પણ નિશ્ચિત મૂલ્ય લેતાં $h(x,y)$ કાં તો
$f$ની જેમ અથવા $g$ની જેમ વર્તે છે. એટલે
$\cfind{x}(x)$ વ્યાખ્યાયિત છે કે નહીં તે નક્કી કરવું
આ બે સ્થિતિમાંથી કઈ છે તે નક્કી કરવા સમાન છે. પરંતુ
તે $h_x(y) \simeq h(x,y)$ એ~$C$માં છે કે નહીં તે
નક્કી કરવા સમાન છે; $A$ સંગણનીય હોય તો તે કરી શકીએ.

વધુ ઔપચારિક રીતે, $h$ આંશિક સંગણનીય હોવાથી કોઈ
સૂચક~$e$ માટે તે વિધેય $\cfind{e}[2]$ છે.
\sourcecorrection{OLCT-010}{સ્થિર અંગ્રેજી મૂળ અહીં
દ્વિ-આર્ગ્યુમેન્ટવાળા $h(x,y)$ માટે એક-આર્ગ્યુમેન્ટવાળું
$\cfind{e}$ લખે છે. એ જ સૂચકનું દ્વિ-આર્ગ્યુમેન્ટવાળું
સ્વરૂપ $\cfind{e}[2]$ અહીં વપરાયું છે.}
$s$-$m$-$n$ પ્રમેયથી એવું આદિમ પુનરાવર્તી વિધેય~$s$
છે કે દરેક~$x$ માટે $\cfind{s(e,x)}(y) = h_x(y)$.
હવે દરેક $x$ માટે $\cfind{x}(x) \fdefined$ હોય, તો
$\cfind{s(e,x)}$ એ~$g$ જેવું જ વિધેય છે અને તેથી
$s(e,x)$ એ $A$માં છે. બીજી તરફ, $\cfind{x}(x)
\uparrow$ હોય, તો $\cfind{s(e,x)}$ એ~$f$ જેવું જ
વિધેય છે અને તેથી $s(e,x)$ એ~$A$માં નથી. બીજા
શબ્દોમાં, દરેક~$x$ માટે $x \in K$ ત્યારે અને માત્ર
ત્યારે જ્યારે $s(e,x) \in A$. $A$ સંગણનીય હોત,
તો $K$ પણ હોત; આ વિરોધાભાસ છે. તેથી $A$
સંગણનીય નથી.
\end{proof}

રાઇસનું પ્રમેય ખૂબ શક્તિશાળી છે. આગળનું તાત્કાલિક પરિણામ
તેના કેટલાક ઉપયોગો બતાવે છે.
\begin{cor}
આગળના ગણો અનિર્ણેય છે.
\begin{enumerate}
\item $\Setabs{x}{\text{$17$ એ $\cfind{x}$ના મૂલ્યગણમાં છે}}$
\item $\Setabs{x}{\text{$\cfind{x}$ અચળ છે}}$
\item $\Setabs{x}{\text{$\cfind{x}$ સંપૂર્ણ છે}}$
\item $\Setabs{x}{\text{જ્યારે $y < y'$, ત્યારે $\cfind{x}(y) \fdefined$,
    અને $\cfind{x}(y') \fdefined$ હોય તો $\cfind{x}(y) < \cfind{x}(y')$}}$
\end{enumerate}
\end{cor}

\begin{proof}
  આ બધા અતુચ્છ સૂચક ગણો છે.
\end{proof}
% END OLP-0248
\endgroup


\begingroup
\makeatletter\def\ol@sectioncs{section}\makeatother

% BEGIN OLP-0249
\olfileid{cmp}{thy}{fix}
\olsection{નિશ્ચિત-બિંદુનું પ્રમેય}
\phantomsection\label{guunit:OLP-0249}


અટકવાની સમસ્યા ફરી વિચારીએ. અસ્થાયી સંકેત તરીકે
$\tuple{x, y}$ માટે $\gn{\cfind{x}(y)}$ લખીએ; તેને
$\cfind{x}(y)$ના મૂલ્યનું ``નામ'' માનો. આ સંકેત વડે
અટકવાની સમસ્યા અનિર્ણેય છે તેની આપણી એક સાબિતી
બીજી રીતે કહી શકાય.

પ્રશ્ન: શું એવું સંગણનીય વિધેય $h$ છે કે દરેક $x$ અને
$y$ માટે
\[
h(\gn{\cfind{x}(y)}) =
\begin{cases}
1 & \text{જો $\cfind{x}(y) \fdefined$} \\
0 & \text{અન્યથા.}
\end{cases}
\]

જવાબ: ના. અન્યથા આંશિક વિધેય
\[
g(x) \simeq
\begin{cases}
0 & \text{જો $h(\gn{\cfind{x}(x)}) = 0$} \\
\text{અનિર્ધારિત} & \text{અન્યથા}
\end{cases}
\]
સંગણનીય હોત અને તેથી તેનો કોઈ સૂચક~$e$ હોત. પરંતુ પછી
\[
\cfind{e}(e) \simeq
\begin{cases}
0 & \text{જો $h(\gn{\cfind{e}(e)}) = 0$} \\
\text{અનિર્ધારિત} & \text{અન્યથા,}
\end{cases}
\]
એટલે $\cfind{e}(e)$ ત્યારે અને માત્ર ત્યારે વ્યાખ્યાયિત
થાય જ્યારે તે અનિર્ધારિત હોય; આ વિરોધાભાસ છે.

હવે $\cfind{e}$વાળું સમીકરણ જુઓ. તેમાં એક અર્થમાં
સ્વ-સંદર્ભ છે: $\cfind{e}(e)$નું મૂલ્ય એક ચોક્કસ રીતે
$\gn{\cfind{e}(e)}$ પર આધાર રાખે એવું આપણે ગોઠવ્યું છે.
નિશ્ચિત-બિંદુનું પ્રમેય કહે છે કે સામાન્ય રીતે પણ આ કરી
શકાય છે---માત્ર વિરોધાભાસ સાબિત કરવા માટે જ નહીં.

\olref{lem:fixed-equiv} નિશ્ચિત-બિંદુનું પ્રમેય બે સમકક્ષ
રીતે રજૂ કરે છે. તર્કની દૃષ્ટિએ બંને વિધાનો સાચાં હોવાથી
તેમની સમકક્ષતા મળે છે; પરંતુ ખરેખર કહેવાનો અર્થ એ છે કે
દરેક વિધાન બીજામાંથી સીધું મળે છે, તેથી તેમને એક જ પ્રમેયનાં
વૈકલ્પિક વિધાનો માની શકાય.

\begin{lem}
\ollabel{lem:fixed-equiv}
આગળનાં વિધાનો સમકક્ષ છે:
\begin{enumerate}
\item દરેક આંશિક સંગણનીય વિધેય $g(x,y)$ માટે એવો
  સૂચક~$e$ છે કે દરેક~$y$ માટે
\[
\cfind{e}(y) \simeq g(e,y).
\]
\item દરેક સંગણનીય વિધેય~$f(x)$ માટે એવો સૂચક~$e$
  છે કે દરેક~$y$ માટે
\[
\cfind{e}(y) \simeq \cfind{f(e)}(y).
\]
\end{enumerate}
\end{lem}

\begin{proof}
$(1) \Rightarrow (2)$: $f$ આપેલું હોય, તો $g$ને
$g(x,y) \simeq \fn{Un}(f(x),y)$ વડે વ્યાખ્યાયિત કરો.
(1) વડે એવો સૂચક~$e$ મેળવો કે દરેક~$y$ માટે
\begin{align*}
\cfind{e}(y) & = \fn{Un}(f(e),y) \\
& = \cfind{f(e)}(y).
\end{align*}

$(2) \Rightarrow (1)$: $g$ આપેલું હોય, તો $s$-$m$-$n$
પ્રમેયથી એવું $f$ મેળવો કે દરેક $x$ અને~$y$ માટે
$\cfind{f(x)}(y) \simeq g(x,y)$. (2) વડે એવો સૂચક~$e$
મેળવો કે
\begin{align*}
\cfind{e}(y) & = \cfind{f(e)}(y) \\
& = g(e,y).
\end{align*}
આથી સાબિતી પૂરી થાય છે.
\end{proof}

\begin{explain}
વિધાન (1) સાચું છે (અને તેથી (2) પણ) એ બતાવતાં પહેલાં
તે કેટલું નવાઈભર્યું છે તે વિચારો. $e$ને કમ્પ્યુટર પ્રોગ્રામ
માનો; વિધાન (1) કહે છે કે કોઈ પણ આંશિક સંગણનીય
$g(x,y)$ માટે $g_e(y) \simeq g(e,y)$ ગણતો કમ્પ્યુટર
પ્રોગ્રામ $e$ શોધી શકાય છે. બીજા શબ્દોમાં, પોતાને જ
સંદર્ભિત કરતું વિધેય ગણતો પ્રોગ્રામ શોધી શકાય છે.
\end{explain}

\begin{thm}
\olref{lem:fixed-equiv}નાં બંને વિધાનો સાચાં છે. ખાસ
કરીને, દરેક આંશિક સંગણનીય વિધેય $g(x,y)$ માટે એવો
સૂચક~$e$ છે કે દરેક~$y$ માટે
\[
\cfind{e}(y) \simeq g(e,y).
\]
\end{thm}

\begin{proof}
જરૂરી ઘટકો ઉપરની અટકવાની સમસ્યાની ચર્ચામાં ગર્ભિત છે.
$\fn{diag}(x)$ એવું સંગણનીય વિધેય હોય કે દરેક $x$ માટે
તે વિધેય $f_x(y) \simeq \cfind{x}[2](x,y)$નો સૂચક
પાછો આપે, એટલે
\[
\cfind{\fn{diag}(x)}(y) \simeq \cfind{x}[2](x,y).
\]
$\fn{diag}$ને એવી પ્રક્રિયા માનો કે જે દ્વિ-આર્ગ્યુમેન્ટવાળા
વિધેયના પ્રોગ્રામમાંથી એક-આર્ગ્યુમેન્ટવાળા વિધેયનો પ્રોગ્રામ
બનાવે છે, જ્યાં મૂળ પ્રોગ્રામ પોતે પ્રથમ આર્ગ્યુમેન્ટ તરીકે
નિશ્ચિત કરાયો છે. $\fn{diag}$ને ઔપચારિક રીતે આ રીતે
વ્યાખ્યાયિત કરી શકાય: પહેલાં $s$ને આમ વ્યાખ્યાયિત કરો:
\[
s(x,y) \simeq \fn{Un}^2(x,x,y),
\]
જ્યાં $\fn{Un}^2$ દ્વિ-આર્ગ્યુમેન્ટવાળાં આંશિક સંગણનીય
વિધેયો માટે સર્વવ્યાપી ત્રિ-આર્ગ્યુમેન્ટવાળું વિધેય છે.
પછી $s$-$m$-$n$ પ્રમેયથી એવું આદિમ પુનરાવર્તી વિધેય
$\fn{diag}$ મળે છે કે
\[
\cfind{\fn{diag}(x)}(y) \simeq s(x,y).
\]
\sourcecorrection{OLCT-011}{સ્થિર અંગ્રેજી મૂળ આ
બાંધકામમાં બે જગ્યાએ $\cfind{x}(x,y)$ અને નીચેની
ગણતરીમાં $\cfind{\gn{l}}(\gn{l},y)$ લખે છે. બંને
દ્વિ-આર્ગ્યુમેન્ટવાળાં વિધેયો છે; અહીં તે મુજબ
$\cfind{x}[2]$ અને $\cfind{\gn{l}}[2]$ રાખ્યાં છે.}

હવે વિધેય $l$ને આમ વ્યાખ્યાયિત કરો:
\[
l(x,y) \simeq g(\fn{diag}(x),y).
\]
અને $\gn{l}$ એ $l$નો સૂચક હોય. છેલ્લે
$e = \fn{diag}(\gn{l})$ રાખો. ત્યારે દરેક $y$ માટે
\begin{align*}
\cfind{e}(y) & \simeq \cfind{\fn{diag}(\gn{l})}(y) \\
& \simeq \cfind{\gn{l}}[2](\gn{l}, y) \\
& \simeq l(\gn{l}, y) \\
& \simeq g(\fn{diag}(\gn{l}),y) \\
& \simeq g(e, y),
\end{align*}
જે જોઈતું હતું.
\end{proof}

\begin{explain}
અહીં શું બને છે? માનો કે પોતાને જ છાપતો કમ્પ્યુટર પ્રોગ્રામ
લખવાનું કામ તમને મળ્યું છે. વધુમાં માનો કે તમે એવી
પ્રોગ્રામિંગ ભાષામાં કામ કરો છો જેમાં શબ્દમાળાઓનાં વિપુલ
અને વિચિત્ર વિધેયોનું ગ્રંથાલય છે. ખાસ કરીને, તેમાં
$\fn{diag}$ નામનું વિધેય હોય જે આ રીતે કામ કરે:
ઇનપુટ શબ્દમાળા~$s$ આપવામાં આવે, તો $\fn{diag}$ એ~$s$માં
આવતાં `x' પ્રતીકનાં દરેક સ્થાનને શોધીને તેના સ્થાને
મૂળ શબ્દમાળાનું અવતરણ-ચિહ્નવાળું
રૂપ મૂકે છે. દાખલા તરીકે, આ શબ્દમાળા આપો:
\begin{quote}
\begin{verbatim}
hello x world
\end{verbatim}
\end{quote}
તો તે પાછું આપે:
\begin{quote}
\begin{verbatim}
hello 'hello x world' world
\end{verbatim}
\end{quote}
આ સ્થિતિમાં ઇચ્છિત પ્રોગ્રામ લખવો સહેલો છે; તપાસો કે
\begin{quote}
\begin{verbatim}
print(diag('print(diag(x))'))
\end{verbatim}
\end{quote}
કામ કરે છે. C++ અને Java જેવી સામાન્ય ભાષાઓમાં પણ
આ જ વિચાર (વધુ જટિલ અમલ સાથે) કામ કરે છે.

નિશ્ચિત-બિંદુના પ્રમેયની સાબિતીથી હવે માત્ર થોડાં પગલાં
દૂર છીએ. માનો કે છાપવાના વિધેયનું એક રૂપ
$\fn{print}(x,y)$ શબ્દમાળા $x$ અને સંખ્યાત્મક આર્ગ્યુમેન્ટ
$y$ લે છે અને શબ્દમાળા $x$ને $y$ વાર છાપે છે. ત્યારે આ
``પ્રોગ્રામ''
\begin{quote}
\begin{verbatim}
getinput(y);
print(diag('getinput(y); print(diag(x), y)'), y)
\end{verbatim}
\end{quote}
ઇનપુટ $y$ પર પોતાને $y$ વાર છાપે છે.
$\fn{getinput}$---$\fn{print}$---$\fn{diag}$ના માળખાને
કોઈ પણ વિધેય $g(x,y)$ વડે બદલીને
\begin{quote}
\begin{verbatim}
g(diag('g(diag(x), y)'), y)
\end{verbatim}
\end{quote}
મળે છે; આ પ્રોગ્રામ ઇનપુટ $y$ પર, પોતે અને $y$
આર્ગ્યુમેન્ટો લઈને $g$ ચલાવે છે. ``અવતરણમાં મૂકવા''ને
``સૂચક વાપરવા'' તરીકે સમજીએ, તો ઉપરની સાબિતી મળે છે.

હમણાં માટે, આ સાબિતી ઔપચારિક યુક્તિ અથવા જાદુ જેવી
લાગે તો ચાલે. પરંતુ ઉપરની દલીલની વિગતો તમે ફરીથી
ઘડી શકો એવું હોવું જોઈએ. અપૂર્ણતાનાં પ્રમેયો (અને
સંબંધિત ``નિશ્ચિત-બિંદુનું પ્રમેય'') સાબિત કરીએ ત્યારે
આ રીત કેમ કામ કરે છે તે સમજવાની બીજી રીતો પણ જોઈશું.
\end{explain}

\begin{digress}
આ જ વિચારથી ``નિશ્ચિત-બિંદુ'' સંયોજક મેળવી શકાય છે.
માનો કે તમારી પાસે લૅમ્બ્ડા પદ $g$ છે અને એવું બીજું
પદ $k$ જોઈએ છે કે $k$, $gk$ સાથે $\beta$-સમકક્ષ હોય.
આપણી સંકેત-પદ્ધતિથી પદો વ્યાખ્યાયિત કરો:
\[
\fn{diag}(x) = xx
\]
અને
\[
l(x) = g(\fn{diag}(x))
\]
બીજા શબ્દોમાં, $l$ એ પદ $\lambd[x][g(xx)]$ છે. $k$
એ પદ $ll$ હોય. ત્યારે
\begin{align*}
k & = (\lambd[x][g(xx)])(\lambd[x][g(xx)]) \\
& \red  g((\lambd[x][g(xx)])(\lambd[x][g(xx)])) \\
& = gk.
\end{align*}
જો
\[
Y = \lambd[g][((\lambd[x][g(xx)])(\lambd[x][g(xx)]))]
\]
રાખીએ, તો $Yg$ અને $g(Yg)$ બંને એક સામાન્ય પદ સુધી
ઘટે છે; એટલે $Yg \equiv_\beta
g(Yg)$. તેને ``કરીનો સંયોજક'' કહે છે. તેના બદલે જો
\[
Y = (\lambd[xg][g(xxg)])(\lambd[xg][g(xxg)])
\]
રાખીએ, તો ખરેખર $Yg$, $g(Yg)$ સુધી ઘટે છે, જે વધુ
મજબૂત વિધાન છે. $Y$ના આ બીજા રૂપને ``ટ્યુરિંગનો
સંયોજક'' કહે છે.
\end{digress}
% END OLP-0249
\endgroup


\begingroup
\makeatletter\def\ol@sectioncs{section}\makeatother

% BEGIN OLP-0250
\olfileid{cmp}{thy}{apf}
\olsection{નિશ્ચિત-બિંદુના પ્રમેયનો ઉપયોગ}
\phantomsection\label{guunit:OLP-0250}


નિશ્ચિત-બિંદુનું પ્રમેય આંશિક સંગણનીય વિધેયોને તેમના
સૂચકોના આધારે વ્યાખ્યાયિત કરવાની છૂટ આપે છે. દાખલા
તરીકે, એવો સૂચક $e$ મળે છે કે દરેક $y$ માટે
\[
\cfind{e}(y) = e + y.
\]
બીજા ઉદાહરણ તરીકે, નિશ્ચિત-બિંદુના પ્રમેયની સાબિતીથી
પોતાને જ છાપતો Java અથવા C++ પ્રોગ્રામ ઘડી શકાય છે.

યાદ કરો કે દરેક $e$ માટે $W_e$ને $\cfind{e}$નો પ્રદેશ
રાખીએ, તો અનુક્રમ $W_0$, $W_1$, $W_2$,~\dots સંગણકીય
રીતે પરિગણનીય ગણોની પરિગણના કરે છે. આ ગણોમાંના કેટલાક
સંગણનીય છે. એવો પ્રશ્ન થાય કે શું એવી કલનવિધિ છે જે
મૂલ્ય $x$ ઇનપુટ તરીકે લે અને, જો $W_x$ સંગણનીય હોય,
તો તેના લક્ષણ વિધેયનો સૂચક પાછો આપે? જવાબ ``ના'' છે;
આવી કલનવિધિ નથી:

\begin{thm}
આગળનો ગુણધર્મ ધરાવતું કોઈ આંશિક સંગણનીય વિધેય $f$
નથી: જ્યારે $W_e$ સંગણનીય હોય, ત્યારે $f(e)$
વ્યાખ્યાયિત હોય અને $\cfind{f(e)}$ તેનું લક્ષણ વિધેય હોય.
\end{thm}

\begin{proof}
$f$ કોઈ પણ આંશિક સંગણનીય વિધેય હોય; આપણે એવો $e$
બનાવીશું કે $W_e$ સંગણનીય હોય, પરંતુ $\cfind{f(e)}$
તેનું લક્ષણ વિધેય ન હોય. નિશ્ચિત-બિંદુના પ્રમેયથી
એવો સૂચક $e$ મળે છે કે
\[
\cfind{e}(y) \simeq
\begin{cases}
0 & \text{જો $y=0$ અને $\cfind{f(e)}(0) \fdefined = 0$} \\
\text{અનિર્ધારિત} & \text{અન્યથા.}
\end{cases}
\]
એટલે $e$ નીચેના વિધેય પર નિશ્ચિત-બિંદુનું પ્રમેય
લાગુ કરીને મળે છે:
\[
g(x,y) \simeq
\begin{cases}
0 & \text{જો $y=0$ અને $\cfind{f(x)}(0) \fdefined = 0$} \\
\text{અનિર્ધારિત} & \text{અન્યથા.}
\end{cases}
\]
અનૌપચારિક રીતે $g$ આંશિક સંગણનીય છે: ઇનપુટ $x$
અને $y$ પર કલનવિધિ પહેલાં તપાસે છે કે $y$ એ~$0$
છે કે નહીં. હોય, તો તે $f(x)$ ગણે છે અને પછી
સર્વવ્યાપી મશીનથી $\cfind{f(x)}(0)$ ગણે છે. આ
છેલ્લી સંગણના થંભીને~$0$ આપે, તો કલનવિધિ~$0$ આપે;
અન્યથા તે થંભતી નથી.
\sourcecorrection{OLCT-012}{સ્થિર અંગ્રેજી મૂળ સાબિતી
શરૂ કરતાં $f$ને સંપૂર્ણ સંગણનીય માને છે, જ્યારે પ્રમેય
બધાં આંશિક સંગણનીય વિધેયોને નકારે છે. ઉપરના $g$ની
ગણતરી $f(x)$ અનિર્ધારિત રહે ત્યારે પણ આંશિક સંગણનીય
રહે છે; તેથી સાબિતી આ વ્યાપક દાવાને પણ આવરે છે.}

હવે નોંધો કે $\cfind{f(e)}(0)$ વ્યાખ્યાયિત થઈને $0$
જેટલું હોય, તો $\cfind{e}(y)$ ત્યારે અને માત્ર ત્યારે
વ્યાખ્યાયિત છે જ્યારે $y$, $0$ હોય; તેથી $W_e =
\{ 0 \}$. જો $\cfind{f(e)}(0)$ અનિર્ધારિત હોય અથવા
વ્યાખ્યાયિત હોય પરંતુ $0$ ન હોય, તો $W_e = \emptyset$.
બંને રીતે $\cfind{f(e)}$, $W_e$નું લક્ષણ વિધેય નથી:
ઇનપુટ $0$ પર તે ખોટો જવાબ આપે છે. જો સૂચક આપતું
વિધેય જ આ ઇનપુટ પર અનિર્ધારિત હોય, તો માગેલી શરત
પહેલેથી જ ભંગાય છે.
\end{proof}
% END OLP-0250
\endgroup


\begingroup
\makeatletter\def\ol@sectioncs{section}\makeatother

% BEGIN OLP-0251
\olfileid{cmp}{thy}{slf}
\olsection{સ્વ-સંદર્ભ વડે વિધેયોની વ્યાખ્યા}
\phantomsection\label{guunit:OLP-0251}


વિધેયોને તેમના પોતાના જ આધારે વ્યાખ્યાયિત કરી શકવાની
ક્ષમતા સામાન્ય રીતે ઉપયોગી છે. દાખલા તરીકે, સંગણનીય
વિધેયો $k$, $l$ અને~$m$ આપેલાં હોય, તો નિશ્ચિત-બિંદુનું
ઉપપ્રમેય કહે છે કે દરેક~$y$ માટે નીચેનું સમીકરણ સંતોષતું
આંશિક સંગણનીય વિધેય~$f$ છે:
\[
f(y) \simeq
\begin{cases}
k(y) & \text{જો $l(y) = 0$} \\
f(m(y)) & \text{અન્યથા.}
\end{cases}
\]
વધુ ચોક્કસ રીતે, $f$ મેળવવા પહેલાં
\[
g(x,y) \simeq
\begin{cases}
k(y) & \text{જો $l(y) = 0$} \\
\cfind{x}(m(y)) & \text{અન્યથા}
\end{cases}
\]
રાખીએ અને પછી નિશ્ચિત-બિંદુના ઉપપ્રમેયથી એવો સૂચક~$e$
મેળવીએ કે $\cfind{e}(y) = g(e,y)$.

ચોક્કસ ઉદાહરણ તરીકે, ``મહત્તમ સમાપવર્તક'' વિધેય
$\fn{gcd}(u,v)$ આ રીતે વ્યાખ્યાયિત કરી શકાય:
\[
\fn{gcd}(u,v) \simeq
\begin{cases}
v & \text{જો $u = 0$} \\
\fn{gcd}(\fn{mod}(v, u), u) & \text{અન્યથા}
\end{cases}
\]
અહીં $\fn{mod}(v, u)$, $v$ને~$u$ વડે ભાગતાં મળતી
શેષ સંખ્યા દર્શાવે છે. નિશ્ચિત-બિંદુના ઉપપ્રમેયથી
$\fn{gcd}$ આંશિક સંગણનીય છે. (હકીકતમાં $y$ જોડી
$\tuple{u, v}$નો કોડ હોય એમ લઈને તેને ઉપરના સ્વરૂપમાં
લાવી શકાય.) ત્યાર પછી~$u$ પર અનુમાનથી બતાવી શકાય કે
$\fn{gcd}$ ખરેખર સંપૂર્ણ છે.

અલબત્ત, હમણાં ચર્ચેલાં ઉદાહરણો કરતાં ઘણાં વધુ જટિલ
સ્વ-સંદર્ભી વ્યાખ્યાઓ આપી શકાય છે. મોટાભાગની પ્રોગ્રામિંગ
ભાષાઓ કોઈ ને કોઈ રીતે વિધેયને તેના પોતાના આધારે વ્યાખ્યાયિત
કરવાની છૂટ આપે છે. આ, વિધેયની ગણતરી કરતી કલનવિધિના
\emph{સૂચક}ના આધારે વિધેય વ્યાખ્યાયિત કરવાની ક્ષમતા કરતાં
થોડું ઓછું અદ્ભુત છે; નિશ્ચિત-બિંદુનું પ્રમેય તેની પૂર્ણ
સામાન્યતામાં પછીની ક્ષમતા આપે છે.
% END OLP-0251
\endgroup


\OLEndChapterHook
% END OLP-0228
\endgroup


\OLEndPartHook
% END OLP-0208
